\chapter{The Lean Library: Foundations, Generators and Rules}\label{ch:lean-core}

\lstdefinestyle{leancore}{style=lean,
  literate={→}{{$\to$}}1 {←}{{$\leftarrow$}}1 {≤}{{$\le$}}1 {≥}{{$\ge$}}1 {≠}{{$\ne$}}1
           {ω}{{$\omega$}}1 {⟩}{{$\rangle$}}1 {⟨}{{$\langle$}}1 {∀}{{$\forall$}}1
           {∃}{{$\exists$}}1 {λ}{{$\lambda$}}1 {∑}{{$\sum$}}1 {·}{{$\cdot$}}1
           {ℂ}{{$\mathbb{C}$}}1 {ℕ}{{$\mathbb{N}$}}1 {ℤ}{{$\mathbb{Z}$}}1 {ℚ}{{$\mathbb{Q}$}}1
           {ℝ}{{$\mathbb{R}$}}1 {↦}{{$\mapsto$}}1 {∘}{{$\circ$}}1
           {α}{{$\alpha$}}1 {β}{{$\beta$}}1 {π}{{$\pi$}}1 {∈}{{$\in$}}1 {∉}{{$\notin$}}1
           {×}{{$\times$}}1 {₀}{{$_0$}}1 {₁}{{$_1$}}1 {₂}{{$_2$}}1 {₃}{{$_3$}}1 {₄}{{$_4$}}1
           {⁻}{{$^{-}$}}1 {¹}{{$^{1}$}}1 {√}{{$\surd$}}1 {∧}{{$\wedge$}}1 {∨}{{$\vee$}}1
           {¬}{{$\neg$}}1 {↔}{{$\leftrightarrow$}}1 {ε}{{$\varepsilon$}}1 {μ}{{$\mu$}}1
           {σ}{{$\sigma$}}1 {τ}{{$\tau$}}1 {φ}{{$\varphi$}}1 {ρ}{{$\rho$}}1 {ν}{{$\nu$}}1
           {δ}{{$\delta$}}1 {θ}{{$\theta$}}1 {ζ}{{$\zeta$}}1 {ι}{{$\iota$}}1 {∏}{{$\prod$}}1
           {•}{{$\bullet$}}1 {▸}{{$\triangleright$}}1,
}
\providecommand{\leanLLone}{\lstinline[style=leancore,frame=none,basicstyle=\ttfamily,%
  breaklines=true,breakatwhitespace=false]}

This chapter specifies the Lean~4 library that ships inside the ArchytasZX package, in the
directory \file{archytaszx/lean}. The library defines the semantics of ArchytasZX diagrams as
labeled tensor networks over $\C$, states every rewrite rule of the engine as a lemma about such
networks, and supplies the decision tools (exact cyclotomic evaluation, ball arithmetic,
contraction certificates) that the engine's Lean exporters invoke. The engine is
Lean-certifying: an exported proof imports these modules and is checked by the Lean
kernel~\cite{demoura2021lean4} against Mathlib~\cite{mathlib2020}. This chapter covers the
library's build configuration and module structure, and in full detail the modules that define
networks, generators, rule lemmas, mixed networks, exact arithmetic, refutation, contraction, and
the diagram data type with its \bbox{} expansion. The modules for \bbox{} families are specified in
\cref{ch:lean-boxes}, the path-sum modules in \cref{ch:lean-pathsum}, and the symbolic and linked
modules in \cref{ch:lean-sym}; the Python side that writes the Lean files is specified in
\cref{ch:lean-export-net} and \cref{ch:lean-export}.

Throughout, $\N=\{0,1,2,\dots\}$, a \emph{label} is a natural number naming one contracted index,
and all leg values are natural numbers. For a proposition $P$, $[P]\in\{0,1\}$ is its indicator.
Lean identifiers appear in typewriter type; every declaration lies in the namespace
\lean|ArchytasZX|, and every module opens \lean|Complex| and \lean|Finset| unless noted.

%======================================================================
\section{The library as a whole}\label{sec:lean-core-library}

\subsection{Toolchain, package and build}\label{sec:lean-core-build}

The library is a Lake package. Its four configuration files fix every version.

\begin{description}[style=nextline]
\item[\file{lean-toolchain}] The single line \code|leanprover/lean4:v4.35.0|: the package builds
  with Lean~4 version 4.35.0\index{Lean toolchain}. \tool{Elan}{https://github.com/leanprover/elan}
  reads this file and installs that toolchain.
\item[\file{lakefile.toml}] The package \code|archytaszx| with \code|defaultTargets = ["ArchytasZX"]|,
  one requirement and one library:
\begin{lstlisting}[style=grammar]
name = "archytaszx"
defaultTargets = ["ArchytasZX"]

[[require]]
name = "mathlib"
git = "https://github.com/leanprover-community/mathlib4"
rev = "v4.35.0"

[[lean_lib]]
name = "ArchytasZX"
\end{lstlisting}
\item[\file{lake-manifest.json}] Manifest format \code|1.3.0|, \code|packagesDir| =
  \code|.lake/packages|, \code|lakeDir| = \code|.lake|, \code|fixedToolchain| = \code|false|. It pins
  the packages of \cref{tab:lean-core-manifest}; only Mathlib is a direct requirement
  (\code|inherited| = \code|false|), the others come in through Mathlib.
\item[\file{ArchytasZX.lean}] The root module. It consists of 87 \lean|import| lines and nothing
  else; the imported modules and their transitive imports make up all 91 modules in
  \file{archytaszx/lean/ArchytasZX/}.
\end{description}

\begin{table}[htbp]
\centering\small
\caption{Packages pinned by \texttt{lake-manifest.json}.}\label{tab:lean-core-manifest}
\begin{tabular}{@{}lllp{4.6cm}@{}}
\toprule
Package & Input revision & Pinned commit & Repository \\
\midrule
\code|mathlib| & \code|v4.35.0| & \code|56a005855ea1| & \url{https://github.com/leanprover-community/mathlib4} \\
\code|plausible| & \code|main| & \code|dfa4f38190d0| & \url{https://github.com/leanprover-community/plausible} \\
\code|LeanSearchClient| & \code|main| & \code|d6b4c7be5be7| & \url{https://github.com/leanprover-community/LeanSearchClient} \\
\code|importGraph| & \code|main| & \code|a41f0e9a4dea| & \url{https://github.com/leanprover-community/import-graph} \\
\code|proofwidgets| & \code|main| & \code|df72b022e558| & \url{https://github.com/leanprover-community/ProofWidgets4} \\
\code|aesop| & \code|master| & \code|4a174c490d87| & \url{https://github.com/leanprover-community/aesop} \\
\code|Qq| & \code|master| & \code|dce0ad79c486| & \url{https://github.com/leanprover-community/quote4} \\
\code|batteries| & \code|main| & \code|66e469188fe0| & \url{https://github.com/leanprover-community/batteries} \\
\code|Cli| & \code|v4.35.0| & \code|c24fe6078c63| & \url{https://github.com/leanprover/lean4-cli} \\
\bottomrule
\end{tabular}
\end{table}

The full pinned commits are:
\begin{center}\small
\begin{tabular}{@{}ll@{}}
\toprule
Package & Commit \\
\midrule
\code|mathlib| & \code|56a005855ea17e0b0958f73ec801ae140b0a1679| \\
\code|plausible| & \code|dfa4f38190d07c37e315953e08b90fe2d95b3c24| \\
\code|LeanSearchClient| & \code|d6b4c7be5be75ec4917c873ac50c756dc1436372| \\
\code|importGraph| & \code|a41f0e9a4dea23188ea9cc1c41ac9f8bbafe0b93| \\
\code|proofwidgets| & \code|df72b022e558978a1c50264cfc4d7fec86fa8e04| \\
\code|aesop| & \code|4a174c490d8707121c4a9290800489f105d85e07| \\
\code|Qq| & \code|dce0ad79c486cb1fb504af04f953c3466b0c6db8| \\
\code|batteries| & \code|66e469188fe07963e5cf666c7d4433f45b8131c5| \\
\code|Cli| & \code|c24fe6078c630e89b489bffe59b3e20ac30d12c8| \\
\bottomrule
\end{tabular}
\end{center}
Each package's \code|configFile| is \code|lakefile.toml|, except Mathlib's, which is
\code|lakefile.lean|.

\paragraph{Building.} From \file{archytaszx/lean}, \code|lake build| fetches the pinned packages
into \file{.lake/packages}, then compiles the default target, the library \code|ArchytasZX|, that is
the root module and every module it imports. Mathlib's command \code|lake exe cache get|
downloads Mathlib's compiled files. An exported proof is a Lean source file whose header imports
the modules it needs (for example \lean|import ArchytasZX| or \lean|import ArchytasZX.PathSumNetW|).
The engine's checker \py|check(source, path, timeout=1800)| in
\module{archytaszx.semantics.lean_net} checks such a file in four steps:
\begin{enumerate}
\item find \code|lake| on the search path, else at \file{~/.elan/bin/lake}; with neither, raise
  \py|NetUnsupported("Lean (lake) is not installed")|;
\item write \py|source| to \py|path| (UTF-8);
\item if \file{archytaszx/lean/.lake/packages/mathlib} is not a directory, run
  \code|lake exe cache get| in \file{archytaszx/lean}; on a nonzero exit return
  \py|(False, stdout + stderr)|;
\item run \code|lake build| in \file{archytaszx/lean}; on a nonzero exit return
  \py|(False, stdout + stderr)|; then run \code|lake env lean <absolute path>| there and return
  \py|(exit code = 0, stdout + stderr)|.
\end{enumerate}
Each subprocess has the time limit \py|timeout| seconds. The repository's continuous integration
does not build the Lean library.

\paragraph{File conventions.} Every module starts with the comment block
\code|Copyright 2026 Arkhip A. Dmitriev / Licensed under the Apache License, Version 2.0.|, then its
imports, then a module docstring \code|/-! ... -/| headed by a Markdown title, then
\lean|namespace ArchytasZX|. Docstrings \code|/-- ... -/| state what each declaration is.

\subsection{Trust base}\label{sec:lean-core-trust}

\begin{itemize}
\item \textbf{No axioms, no gaps.} No module of \file{archytaszx/lean/ArchytasZX/} declares an
  \lean|axiom|, and none contains \lean|sorry|, \lean|admit| or \lean|native_decide|. Every theorem
  therefore rests on the Lean kernel and the axioms of Lean's core library as Mathlib uses them
  (propositional extensionality, the axiom of choice, quotient soundness); the real and complex
  numbers of Mathlib are classical.
\item \textbf{Computation in proofs.} Decidable checks (\lean|isoOK|, \lean|Diagram.check|,
  \lean|leafOK|, \lean|chainOK|, coefficient-vector reductions) are \lean|Bool| functions; an exported
  proof discharges a hypothesis \lean|h : f x = true| with \lean|decide| or \lean|rfl|, so the Lean
  kernel evaluates the check by definitional unfolding.
\item \textbf{Unsafe code.} Three modules outside this chapter (\texttt{PathSumChain},
  \texttt{PathSumQuditChain}, \texttt{SymF2Chain}) declare \lean|unsafe| metaprograms of type
  \lean|MetaM| and attach them with \lean|@[implemented_by]| to the elaboration-time functions that
  enumerate chain states; these functions run while a proof is elaborated, and the proof terms they
  help produce are checked by the kernel. No module of this chapter uses \lean|unsafe|.
\item \textbf{Noncomputable definitions.} Every definition with a complex or real value (tensor
  entries, sums, evaluations) is \lean|noncomputable|; theorems about them are proved, never
  evaluated.
\end{itemize}

\subsection{Module structure}\label{sec:lean-core-modules}

\Cref{tab:lean-core-imports} lists every module of \file{archytaszx/lean/ArchytasZX/} with its
direct imports, its line count, its number of top-level theorems and of top-level definitions
(\lean|def|, \lean|abbrev|, \lean|structure|, \lean|inductive|), and the chapter that specifies it.
Imports from Mathlib are counted, not named, except in \texttt{Net}, which carries all 43
Mathlib imports of the foundation. The library has 79{,}797 lines, 3{,}971 top-level theorems and
1{,}994 top-level definitions. \texttt{Light} has no declarations: it imports 87 modules, none of which imports \texttt{Cyclo} or \texttt{NumField} directly or indirectly; the benchmark's Lean check tries it in place of the full library (\cref{ch:bench-pipeline}). The modules \texttt{AdderFamEnv}, \texttt{FamDiagram}, \texttt{FamExpandEq}, \texttt{FamFcaExpand}, \texttt{LinkedMerge}, and \texttt{LinkedPair} also avoid \texttt{Cyclo} and \texttt{NumField} and are not imported by \texttt{Light}. \Cref{fig:lean-core-spine-a,fig:lean-core-spine} draw the import structure of the modules
of this chapter.

{\small
\begin{xltabular}{\textwidth}{@{}l>{\raggedright\arraybackslash}Xrrrc@{}}
\caption{Every module of the library: direct imports (modules of \texttt{ArchytasZX}, then the number
of Mathlib imports), lines, theorems, definitions, and the specifying chapter: C this chapter,
B \cref{ch:lean-boxes}, P \cref{ch:lean-pathsum}, S \cref{ch:lean-sym}.}\label{tab:lean-core-imports}\\
\toprule
Module & Imports & Lines & Thm. & Def. & Ch. \\
\midrule
\endfirsthead
\toprule
Module & Imports & Lines & Thm. & Def. & Ch. \\
\midrule
\endhead
\bottomrule
\endfoot
\texttt{AdderFam} & Closed, ChainInduct & 504 & 26 & 11 & B \\
\texttt{AdderFamDraper} & AdderFamFca & 1517 & 115 & 42 & B \\
\texttt{AdderFamDraperEng} & AdderFamDraper, AdderFamFcaEng & 1610 & 105 & 57 & B \\
\texttt{AdderFamEnv} & AdderFam, ChainMixed & 69 & 2 & 3 & B \\
\texttt{AdderFamFca} & AdderFamEnv, Ball, Weights, PairEval & 1672 & 131 & 41 & B \\
\texttt{AdderFamFcaEng} & AdderFamFca & 1304 & 91 & 53 & B \\
\texttt{Assume} & Bundle & 96 & 4 & 4 & C \\
\texttt{Ball} & Cyclo, 2 Mathlib & 280 & 36 & 4 & C \\
\texttt{Bialgebra} & ZXRules & 189 & 4 & 0 & C \\
\texttt{BoxedBridge} & FamK, DiagramBridge & 631 & 27 & 18 & B \\
\texttt{BoxFree} & Induct & 154 & 7 & 0 & B \\
\texttt{Bridge} & Expand & 133 & 0 & 16 & B \\
\texttt{Bundle} & WZ & 1876 & 109 & 46 & B \\
\texttt{ChainBox} & Bridge & 85 & 2 & 7 & B \\
\texttt{ChainBridge} & ChainForm, ChainLabel, PeriodicChain & 244 & 14 & 14 & B \\
\texttt{ChainEvery} & ChainLink & 313 & 26 & 9 & B \\
\texttt{ChainFam} & ChainBridge & 520 & 22 & 3 & B \\
\texttt{ChainForm} & ChainBox & 159 & 13 & 4 & B \\
\texttt{ChainInduct} & Bundle & 313 & 21 & 5 & B \\
\texttt{ChainLabel} & Elim & 157 & 3 & 2 & B \\
\texttt{ChainLink} & ChainBox, ChainInduct & 99 & 4 & 5 & B \\
\texttt{ChainMixed} & ChainInduct & 422 & 27 & 3 & B \\
\texttt{ChainNode} & ChainLink & 489 & 41 & 14 & B \\
\texttt{ChainPair} & ChainFam & 159 & 7 & 1 & B \\
\texttt{Closed} & Rules & 152 & 17 & 1 & C \\
\texttt{Completeness} & Net, Rules, WRules, Closed, ZXRules & 159 & 7 & 0 & C \\
\texttt{Connective} & WRules & 255 & 10 & 0 & C \\
\texttt{Contract} & Refute & 502 & 22 & 19 & C \\
\texttt{CopyBridge} & CopyPin, Linked & 867 & 33 & 14 & B \\
\texttt{CopyPin} & Ball, Mixed, FamK & 128 & 6 & 0 & B \\
\texttt{Cyclo} & Bundle, 1 Mathlib & 451 & 47 & 11 & C \\
\texttt{Diagram} & Net & 246 & 6 & 20 & C \\
\texttt{DiagramBridge} & Diagram & 476 & 17 & 14 & B \\
\texttt{Digits} & Elim & 196 & 13 & 12 & C \\
\texttt{Elim} & Bridge, WZ & 331 & 9 & 8 & C \\
\texttt{Expand} & Diagram & 159 & 0 & 20 & C \\
\texttt{ExpandFacts} & Digits & 230 & 13 & 2 & C \\
\texttt{FamBridge} & NetDigits & 268 & 11 & 3 & B \\
\texttt{FamCheck} & FamBridge & 153 & 8 & 10 & B \\
\texttt{FamDiagram} & OrderK, PairBox & 172 & 8 & 19 & B \\
\texttt{FamExpandEq} & FamDiagram & 68 & 7 & 0 & B \\
\texttt{FamFcaExpand} & FamExpandEq & 135 & 14 & 4 & B \\
\texttt{FamK} & OrderK & 243 & 12 & 2 & B \\
\texttt{FuseB} & Unfuse & 284 & 14 & 1 & C \\
\texttt{Gather} & PortBox & 246 & 22 & 1 & C \\
\texttt{GaussSign} & Cyclo, PathSumQudit & 174 & 8 & 4 & C \\
\texttt{Hopf} & Bialgebra & 245 & 10 & 0 & C \\
\texttt{Induct} & NodeBox & 299 & 15 & 5 & B \\
\texttt{Iso} & Net & 156 & 7 & 5 & C \\
\texttt{Light} & Net, Iso, Rules, Closed, ZXRules, Bialgebra, Hopf, WRules, WAssoc, WFuse, Connective, WZ, Bundle, Unfuse, FuseB, SymNet, SymChunk, SymHint, SymF2, SymF2Chain, SymF2Bridge, DiagramBridge, Port, Assume, Completeness, Weights, Diagram, Expand, Bridge, ChainBox, Elim, Digits, ExpandFacts, Relabel, PortBox, Gather, Split, NodeBox, Induct, BoxFree, Positional, Tower, NodeVal, NodeBridge, Pin, NetDigits, FamBridge, FamCheck, Order, OrderK, FamK, BoxedBridge, Linked, LinkedPeel, LinkedSlide, LinkedRegroup, LinkedSplit, LinkedForm, LinkedGhost, LinkedSplice, LinkedFirst, LinkedDropX, LinkedRead, Mixed, ChainInduct, ChainLink, ChainEvery, Periodic, PeriodicChain, ChainMixed, ChainLabel, ChainForm, ChainBridge, ChainFam, ChainPair, TreeInduct, PhaseMap, ChainNode, PairBox, PairNet, PairEval, PairFam, AdderFam, PairSum, Sym, SymCanon, SymNorm & 96 & 0 & 0 & B \\
\texttt{Linked} & Relabel & 515 & 10 & 23 & S \\
\texttt{LinkedDropX} & LinkedGhost & 302 & 14 & 4 & S \\
\texttt{LinkedFirst} & LinkedSplice & 656 & 42 & 7 & S \\
\texttt{LinkedForm} & LinkedSlide, LinkedRegroup & 170 & 12 & 8 & S \\
\texttt{LinkedGhost} & LinkedForm & 958 & 58 & 15 & S \\
\texttt{LinkedMerge} & LinkedPair & 294 & 26 & 6 & S \\
\texttt{LinkedPair} & LinkedFirst & 455 & 17 & 11 & S \\
\texttt{LinkedPeel} & Linked, PortBox & 1717 & 86 & 18 & S \\
\texttt{LinkedRead} & LinkedDropX, LinkedSplice & 1165 & 55 & 19 & S \\
\texttt{LinkedRegroup} & Linked & 134 & 6 & 6 & S \\
\texttt{LinkedSlide} & BoxFree, LinkedPeel, Weights & 576 & 45 & 6 & S \\
\texttt{LinkedSplice} & LinkedGhost, LinkedPeel & 1267 & 58 & 24 & S \\
\texttt{LinkedSplit} & LinkedPeel, Iso & 1812 & 86 & 15 & S \\
\texttt{Mixed} & Closed & 210 & 16 & 3 & C \\
\texttt{Net} & 43 Mathlib & 480 & 20 & 40 & C \\
\texttt{NetDigits} & Pin & 244 & 6 & 2 & C \\
\texttt{NodeBox} & Split & 1440 & 85 & 7 & B \\
\texttt{NodeBridge} & NodeVal & 339 & 15 & 1 & B \\
\texttt{NodeVal} & Tower & 406 & 34 & 11 & B \\
\texttt{NumField} & Weights, 4 Mathlib & 110 & 7 & 3 & C \\
\texttt{Order} & FamCheck & 280 & 20 & 7 & C \\
\texttt{OrderK} & Order & 246 & 13 & 8 & C \\
\texttt{PairBox} & ChainLabel, Relabel & 229 & 0 & 31 & B \\
\texttt{PairEval} & PairNet & 646 & 34 & 4 & B \\
\texttt{PairFam} & PairEval & 93 & 3 & 0 & B \\
\texttt{PairNet} & PairBox & 815 & 59 & 3 & B \\
\texttt{PairSum} & PairFam & 116 & 9 & 1 & B \\
\texttt{Paper} & Net, Rules, Closed, ZXRules, Cyclo & 172 & 12 & 9 & C \\
\texttt{PathSum} & Sym, Cyclo & 183 & 17 & 2 & P \\
\texttt{PathSumAmp} & PathSumRules & 1964 & 81 & 85 & P \\
\texttt{PathSumChain} & PathSumSupport, PathSumParity & 308 & 6 & 19 & P \\
\texttt{PathSumDigits} & PathSumNet, Positional & 1059 & 47 & 21 & P \\
\texttt{PathSumMixed} & PathSumQudit, PathSumAmp & 492 & 25 & 9 & P \\
\texttt{PathSumMixedNet} & PathSumMixed, PathSumNetAmp, PathSumQuditNet, PathSumQuditDigits & 1087 & 56 & 21 & P \\
\texttt{PathSumNet} & PathSumRules & 676 & 42 & 21 & P \\
\texttt{PathSumNetAmp} & PathSumNetW, PathSumAmp & 409 & 19 & 7 & P \\
\texttt{PathSumNetW} & PathSumW, PathSumChain & 91 & 7 & 3 & P \\
\texttt{PathSumParity} & PathSumRules & 109 & 2 & 4 & P \\
\texttt{PathSumQudit} & PathSum, 1 Mathlib & 1138 & 68 & 55 & P \\
\texttt{PathSumQuditChain} & PathSumChain, PathSumQudit & 103 & 3 & 5 & P \\
\texttt{PathSumQuditDigits} & PathSumQuditNet & 636 & 28 & 12 & P \\
\texttt{PathSumQuditNet} & PathSumQudit, PathSumDigits, GaussSign, PathSumQuditChain & 462 & 29 & 14 & P \\
\texttt{PathSumRules} & PathSumSem & 2198 & 91 & 63 & P \\
\texttt{PathSumSem} & PathSum, Ball & 2173 & 139 & 101 & P \\
\texttt{PathSumSupport} & PathSumAmp & 58 & 2 & 1 & P \\
\texttt{PathSumW} & PathSumDigits & 224 & 9 & 3 & P \\
\texttt{Periodic} & ChainInduct & 42 & 1 & 0 & B \\
\texttt{PeriodicChain} & Periodic & 301 & 22 & 6 & B \\
\texttt{PhaseMap} & Bundle & 114 & 5 & 4 & B \\
\texttt{Pin} & NodeBridge & 152 & 8 & 4 & C \\
\texttt{Port} & Bundle & 1117 & 61 & 11 & B \\
\texttt{PortBox} & Relabel, ExpandFacts & 934 & 63 & 7 & B \\
\texttt{Positional} & BoxFree, Bundle & 141 & 12 & 3 & C \\
\texttt{Refute} & Cyclo & 547 & 26 & 33 & C \\
\texttt{RefuteMod} & Contract & 1789 & 99 & 61 & C \\
\texttt{RefuteSym} & RefuteMod, 1 Mathlib & 1897 & 92 & 63 & C \\
\texttt{Relabel} & Digits & 135 & 4 & 4 & C \\
\texttt{Rules} & Iso, 1 Mathlib & 324 & 19 & 0 & C \\
\texttt{Split} & Gather & 195 & 6 & 2 & C \\
\texttt{Sym} & Closed & 1032 & 78 & 72 & S \\
\texttt{SymCanon} & SymNorm & 851 & 35 & 45 & S \\
\texttt{SymCarry} & SymRange & 728 & 29 & 17 & S \\
\texttt{SymChunk} & SymNet & 283 & 19 & 16 & S \\
\texttt{SymDig} & SymCarry, Cyclo & 845 & 34 & 24 & S \\
\texttt{SymF2} & SymNet & 1231 & 87 & 58 & S \\
\texttt{SymF2Bridge} & SymF2, DiagramBridge & 160 & 3 & 2 & S \\
\texttt{SymF2Chain} & SymF2 & 136 & 0 & 7 & S \\
\texttt{SymFold} & SymCarry & 1112 & 34 & 25 & S \\
\texttt{SymFourier} & SymHint, PathSumSem & 851 & 39 & 33 & S \\
\texttt{SymHint} & SymChunk & 795 & 30 & 28 & S \\
\texttt{SymNet} & SymCanon & 1581 & 72 & 54 & S \\
\texttt{SymNorm} & Sym, Bundle & 926 & 39 & 24 & S \\
\texttt{SymRange} & SymFourier & 1210 & 46 & 41 & S \\
\texttt{Tower} & Positional & 161 & 11 & 1 & C \\
\texttt{TreeInduct} & ChainInduct & 330 & 13 & 7 & B \\
\texttt{Unfuse} & Bundle & 221 & 12 & 0 & C \\
\texttt{WAssoc} & WRules & 120 & 6 & 2 & C \\
\texttt{Weights} & Net, Rules, Closed, ZXRules, Completeness, WZ, Hopf, Connective & 9479 & 213 & 64 & C \\
\texttt{WFuse} & WAssoc, FuseB, NodeBox, Weights & 225 & 10 & 0 & C \\
\texttt{WRules} & Hopf & 289 & 17 & 1 & C \\
\texttt{WZ} & Connective & 242 & 6 & 0 & C \\
\texttt{ZXRules} & Closed & 398 & 28 & 0 & C \\
\end{xltabular}
}

\begin{figure}[htbp]
\centering
\begin{tikzpicture}[x=1.55cm,y=0.95cm,every node/.style={font=\scriptsize,draw,rounded corners=1pt,inner sep=2pt},>={Stealth[length=1.6mm]}]
\node (net) at (3,0) {Net};
\node (iso) at (1,1) {Iso};
\node (diag) at (5,1) {Diagram};
\node (rules) at (1,2) {Rules};
\node (expand) at (5,2) {Expand};
\node (closed) at (1,3) {Closed};
\node (zx) at (0,4) {ZXRules};
\node (mixed) at (2,4) {Mixed};
\node (bia) at (0,5) {Bialgebra};
\node (hopf) at (0,6) {Hopf};
\node (wr) at (0,7) {WRules};
\node (conn) at (-1,8) {Connective};
\node (wassoc) at (1,8) {WAssoc};
\node (compl) at (2,7) {Completeness};
\node (wz) at (-1,9) {WZ};
\node (weights) at (1,10) {Weights};
\node (numf) at (0,11) {NumField};
\node (bundle) at (-1,10.2) {[Bundle]};
\node (unf) at (-2,11) {Unfuse};
\node (fuseb) at (-2,12) {FuseB};
\node (cyclo) at (3,11) {Cyclo};
\node (refute) at (3,12) {Refute};
\node (contract) at (3,13) {Contract};
\node (ball) at (4.3,12) {Ball};
\node (paper) at (2,12) {Paper};
\node (gauss) at (5.2,12.6) {GaussSign};
\node (assume) at (-3,11) {Assume};
\draw[->] (net) -- (iso); \draw[->] (net) -- (diag); \draw[->] (iso) -- (rules);
\draw[->] (diag) -- (expand); \draw[->] (rules) -- (closed); \draw[->] (closed) -- (zx);
\draw[->] (closed) -- (mixed); \draw[->] (zx) -- (bia); \draw[->] (bia) -- (hopf);
\draw[->] (hopf) -- (wr); \draw[->] (wr) -- (conn); \draw[->] (wr) -- (wassoc);
\draw[->] (wr) -- (compl); \draw[->] (conn) -- (wz); \draw[->] (compl) -- (weights);
\draw[->] (wz) -- (weights); \draw[->] (weights) -- (numf); \draw[->] (wz) -- (bundle);
\draw[->] (bundle) -- (unf); \draw[->] (unf) -- (fuseb); \draw[->] (bundle) -- (cyclo);
\draw[->] (cyclo) -- (refute); \draw[->] (refute) -- (contract); \draw[->] (cyclo) -- (ball);
\draw[->] (cyclo) -- (paper); \draw[->] (cyclo) -- (gauss); \draw[->] (bundle) -- (assume);
\end{tikzpicture}
\caption{Import structure of the network and rule-lemma modules of this chapter; an arrow points
from an imported module to an importing one; bracketed modules belong to \cref{ch:lean-boxes}.
Further imports are listed in \cref{tab:lean-core-imports}.}\label{fig:lean-core-spine-a}
\end{figure}

\begin{figure}[htbp]
\centering
\begin{tikzpicture}[x=1.7cm,y=0.9cm,every node/.style={font=\scriptsize,draw,rounded corners=1pt,inner sep=2pt},>={Stealth[length=1.6mm]}]
\node (expand) at (0,0) {Expand};
\node (bridge) at (0,1) {[Bridge]};
\node (elim) at (0,2) {Elim};
\node (digits) at (0,3) {Digits};
\node (ef) at (-1,4) {ExpandFacts};
\node (rel) at (1,4) {Relabel};
\node (portbox) at (0,5) {[PortBox]};
\node (gather) at (0,6) {Gather};
\node (split) at (0,7) {Split};
\node (nodebox) at (0,8) {[NodeBox]};
\node (induct) at (0,9) {[Induct, BoxFree]};
\node (pos) at (2,9.5) {Positional};
\node (tower) at (2,10.5) {Tower};
\node (nv) at (2,11.5) {[NodeVal, NodeBridge]};
\node (pin) at (2,12.5) {Pin};
\node (nd) at (2,13.5) {NetDigits};
\node (fam) at (0,13.5) {[FamBridge, FamCheck]};
\node (order) at (0,14.5) {Order};
\node (orderk) at (0,15.5) {OrderK};
\node (wfuse) at (-2,9) {WFuse};
\draw[->] (expand) -- (bridge); \draw[->] (bridge) -- (elim); \draw[->] (elim) -- (digits);
\draw[->] (digits) -- (ef); \draw[->] (digits) -- (rel); \draw[->] (ef) -- (portbox);
\draw[->] (rel) -- (portbox); \draw[->] (portbox) -- (gather); \draw[->] (gather) -- (split);
\draw[->] (split) -- (nodebox); \draw[->] (nodebox) -- (induct); \draw[->] (induct) -- (pos);
\draw[->] (pos) -- (tower); \draw[->] (tower) -- (nv); \draw[->] (nv) -- (pin);
\draw[->] (pin) -- (nd); \draw[->] (nd) -- (fam); \draw[->] (fam) -- (order);
\draw[->] (order) -- (orderk); \draw[->] (nodebox) -- (wfuse);
\end{tikzpicture}
\caption{Import structure of the diagram and \bbox{} expansion modules of this chapter (\texttt{Elim} also imports \texttt{WZ}; \texttt{WFuse} also imports \texttt{WAssoc}, \texttt{FuseB} and \texttt{Weights}; \texttt{Positional} also imports \texttt{Bundle}); bracketed modules belong to \cref{ch:lean-boxes}.}\label{fig:lean-core-spine}
\end{figure}

\subsection{Relation to the Python engine}\label{sec:lean-core-python}

The library mirrors the engine's data in three layers.
\begin{enumerate}
\item \textbf{Generators.} The constructors \lean|Z|, \lean|X|, \lean|F|, \lean|T|, \lean|Ti|,
  \lean|W|, \lean|B|, \lean|S| of \lean|Kind| (\cref{sec:lean-core-net}) are the engine's
  generators Z spider, X spider, Fourier box $F$, triangle, inverse triangle, W node, dimension
  binder $B$ and dimension splitter $S$; the exporter's table \py|_KINDS| maps each engine kind to
  the constructor of the same name. \lean|I| is the bare wire. \lean|ZW| and \lean|XW| are the
  Z and X spiders with arbitrary complex coefficients, the engine's weighted Z box and its X
  counterpart. \lean|ZB|, \lean|XB|, \lean|WB| and \lean|Bun| carry \bbox{} families: spiders whose
  last ports bundle copies, and copies of a generator along a multiplicity. The discard $G$ has no
  constructor; a diagram with discards is held as its purification (\cref{sec:lean-core-mixed}).
\item \textbf{Diagrams.} \lean|Diagram| (\cref{sec:lean-core-diagram}) is the engine's diagram
  as data, with ports, wires, boundary and \bboxes; \lean|DEqual_iff| equates diagram equality with
  network equality, so the exporter states its final theorem about diagrams.
\item \textbf{Rules.} Every rewrite rule the exporter supports has one or more lemmas whose
  conclusion is a \lean|LocalEq| between the rule's two sides. \Cref{tab:lean-core-rulemap} lists
  the engine rule names (keys of the exporter's dispatch table \py|_RULES| in
  \module{archytaszx.semantics.lean_net}) with the lemmas their steps invoke. A rule step becomes
  \lean|iso| into the shape $C\mathbin{+\!\!+}L$, then \lean|rewrite| with the rule's lemma, then
  \lean|wire_elim| for each bare wire the step leaves.
\end{enumerate}

{\small
\begin{xltabular}{\textwidth}{@{}>{\raggedright\arraybackslash}p{6.9cm}>{\raggedright\arraybackslash}X@{}}
\caption{Engine rules and the lemmas their exported steps invoke (module in parentheses). A
converse rule not listed invokes its forward rule's lemma under
\lean|LocalEq.symm'|.}\label{tab:lean-core-rulemap}\\
\toprule
Engine rule & Lemmas \\
\midrule
\endfirsthead
\toprule
Engine rule & Lemmas \\
\midrule
\endhead
\bottomrule
\endfoot
\code|spider_fusion|, \code|z_fusion_prime_d| & \lean|fusionZ|, \lean|fusionX| (Rules); \lean|zcFusion|, \lean|xcFusion| (Weights); \lean|fusionZG|, \lean|fusionXG| (FuseB) for bundled ports \\
\code|identity_removal| & \lean|identityZ|, \lean|identityX| (Rules) \\
\code|zx_cap|, \code|zx_cap_swapped| & \lean|cap|, \lean|capSwapped| (ZXRules) \\
\code|fourier_state_color_change|, \code|..._swapped| & \lean|colorState|, \lean|colorEffect|, \lean|colorStateX|, \lean|colorEffectX| (ZXRules) \\
\code|state_copy|, \code|state_copy_swapped| & \lean|stateCopy|, \lean|stateCopySwapped| (ZXRules) \\
\code|triangle_inverse_cancellation| & \lean|triangle|, \lean|triangle'| (ZXRules) \\
\code|fourier_cancellation| & \lean|fourier4| (ZXRules) \\
\code|bialgebra|, \code|bialgebra_swapped|, \code|zxw_b2_converse| & \lean|bialgebra|, \lean|bialgebraSwapped| (Bialgebra) \\
\code|hopf|, \code|hopf_swapped| & \lean|hopf|, \lean|hopfSwapped| (Hopf); \lean|hopfW|, \lean|hopfSwappedW|, \lean|hopfXW|, \lean|hopfSwappedXW|, \lean|hopfWXW|, \lean|hopfSwappedWXW| (Weights) \\
\code|w_fusion|, \code|qf_wf| & \lean|wFusion| (WRules) \\
\code|w_identity| & \lean|wIdentity| (WRules) \\
\code|w_zero_copy|, \code|triangle_zero_state|, \code|zxw_bs0| & \lean|wZeroCopy|, \lean|triZeroState|, \lean|triInvZeroState| (WRules) \\
\code|triangle_zero_effect| & \lean|triZeroEffect| (WRules) \\
\code|w_z_effect|, \code|w_zero_effect| & \lean|wZEffect|, \lean|wZeroEffect| (WRules) \\
\code|connective_states|, \code|connective_states_swapped| & \lean|connStates|, \lean|connEffects|, \lean|connStatesX|, \lean|connEffectsX| (Connective) \\
\code|connective_inverse|, \code|connective_inverse_bind_first| & \lean|connInverse|, \lean|connInverseBind| (Connective) \\
\code|wz_bialgebra|, \code|zxw_bzw_converse| & \lean|wzBialgebra| (WZ) \\
\code|port_box_unfusion| & \lean|unfuseZB|, \lean|unfuseXB|, \lean|unfuseXBi| (Unfuse) \\
\code|zxw_s2|, \code|zxw_s2_pink| & \lean|identityZ|, \lean|identityX| (Rules), \lean|identityXZ| (Weights) \\
\code|zxw_s1| & \lean|zcFusionN| (Weights) \\
\code|zxw_sym|, \code|qf_w_sym| (and converses) & \lean|wPerm| (Completeness) \\
\code|zxw_aso| (and converse) & \lean|wAso| (Completeness) \\
\code|qf_z_flex| (and converse) & \lean|zPerm| (Completeness) \\
\code|zxw_k0| & \lean|kCopy| (Completeness) \\
\code|zxw_ept|, \code|zxw_zer|, \code|zxw_ad|, \code|zxw_pcy|, \code|zxw_k2| & \lean|zcEpt|, \lean|zcZer|, \lean|zcAd|, \lean|zcPcy|, \lean|zcK2| (Weights) \\
\code|triangle_definition_converse|, \code|triangle_inverse_definition_converse| & \lean|zcT|, \lean|zcTi| (Weights) \\
\code|phase_to_weight|, \code|weight_to_phase| & \lean|zcCongr| (Weights) \\
\code|scalar_box| & \lean|zcScalar| (Weights) \\
\code|zxw_p1|, \code|zxw_p1_multiplier| & \lean|fourierSq|, \lean|fourierSq'|, \lean|fourierMult| (Weights) \\
\code|zxw_ta|, \code|zxw_bsj|, \code|zxw_k1|, \code|zxw_d1| & \lean|zxwTA|, \lean|zxwBsj|, \lean|kOne|, \lean|zxwD1| (Weights) \\
\code|x_spider_definition|, \code|zxw_h1| & \lean|xDef|, \lean|xDefT|, \lean|kjNeg|, \lean|xDefW| (Weights) \\
\code|binder_definition| & \lean|bindDef|, \lean|bindDef0|, \lean|bindDefP| (Weights) \\
\code|zxw_vw|, \code|zxw_zv|, \code|zxw_va|, \code|zxw_kz|, \code|zxw_ww|, \code|zxw_hd| (and converses) & \lean|zxwVW|, \lean|zxwZV|, \lean|zxwVA|, \lean|zxwKZ|, \lean|zxwWW|, \lean|zxwHD| (Weights) \\
\code|qf_mw|, \code|qf_e|, \code|qf_k0|, \code|qf_d1|, \code|qf_va|, \code|qf_hd|, \code|qf_ww|, \code|qf_dd|, \code|qf_bzw|, \code|qf_dzx|, \code|qf_ta|, \code|qf_s1| (and converses) & \lean|qfMW|, \lean|qfE|, \lean|qfK0|, \lean|qfD1|, \lean|qfVA|, \lean|qfHD|, \lean|qfWW|, \lean|qfDD|, \lean|qfBZW|, \lean|qfDZX|, \lean|qfTA|, \lean|qfS1|, \lean|qfS1c| (Weights) \\
\end{xltabular}
}

Besides the rule lemmas, the exporters use: \lean|iso|, \lean|iso_scalar| (\texttt{Iso}) to reorder
nodes and rename summed labels; \lean|rewrite| and \lean|Equal.trans|, \lean|Equal.mul_trans|,
\lean|Equal.rescale| (\texttt{Net}) to chain steps; \lean|wire_elim| (\texttt{Rules}) to remove bare
wires; \lean|bunI| (\texttt{FuseB}) and \lean|wFuseZ| (\texttt{WFuse}) for copied bare wires and W
fusion over bundled ports; \lean|netCoef| and \lean|root_unique| (\texttt{NumField}) for spider
coefficients given by roots of integer polynomials; \lean|Equal.of_mapDim'| (\texttt{Assume}) for
dimensions assumed equal; \lean|DEqual_iff| (\texttt{Diagram}) for the restatement over diagrams;
\lean|refute|, \lean|eval_cyc|, \lean|inRange_of| (\texttt{Refute}) and \lean|eval_ve_at|,
\lean|eval_ve_groups| (\texttt{Contract}) for refutations; \lean|gauss_rt| (\texttt{GaussSign}) for
quadratic Gauss sums.

%======================================================================
\section{Tensor networks (\texttt{Net})}\label{sec:lean-core-net}

\texttt{ArchytasZX.Net} is the foundation: it defines generator kinds, network nodes, networks,
environments, the tensor of every generator, the evaluation of a network as a sum over its internal
labels, the equality of scaled networks, and the frame theorem that makes a local equation hold in
every context. The semantics is the standard one of the qudit ZX calculus and of the ZXW calculus of
\citet{poor2023zxw} and \citet{wang2023qufinitezxw}: each generator is a tensor with one index per
port, and a network contracts shared indices~\cite{coecke2011interacting}.

\subsection{Data}\label{sec:lean-core-net-data}

\begin{lstlisting}[style=leancore]
inductive Kind
  | Z | X | F | T | Ti | W | B | S | I | ZW
  | ZB (cs : List (List ℕ)) | XB (co ci : List (List ℕ)) | WB (cs : List (List ℕ))
  | Bun (m : ℕ) (base : Kind)
  | XW
  | ZP (dx dy : ℕ) (c : ℤ)
  deriving DecidableEq, Repr

structure Node where
  kind : Kind
  outs : List ℕ      -- output port labels
  ins  : List ℕ      -- input port labels
  dim  : ℕ           -- dimension index
  ph   : ℕ           -- phase index (second dimension index of B and S)
  deriving DecidableEq, Repr

structure Net where
  nodes : List Node
  int   : List ℕ     -- summed (internal) labels
  bnd   : List ℕ     -- boundary labels
  deriving DecidableEq, Repr

structure Env where
  dimv  : ℕ → ℕ                       -- dimension of each dimension index
  phase : ℕ → ℕ → ℝ                   -- phase index, k: phase in turns
  labs  : List (ℕ × ℕ)                -- (label, dimension index)
  mult  : ℕ → ℕ := fun _ => 0         -- value of each multiplicity symbol
  coef  : ℕ → ℕ → ℂ := fun _ _ => 1   -- coefficient index, k: coefficient
\end{lstlisting}

\begin{definition}[Kinds]\label{def:lean-core-kind}
A \term{kind} is one of: \lean|Z| (Z spider), \lean|X| (X spider), \lean|F| (Fourier box),
\lean|T| (triangle), \lean|Ti| (inverse triangle), \lean|W| (W node), \lean|B| (dimension binder),
\lean|S| (dimension splitter), \lean|I| (bare wire), \lean|ZW| (Z spider with complex
coefficients), \lean|XW| (X spider with complex coefficients), \lean|ZB cs| (Z spider whose last
\lean|cs.length| inputs bundle along the chains \lean|cs|), \lean|XB co ci| (X spider whose last
outputs bundle along \lean|co| and last inputs along \lean|ci|), \lean|WB cs| (W node whose last
outputs bundle along \lean|cs|), \lean|Bun m base| (copies of \lean|base| along the
multiplicity symbol \lean|m|), and \lean|ZP dx dy c| (the pair phase $e^{2\pi i\,c\,x\,y/M}$ on two
legs of dimension indices \lean|dx| and \lean|dy|, $M$ the node's dimension). A \term{chain} is a list of multiplicity symbols, outermost first.
Equality of kinds and nodes is structural (\lean|DecidableEq|).
\end{definition}

A node's ports are its outputs followed by its inputs (\lean|Node.labels n = n.outs ++ n.ins|); a
label shared by several ports is one contracted index. \lean|labelsOf ns| is
\lean|ns.flatMap Node.labels|. \lean|Node.rename f n| applies \lean|f| to every port label.

\begin{definition}[Label dimension]
For an environment $E$, \lean|labDim E l = E.dimv ((E.labs.lookup l).getD 0)|: the dimension of
the index named by the first entry for $l$ in \lean|E.labs|, the index $0$ when $l$ has no entry.
\lean|labIdx labs l = (labs.lookup l).getD 0| (in \texttt{Iso}) is the index itself.
\end{definition}

\subsection{Scalars and generator tensors}\label{sec:lean-core-net-tensors}

With $\theta\in\R$ and $d\in\N$,
\[
  \mathrm{ex}(\theta)=e^{2\pi i\theta},\qquad
  \omega_d=\mathrm{om}(d)=e^{2\pi i/d},\qquad
  F_d[j,k]=\mathrm{fou}(d,j,k)=\frac{\omega_d^{jk}}{\sqrt d}.
\]
For $d=0$, Lean's division by zero makes $\omega_0=e^{0}=1$ and $F_0[j,k]=0$. The generator
entries, for leg values $v_i\in\N$ (not restricted to $[0,d)$; ranges are imposed by evaluation), a
phase function $\varphi:\N\to\R$ in turns and a coefficient function $c:\N\to\C$:
\begin{align*}
\mathrm{tZ}(d,\varphi,v_1\dots v_n)&=\sum_{k<d}[v_1=\dots=v_n=k]\,\mathrm{ex}(\varphi(k)),\\
\mathrm{tZw}(d,c,v_1\dots v_n)&=\sum_{k<d}[v_1=\dots=v_n=k]\,c(k),\\
\mathrm{tX}(d,\varphi,o_1\dots o_m;i_1\dots i_n)&=\sum_{k<d}\mathrm{ex}(\varphi(k))
  \prod_{a}F_d[o_a,k]\prod_{b}\overline{F_d[i_b,k]},\\
\mathrm{tXw}(d,c,o;i)&=\sum_{k<d}c(k)\prod_{a}F_d[o_a,k]\prod_{b}\overline{F_d[i_b,k]},\\
\mathrm{tT}(o,i)&=[o=0\lor o=i],\qquad
\mathrm{tTi}(o,i)=\begin{cases}1&o=i\\-1&o=0\land i\ge 1\land o\ne i\\0&\text{otherwise,}\end{cases}\\
\mathrm{tW}(d,o_1\dots o_m,i)&=\begin{cases}[o_1=\dots=o_m=0]&i=0\\
  [i<d]\,[\mathrm{nz}(o)=(i)]&i\ne 0,\end{cases}\\
\mathrm{tB}(s,t,k,a,b)&=[a<s]\,[b<t]\,[k=at+b],\qquad
\mathrm{tI}(d,a,b)=[a=b]\,[a<d],
\end{align*}
where $\mathrm{nz}(o)$ is the list of the nonzero entries of $o$ in order (\lean|o.filter (· ≠ 0)|).
So the X spider is the Z spider with $F$ on every output and $\overline{F}$ on every input, the W
node sends $\ket{0}$ to $\ket{0\cdots0}$ and $\ket{i}$, $0<i<d$, to the sum of the $m$ one-hot
vectors $\ket{0\cdots i\cdots 0}$, and the binder $B$ maps $\ket{a}\otimes\ket{b}$ to $\ket{at+b}$.

\subsection{Digits and bundled legs}\label{sec:lean-core-net-digits}

A bundled leg carries in one value the values of many copies of a port; the copies are the digits of
the value.
\begin{align*}
\mathrm{digit}(d,i,v)&=\lfloor v/d^i\rfloor\bmod d,\\
\mathrm{digs}(d,k,v)&=(\mathrm{digit}(d,0,v),\dots,\mathrm{digit}(d,k-1,v)),\\
\mathrm{digsC}(\mu,d,[\,],v)&=(v),\\
\mathrm{digsC}(\mu,d,m{::}c,v)&=\mathrm{flatMap}\bigl(\mathrm{digsC}(\mu,d,c,\cdot)\bigr)
  \bigl(\mathrm{digs}(d^{\prod_{x\in c}\mu(x)},\mu(m),v)\bigr).
\end{align*}
Here $\mu$ is \lean|E.mult|. So a value bundling along the chain $m{::}c$ splits into $\mu(m)$ digits
in base $d^{\prod_c\mu}$, least significant first, and each digit splits further along $c$.
\begin{itemize}
\item \lean|tailVals μ d vs cs| is the concatenation over the pairs $(v,c)$ of \lean|zipWith| of
  $\mathrm{digsC}(\mu,d,c,v)$.
\item \lean|bundled μ d vs cs| keeps the first $|vs|-|cs|$ values and replaces the last $|cs|$ by
  \lean|tailVals μ d (last values) cs|.
\item \lean|Kind.outCs k n| and \lean|Kind.inCs k n| are the chains of the $n$ outputs, inputs of a
  copied generator: for \lean|XB co _| the outputs get $n-|co|$ empty chains then \lean|co|, for
  \lean|WB cs| likewise with \lean|cs|, for \lean|ZB cs| the inputs get $n-|cs|$ empty chains then
  \lean|cs|, for \lean|Bun m b| every chain of \lean|b| gets \lean|m| prepended, and every other
  case gives $n$ empty chains.
\item \lean|copyLegs μ d i cs vs| is \lean|zipWith (fun c v => digit (d ^ (c.map μ).prod) i v) cs vs|:
  digit $i$ of each value in base $d$ to the multiplicities of its chain.
  \lean|copyOuts E b d i os| and \lean|copyIns E b d i is| apply it with \lean|b.outCs os.length|,
  \lean|b.inCs is.length|.
\end{itemize}

\begin{definition}[Tensor of a kind]\label{def:lean-core-tensor}
\lean|tensor E : Kind → ℕ → ℕ → List ℕ → List ℕ → ℂ| takes a kind, a dimension $d$, a phase index
$p$ and output and input values:
\begin{center}\small
\begin{tabular}{@{}ll@{}}
\toprule
Kind and arity & Entry \\
\midrule
\lean|Z| & $\mathrm{tZ}(d,\varphi_p,o\mathbin{+\!\!+}i)$, with $\varphi_p=$ \lean|E.phase p| \\
\lean|X| & $\mathrm{tX}(d,\varphi_p,o;i)$ \\
\lean|F|, one output $o$, one input $i$ & $F_d[o,i]$ \\
\lean|T|, \lean|Ti|, one output, one input & $\mathrm{tT}(o,i)$, $\mathrm{tTi}(o,i)$ \\
\lean|W|, one input $i$ & $\mathrm{tW}(d,o,i)$ \\
\lean|I|, one output $a$, one input $b$ & $\mathrm{tI}(d,a,b)$ \\
\lean|ZW|, \lean|XW| & $\mathrm{tZw}(d,c_p,o\mathbin{+\!\!+}i)$, $\mathrm{tXw}(d,c_p,o;i)$, $c_p=$ \lean|E.coef p| \\
\lean|ZB cs| & $\mathrm{tZ}(d,\varphi_p,o\mathbin{+\!\!+}\mathrm{bundled}(\mu,d,i,cs))$ \\
\lean|XB co ci| & $\mathrm{tX}(d,\varphi_p,\mathrm{bundled}(\mu,d,o,co);\mathrm{bundled}(\mu,d,i,ci))$ \\
\lean|WB cs|, one input $w$ & $\mathrm{tW}(d,\mathrm{bundled}(\mu,d,o,cs),w)$ \\
\lean|Bun m b| & $\prod_{j<\mu(m)}\mathrm{tensor}(E,b,d,p,\mathrm{copyOuts}(E,b,d,j,o),\mathrm{copyIns}(E,b,d,j,i))$ \\
\lean|ZP dx dy c|, legs $o\mathbin{+\!\!+}i=[x,y]$ & $\mathrm{tP}(d,c,[x,y])=\mathrm{ex}(c\,x\,y/d)$ ($0$ for any other leg count) \\
any other case (wrong arity, \lean|B|, \lean|S|) & $0$ \\
\bottomrule
\end{tabular}
\end{center}
\end{definition}

\begin{definition}[Node value]\label{def:lean-core-val}
\lean|val E σ n| is the entry of node $n$ under the label assignment $\sigma:\N\to\N$. With
$d=$ \lean|E.dimv n.dim|, $o=$ \lean|n.outs.map σ|, $i=$ \lean|n.ins.map σ|, it agrees with
\lean|tensor E n.kind d n.ph o i| on every kind except the binder and splitter, which only
\lean|val| evaluates: \lean|B| with $o=(k)$, $i=(a,b)$ is $\mathrm{tB}(d,\mathrm{dimv}(n.ph),k,a,b)$,
and \lean|S| with $o=(a,b)$, $i=(k)$ is $\mathrm{tB}(d,\mathrm{dimv}(n.ph),k,a,b)$; every other arity
of \lean|B| or \lean|S| is $0$. So for \lean|B| and \lean|S| the field \lean|dim| indexes the first
factor dimension $s$ and \lean|ph| the second factor dimension $t$; the product leg has dimension
$st$. Copies of a binder or splitter (\lean|Bun m B|) evaluate to $0$.
\end{definition}

\subsection{Evaluation and equality}\label{sec:lean-core-net-eval}

\begin{definition}[Sums over labels]
For $D:\N\to\N$, a label list $ls$, an integrand $f:(\N\to\N)\to\C$ and an assignment $\sigma$,
\[
  \mathrm{sumOver}(D,[\,],f,\sigma)=f(\sigma),\qquad
  \mathrm{sumOver}(D,l{::}ls,f,\sigma)=\sum_{j<D(l)}\mathrm{sumOver}\bigl(D,ls,f,\sigma[l\mapsto j]\bigr).
\]
A repeated label is summed once per occurrence, the inner occurrence overwriting the outer.
\end{definition}

\begin{definition}[Evaluation]\label{def:lean-core-eval}
\lean|prodVal E ns σ| is $\prod_{n\in ns}\mathrm{val}(E,\sigma,n)$ and
\[
  \mathrm{eval}(E,N,\sigma)=\mathrm{sumOver}\bigl(\mathrm{labDim}\,E,\;N.\mathrm{int},\;
  \mathrm{prodVal}(E,N.\mathrm{nodes}),\;\sigma\bigr).
\]
The boundary values are read from $\sigma$; $\sigma$ on labels in neither \lean|int| nor
\lean|bnd| is read as well, so a label that appears on a node but in neither list acts as a fixed
boundary value.
\end{definition}

\begin{definition}[Equality of scaled networks]\label{def:lean-core-equal}
\lean|InRange E ls σ| holds when $\sigma(l)<\mathrm{labDim}(E,l)$ for every $l\in ls$.
\begin{lstlisting}[style=leancore]
def Equal (E : Env) (s₁ : ℂ) (N₁ : Net) (s₂ : ℂ) (N₂ : Net) : Prop :=
  N₁.bnd.Perm N₂.bnd ∧ ∀ σ, InRange E N₁.bnd σ → s₁ * eval E N₁ σ = s₂ * eval E N₂ σ
\end{lstlisting}
That is, $s_1\cdot N_1=s_2\cdot N_2$ as tensors: the boundaries are permutations of each other and
the scaled entries agree at every in-range boundary assignment.
\end{definition}

\begin{lemma}[Equality is an equivalence, and scalars chain]
\begin{itemize}
\item \lean|InRange.perm|: \lean|a.Perm b → InRange E a σ → InRange E b σ|.
\item \lean|Equal.refl E s N|, \lean|Equal.symm|, \lean|Equal.trans| (same middle scalar).
\item \lean|Equal.mul_trans|: from $a X=bY$ and $cY=eZ$, $(ac)X=(be)Z$.
\item \lean|Equal.rescale|: from $aX=bY$, $a\ne0$ and $sb=ta$, $sX=tY$.
\end{itemize}
\end{lemma}

\subsection{Lemmas about sums}\label{sec:lean-core-net-sums}

Fix $D:\N\to\N$. A function $f$ \term{reads} the labels $S$ (\lean|Reads S f|) when
$f(\tau)=f(\tau')$ for all $\tau,\tau'$ that agree on $S$.
\begin{itemize}
\item \lean|sumOver_congr|: if $f(\tau)=g(\tau)$ for every $\tau$ that agrees with $\sigma$ off $ls$
  and has $\tau(x)<D(x)$ on $ls$, the two sums over $ls$ agree.
\item \lean|sumOver_ext|: if $f$ reads $S$ and $\sigma,\sigma'$ agree on $S\setminus ls$, the sums at
  $\sigma$ and $\sigma'$ agree.
\item \lean|sumOver_append|: $\mathrm{sumOver}(a\mathbin{+\!\!+}b,f)=\mathrm{sumOver}(a,\mathrm{sumOver}(b,f))$.
\item \lean|mul_sumOver|: $k\cdot\mathrm{sumOver}(ls,f,\sigma)=\mathrm{sumOver}(ls,k f,\sigma)$.
\item \lean|sumOver_const_mul|: if $g$ reads $S$ and $S\cap ls=\emptyset$,
  $\mathrm{sumOver}(ls,g f,\sigma)=g(\sigma)\,\mathrm{sumOver}(ls,f,\sigma)$.
\item \lean|sumOver_perm|: a permutation of $ls$ leaves the sum unchanged (by Fubini,
  \lean|Finset.sum_comm|, on each adjacent swap).
\item \lean|sumOver_map|: for an injective $f$ with $D(f(x))=D(x)$ on $ls$,
  $\mathrm{sumOver}(ls.\mathrm{map}\,f,\tau\mapsto g(\tau\circ f),\sigma)=\mathrm{sumOver}(ls,g,\sigma\circ f)$.
\end{itemize}
Node lemmas: \lean|val_reads| (a node's value reads its labels), \lean|prodVal_reads| (reads
\lean|labelsOf ns|), \lean|prodVal_append| (product over $a\mathbin{+\!\!+}b$),
\lean|val_rename| ($\mathrm{val}(\sigma,n.\mathrm{rename}\,f)=\mathrm{val}(\sigma\circ f,n)$) and
\lean|prodVal_rename|.

\subsection{Local equations and the frame theorem}\label{sec:lean-core-net-frame}

\begin{definition}[Local equation]\label{def:lean-core-localeq}
\begin{lstlisting}[style=leancore]
def LocalEq (E : Env) (L : List Node) (iL : List ℕ) (R : List Node) (iR : List ℕ)
    (iface : List ℕ) (kL kR : ℂ) : Prop :=
  ∀ τ, InRange E iface τ →
    kL * sumOver (labDim E) iL (prodVal E L) τ = kR * sumOver (labDim E) iR (prodVal E R) τ
\end{lstlisting}
In words: the network fragment $L$ with internal labels $iL$, scaled by $k_L$, equals the fragment
$R$ with internal labels $iR$, scaled by $k_R$, at every assignment in range on the interface labels
\lean|iface|. Every rule lemma of this library concludes \lean|LocalEq E L iL R iR iface kL kR|;
most have $k_L=1$, so they read $L=k_R\cdot R$.
\end{definition}

\begin{theorem}[\lean|frame|]\label{thm:lean-core-frame}
Let $C$ be a context node list with summed labels $iC$, and suppose
\lean|LocalEq E L iL R iR iface kL kR|, no label of $C$ lies in $iL$ or $iR$, and every interface
label lies in $iC$ or in \lean|bnd|. Then for every $\sigma$ in range on \lean|bnd|,
\[
  k_L\,\mathrm{eval}(E,\langle C\mathbin{+\!\!+}L,\;iC\mathbin{+\!\!+}iL,\;bnd\rangle,\sigma)
  =k_R\,\mathrm{eval}(E,\langle C\mathbin{+\!\!+}R,\;iC\mathbin{+\!\!+}iR,\;bnd\rangle,\sigma).
\]
\end{theorem}
\emph{Proof technique.} For $X\in\{L,R\}$ both sides split as
\[\mathrm{sumOver}\bigl(iC,\ \tau\mapsto
\mathrm{prodVal}(C,\tau)\cdot\mathrm{sumOver}(iX,\mathrm{prodVal}(X),\tau)\bigr)\]
by \lean|sumOver_append| and \lean|sumOver_const_mul| (the context reads no label of $iX$); inside the
outer sum every interface label is in range (summed in $iC$ or fixed on the boundary), so the local
equation applies pointwise.

\begin{theorem}[\lean|rewrite|]\label{thm:lean-core-rewrite}
Under the hypotheses of \lean|frame| with $k_L=1$, $k_R=k$, and $s'=s k$:
\lean|Equal E s ⟨C ++ L, iC ++ iL, bnd⟩ s' ⟨C ++ R, iC ++ iR, bnd⟩|.
\end{theorem}

\begin{example}[A local equation in context]
The cap lemma (\cref{sec:lean-core-zxrules}) states
\lean|LocalEq E [⟨.Z,[w],[],dz,a⟩, ⟨.X,[],[w],dz,b⟩] [w] [] [] iface 1 (rt d)|: a phaseless Z
state on label $w$ followed by a phaseless X effect on $w$, summed over $w$, is the empty network
times $\sqrt d$. With a context $C$ that does not mention $w$, \lean|rewrite| with $k=\sqrt d$ turns
$s\cdot(C\mathbin{+\!\!+}L)$ into $(s\sqrt d)\cdot C$, the summed label list losing $w$.
\end{example}

%======================================================================
\section{Isomorphic networks (\texttt{Iso})}\label{sec:lean-core-iso}

\texttt{Iso} decides that one network is another up to a bijective renaming of its summed labels,
a reordering of its nodes, and the port symmetries of spiders.

\begin{definition}[Renaming]
\lean|ren π x| for a list of pairs $\pi$ is the second component of the first pair whose first
component is $x$, and $x$ itself when no pair has first component $x$. \lean|ren_of_not_mem| states
the latter case; \lean|ren_injective| states that \lean|ren π| is injective when the keys
$K=\pi.\mathrm{map}\,\mathrm{fst}$ are distinct and $K.\mathrm{map}(\mathrm{ren}\,\pi)$ is a
permutation of $K$.
\end{definition}

\begin{definition}[Node equivalence]
\lean|nodeEquiv n m| is true when $n$ and $m$ have equal kind, dimension index and phase index and
\begin{itemize}
\item for \lean|Z| and \lean|ZW|: \lean|n.outs ++ n.ins| is a permutation of \lean|m.outs ++ m.ins|;
\item for \lean|X| and \lean|XW|: the outputs are permutations of each other and so are the inputs;
\item for every other kind: the output lists are equal and the input lists are equal.
\end{itemize}
\lean|val_nodeEquiv| proves that equivalent nodes have equal values under every assignment.
\end{definition}

\begin{definition}[\lean|isoOK|]
Let \lean|Node.none| be the placeholder node \lean|⟨.I, [], [], 0, 0⟩|, $f=\mathrm{ren}\,\pi$, $K$ the keys of $\pi$,
$A'=A.\mathrm{nodes}.\mathrm{map}(\mathrm{rename}\,f)$ and $P=p.\mathrm{map}(j\mapsto
A'.\mathrm{getD}\,j\ \mathrm{none})$, \lean|isoOK labs A B π p| is the conjunction of:
\begin{enumerate}
\item $K$ has no duplicates;
\item $K.\mathrm{map}\,f$ is a permutation of $K$;
\item every label on a node of $A$ is in $A.\mathrm{int}$ or not in $K$ (the renaming touches only
  summed labels);
\item every $x\in A.\mathrm{int}$ has \lean|labIdx labs (f x) == labIdx labs x|;
\item $A.\mathrm{bnd}=B.\mathrm{bnd}$ (as lists);
\item $B.\mathrm{int}$ is a permutation of $A.\mathrm{int}.\mathrm{map}\,f$;
\item $P$ is a permutation of $A'$ (so $p$ is a permutation of node positions);
\item $|P|=|B.\mathrm{nodes}|$;
\item \lean|nodeEquiv| holds for each pair of \lean|List.zip P B.nodes|.
\end{enumerate}
The certificate is $(\pi,p)$: the label renaming and the position in $A'$ of each node of $B$.
\end{definition}

\begin{theorem}
\lean|prodVal_equiv|: lists of equal length that are pairwise \lean|nodeEquiv| have equal products.
\lean|iso_eval|: \lean|isoOK E.labs A B π p = true → eval E A σ = eval E B σ| for every $\sigma$.
\lean|iso|: under the same hypothesis \lean|Equal E s A s B| for every $s$; \lean|iso_scalar| adds a
proof $s=s'$ and gives \lean|Equal E s A s' B|.
\end{theorem}
\emph{Proof technique.} \lean|iso_eval| permutes the summed labels of $B$ into $A.\mathrm{int}.\mathrm{map}\,f$
(\lean|sumOver_perm|), rewrites the product of $B$ as the product of the renamed nodes of $A$,
pulls the renaming out with \lean|sumOver_map|, and closes with \lean|sumOver_ext| because $f$ fixes
every label of $A$ outside $A.\mathrm{int}$.

%======================================================================
\section{Generator facts and closed forms (\texttt{Rules}, \texttt{Closed})}\label{sec:lean-core-closed}

\subsection{Phases, roots of unity and the Fourier box}

\texttt{Rules} opens with facts about the scalars (\texttt{Rules} also imports
the Mathlib module \nolinkurl{Mathlib.RingTheory.RootsOfUnity.Complex}):
\begin{itemize}
\item \lean|ex_zero|: $\mathrm{ex}(0)=1$; \lean|ex_add|: $\mathrm{ex}(a+b)=\mathrm{ex}(a)\mathrm{ex}(b)$;
  \lean|ex_int|: $\mathrm{ex}(a+z)=\mathrm{ex}(a)$ for $z\in\Z$; \lean|ex_congr_int|: $x=y+z$ with
  $z\in\Z$ gives $\mathrm{ex}(x)=\mathrm{ex}(y)$.
\item \lean|om_prim|: for $d>0$, $\omega_d$ is a primitive $d$-th root of unity; \lean|conj_om|:
  $\overline{\omega_d}=\omega_d^{-1}$; \lean|conj_sqrt|, \lean|sqrt_mul_sqrt|: $\sqrt d$ is real and
  $\sqrt d\sqrt d=d$.
\item \lean|fou_unitary| ($d>0$, $k,k'<d$): $\sum_{j<d}F_d[j,k]\overline{F_d[j,k']}=[k=k']$.
  \emph{Technique:} each term is $u^j/d$ with $u=\omega^k\omega^{-k'}$; for $k\ne k'$, $u\ne1$ and
  $u^d=1$, so the geometric sum vanishes.
\item \lean|fou_symm|: $F_d[j,k]=F_d[k,j]$.
\item \lean|tZ_at|: if $j$ is one of the values $vs$, $\mathrm{tZ}(d,\varphi,vs)=[j<d\land\forall v\in vs,
  v=j]\,\mathrm{ex}(\varphi(j))$; \lean|tZ_perm|: $\mathrm{tZ}$ is invariant under permutation of its
  values.
\end{itemize}

\texttt{Closed} defines \lean|rt d = (Real.sqrt d : ℂ)| (an \lean|abbrev|) and proves:
\begin{itemize}
\item \lean|rt_ne| ($\sqrt d\ne0$ for $d>0$), \lean|rt_cancel|
  ($\sqrt d\,(\sqrt d)^{-n}(\sqrt d)^n=\sqrt d$), \lean|nat_inv_mul| ($d^{-1}d=1$), \lean|rt_sq|,
  \lean|rt_sq'| ($\sqrt d^2=d$), \lean|om_ne_zero|, \lean|om_pow_self| ($\omega_d^d=1$),
  \lean|om_pow_mod| ($\omega_d^n=\omega_d^{n\bmod d}$), \lean|om_pow_eq_iff|
  ($\omega_d^a=\omega_d^b\iff a\equiv b\pmod d$).
\item \lean|om_geom|: for $0<m$, $1<p$, $p\mid m$: $\sum_{j<p}\omega_m^{j(m/p)}=0$.
\item \lean|char_sum| (the character sum): $\sum_{k<d}\omega_d^{ka}\,\overline{\omega_d}^{kb}=d\,[a\equiv
  b\pmod d]$.
\item \lean|prod_fou|: $\prod_{o\in os}F_d[o,k]=\omega_d^{k\sum os}\,(\sqrt d)^{-|os|}$, and
  \lean|prod_conj_fou| the conjugate form.
\item \lean|tX_flat| (a phaseless X spider is a scaled delta of sums): if
  $\mathrm{ex}(\varphi(k))=1$ for all $k<d$,
  \[\mathrm{tX}(d,\varphi,os;is)=(\sqrt d)^{-(|os|+|is|)}\;d\;\Bigl[\textstyle\sum os\equiv\sum is\pmod d\Bigr].\]
\item \lean|tZ_one|: a phaseless Z spider on one value $x$ is $[x<d]$.
\item \lean|fou_sq|: $\sum_{w<d}F_d[x,w]F_d[w,y]=[(x+y)\bmod d=0]$: the Fourier box squared is the
  antipode (dualiser) $\ket{x}\mapsto\ket{-x}$.
\item \lean|sum_mod_eq|: $\sum_{w<d}[s\equiv w]\,f(w)=[0<d]\,f(s\bmod d)$.
\end{itemize}

\subsection{Bare wires}

\begin{theorem}[\lean|wire_elim|]\label{thm:lean-core-wireelim}
Let a network have nodes $ns_1\mathbin{+\!\!+}\langle\mathtt{I},[a],[b],dz,0\rangle{::}ns_2$, summed
labels \lean|int| and boundary \lean|bnd|, with $b\in\mathit{int}$, $a\ne b$, $a\in\mathit{int}\cup\mathit{bnd}$
and both $a$ and $b$ of dimension index $dz$. Then for every $s$
\[
  \mathtt{Equal}\;E\;s\;\langle ns_1\mathbin{+\!\!+}\mathtt{I}_{a,b}{::}ns_2,\;\mathit{int},\;\mathit{bnd}\rangle
  \;s\;\langle(ns_1\mathbin{+\!\!+}ns_2).\mathrm{map}(\mathrm{rename}(\mathrm{ren}[(b,a)])),\;
  \mathit{int}.\mathrm{erase}\,b,\;\mathit{bnd}\rangle:
\]
the wire is removed and its summed end $b$ is renamed to $a$ everywhere.
\end{theorem}

\section{Spider and ZX rules (\texttt{Rules}, \texttt{ZXRules}, \texttt{Bialgebra}, \texttt{Hopf})}\label{sec:lean-core-zxrules}

Each lemma below is a \lean|LocalEq| (\cref{def:lean-core-localeq}); the statement lists the left
nodes, the left summed labels, the right nodes and summed labels, and the scalars. A node is written
$\langle\mathit{kind},\mathit{outs},\mathit{ins},\mathit{dz},\mathit{ph}\rangle$, $d$ abbreviates
\lean|E.dimv dz|, and ``phaseless $a$'' abbreviates the hypothesis
\lean|∀ k < E.dimv dz, ex (E.phase a k) = 1|. Unless stated otherwise every named label has
dimension index $dz$ (hypotheses \lean|labIdx E.labs x = dz|), every boundary label of a fragment
lies in \lean|iface|, distinct names denote distinct labels where the lemma has a hypothesis saying
so, and $d>0$ where the proof divides by $d$. These hypotheses are listed in each Lean statement;
the summaries give the equation. \Cref{fig:lean-core-rules} draws the main ones.

\begin{figure}[htbp]
\centering
\begin{tabular}{@{}c@{\qquad}c@{\qquad}c@{}}
\begin{tikzpicture}[zxd]
\node[zs] (a) at (0,0) {$a$}; \node[zs] (b) at (0,1) {$b$};
\draw (a) -- node[right,dimlabel]{$w$} (b);
\draw (a) -- (-0.4,-0.7); \draw (a) -- (0.4,-0.7);
\draw (b) -- (-0.4,1.7); \draw (b) -- (0.4,1.7);
\end{tikzpicture}
$=$
\begin{tikzpicture}[zxd]
\node[zs] (c) at (0,0.5) {$c$};
\draw (c) -- (-0.4,-0.2); \draw (c) -- (0.4,-0.2);
\draw (c) -- (-0.4,1.2); \draw (c) -- (0.4,1.2);
\end{tikzpicture}
&
\begin{tikzpicture}[zxd]
\node[zdot] (z) at (0,0) {}; \node[xdot] (x) at (0,1) {};
\draw (z) -- node[right,dimlabel]{$w$} (x);
\end{tikzpicture}
$=\sqrt d$
&
\begin{tikzpicture}[zxd]
\node[xdot] (x) at (0,0) {}; \node[zdot] (z) at (0,0.8) {};
\draw (x) -- (z); \draw (z) -- (-0.4,1.5); \draw (z) -- (0.4,1.5);
\end{tikzpicture}
$=k$
\begin{tikzpicture}[zxd]
\node[xdot] (x1) at (-0.3,0.6) {}; \node[xdot] (x2) at (0.3,0.6) {};
\draw (x1) -- (-0.3,1.3); \draw (x2) -- (0.3,1.3);
\end{tikzpicture}
\\[2pt]
\small\lean|fusionZ| ($e(c)=e(a)e(b)$) & \small\lean|cap| & \small\lean|stateCopy| ($k\sqrt d^{\,n}=\sqrt d$) \\[8pt]
\begin{tikzpicture}[zxd]
\node[xdot] (x) at (0,0) {}; \node[zdot] (z) at (0,0.8) {};
\draw (x) -- (z); \draw (x) -- (-0.4,-0.6); \draw (x) -- (0.4,-0.6);
\draw (z) -- (-0.4,1.4); \draw (z) -- (0.4,1.4);
\end{tikzpicture}
$=\sqrt d$
\begin{tikzpicture}[zxd]
\node[zdot] (z1) at (-0.4,0) {}; \node[zdot] (z2) at (0.4,0) {};
\node[xdot] (x1) at (-0.4,0.9) {}; \node[xdot] (x2) at (0.4,0.9) {};
\draw (z1) -- (x1); \draw (z1) -- (x2); \draw (z2) -- (x1); \draw (z2) -- (x2);
\draw (z1) -- (-0.4,-0.6); \draw (z2) -- (0.4,-0.6);
\draw (x1) -- (-0.4,1.5); \draw (x2) -- (0.4,1.5);
\end{tikzpicture}
&
\begin{tikzpicture}[zxd]
\node[zs] (z) at (0,0) {}; \node[xs] (x) at (0,1.6) {};
\node[fbox] (f1) at (0.5,0.55) {$F$}; \node[fbox] (f2) at (0.5,1.05) {$F$};
\draw (z) to[bend left=30] (x);
\draw (z) -- (f1); \draw (f1) -- (f2); \draw (f2) -- (x);
\draw (z) -- (0,-0.6); \draw (x) -- (0,2.2);
\end{tikzpicture}
$=\dfrac1d$
\begin{tikzpicture}[zxd]
\node[zs] (z) at (0,0) {}; \node[xs] (x) at (0,1.2) {};
\draw (z) -- (0,-0.6); \draw (x) -- (0,1.8);
\end{tikzpicture}
&
\begin{tikzpicture}[zxd]
\node[tri] (t) at (0,0) {}; \node[trii] (ti) at (0,0.8) {};
\draw (0,-0.6) -- (t); \draw (t) -- (ti); \draw (ti) -- (0,1.4);
\end{tikzpicture}
$=$
\begin{tikzpicture}[zxd]
\draw (0,-0.6) -- (0,1.4);
\end{tikzpicture}
\\[2pt]
\small\lean|bialgebra| & \small\lean|hopf| ($kd=1$) & \small\lean|triangle| \\[8pt]
\begin{tikzpicture}[zxd]
\node[wnode] (w1) at (0,0) {}; \node[wnode] (w2) at (0.4,0.8) {};
\draw (0,-0.6) -- (w1); \draw (w1) -- (-0.4,1.5); \draw (w1) -- (w2);
\draw (w2) -- (0.1,1.5); \draw (w2) -- (0.7,1.5);
\end{tikzpicture}
$=$
\begin{tikzpicture}[zxd]
\node[wnode] (w) at (0,0.4) {};
\draw (0,-0.3) -- (w); \draw (w) -- (-0.4,1.2); \draw (w) -- (0,1.2); \draw (w) -- (0.4,1.2);
\end{tikzpicture}
&
\begin{tikzpicture}[zxd]
\node[snode] (s) at (0,0) {$S$}; \node[bnode] (b) at (0,1) {$B$};
\draw (0,-0.6) -- (s); \draw (s.north west) -- (b.south west); \draw (s.north east) -- (b.south east);
\draw (b) -- (0,1.6);
\end{tikzpicture}
$=$
\begin{tikzpicture}[zxd]
\draw (0,-0.6) -- (0,1.6);
\end{tikzpicture}
&
\begin{tikzpicture}[zxd]
\node[fbox] (f1) at (0,0) {$F$}; \node[fbox] (f2) at (0,0.55) {$F$};
\node[fbox] (f3) at (0,1.1) {$F$}; \node[fbox] (f4) at (0,1.65) {$F$};
\draw (0,-0.5) -- (f1) -- (f2) -- (f3) -- (f4) -- (0,2.15);
\end{tikzpicture}
$=$
\begin{tikzpicture}[zxd]
\draw (0,-0.5) -- (0,2.15);
\end{tikzpicture}
\\[2pt]
\small\lean|wFusion| & \small\lean|connInverse| & \small\lean|fourier4|
\end{tabular}
\caption{Rule lemmas of \cref{sec:lean-core-zxrules,sec:lean-core-wrules} as pictures (inputs at the
bottom, outputs at the top; small dots are phaseless spiders).}\label{fig:lean-core-rules}
\end{figure}

\subsection{Fusion and identity (\texttt{Rules})}

\begin{lemma}[\lean|fusionZ|]
Hypotheses: $w\in o_1\mathbin{+\!\!+}i_1$, $w\in o_2\mathbin{+\!\!+}i_2$, $o_3\mathbin{+\!\!+}i_3$ is a
permutation of $(o_1\mathbin{+\!\!+}i_1\mathbin{+\!\!+}o_2\mathbin{+\!\!+}i_2)$ with every occurrence of $w$
removed, $w$ of index $dz$, and $\mathrm{ex}(\varphi_c(k))=\mathrm{ex}(\varphi_a(k))\,\mathrm{ex}(\varphi_b(k))$ for
$k<d$. Conclusion:
\[\mathtt{LocalEq}\;[\langle Z,o_1,i_1,dz,a\rangle,\langle Z,o_2,i_2,dz,b\rangle]\;[w]\;
[\langle Z,o_3,i_3,dz,c\rangle]\;[\,]\;\mathit{iface}\;1\;1.\]
\end{lemma}

\begin{lemma}[\lean|fusionX|]
Hypotheses: $d>0$; $w$ occurs exactly once in $o_1$ and not in $i_1$; exactly once in $i_2$ and not in
$o_2$; $o_3\sim o_1.\mathrm{erase}\,w\mathbin{+\!\!+}o_2$ and $i_3\sim i_1\mathbin{+\!\!+}i_2.\mathrm{erase}\,w$
($\sim$ is \lean|List.Perm|); phases multiply as in \lean|fusionZ|. Conclusion: two X spiders joined
by $w$ (an output of the first, an input of the second) fuse into
$\langle X,o_3,i_3,dz,c\rangle$, scalars $1,1$. \emph{Technique:} each spider is a sum over its
index $k$; the sum over $w$ of $F[w,k]\overline{F[w,k']}$ is $[k=k']$ by \lean|fou_unitary|.
\end{lemma}

\lean|map_update_absent| (setting a label absent from a list leaves its values) and
\lean|prod_update_once| (a product over a list in which $w$ occurs once, at $w\mapsto j$, is $g(j)$
times the product over the rest) serve the X proofs.

\begin{lemma}[\lean|identityZ|, \lean|identityX|]
With $(p,q)=(w,x)$ or $(p,q)=(x,w)$ and $a$ phaseless:
$\langle Z,[p],[q],dz,a\rangle=\langle I,[x],[w],dz,0\rangle$ (no summed labels, scalars $1,1$).
The X version adds $d>0$ and $p,q\in\mathit{iface}$ of index $dz$; it follows from
\lean|fou_unitary| and \lean|fou_symm|.
\end{lemma}

\subsection{\texttt{ZXRules}}

Helpers: \lean|labDim_of| ($\mathrm{labIdx}=dz\Rightarrow\mathrm{labDim}=\mathrm{dimv}\,dz$),
\lean|rt_inv_mul| ($(\sqrt d)^{-1}d=\sqrt d$), \lean|sum_range_ite_mod|, \lean|sum_range_ite_mod'|,
\lean|sum_mod0|, \lean|sum_mod0'| (sums of indicators of residues), \lean|prod_ite_zero|
($\prod_{o\in l}[\tau o=0]c=[\forall o, \tau o=0]c^{|l|}$), \lean|char_sum0|
($\sum_{k<d}\omega^{ka}=d[a\equiv0]$), \lean|sum_fou_left|, \lean|sum_fou_right|
($\sum_jF[x,j]=(\sqrt d)^{-1}d[x\equiv0]$), \lean|fou_zero|, \lean|fou_zero'| ($F[x,0]=F[0,x]=(\sqrt
d)^{-1}$), \lean|tri_sum|, \lean|tri_sum'| ($\sum_j\mathrm{tT}(j,i)\mathrm{tTi}(o,j)=[o=i]$ and the
reverse, for $o,i<d$), \lean|antipode_iff| ($(j+a)\bmod d=0\iff j=(d-a)\bmod d$ for $a,j<d$),
\lean|antipode_inj|, \lean|antipode_sq| (the antipode squared is the identity).

\begin{description}[style=nextline,font=\normalfont\ttfamily]
\item[cap, capSwapped] $[\langle Z,[w],[\,],dz,a\rangle,\langle X,[\,],[w],dz,b\rangle]$, summed $[w]$,
  equals the empty fragment times $\sqrt d$ ($k_L=1$, $k_R=\sqrt d$); $a$, $b$ phaseless, $d>0$. The
  swapped form exchanges the colors: $\langle X,[w],[\,]\rangle,\langle Z,[\,],[w]\rangle$.
\item[colorState, colorEffect] $\langle Z,[w],[\,],dz,a\rangle$ then $\langle F,[x],[w],dz,0\rangle$,
  summed $[w]$, is $\langle X,[x],[\,],dz,b\rangle$; the effect form is
  $\langle Z,[\,],[w]\rangle,\langle F,[w],[x]\rangle\mapsto\langle X,[\,],[x]\rangle$; $a,b$ phaseless,
  $w\ne x$, scalars $1,1$.
\item[colorStateX, colorEffectX] The same with X and Z exchanged; these add $x\in\mathit{iface}$ of
  index $dz$.
\item[stateCopy] $[\langle X,[w],[\,],dz,a\rangle,\langle Z,os,[w],dz,b\rangle]$, summed $[w]$, with
  $w\notin os$ and every $o\in os$ in \lean|iface| of index $dz$, equals
  $os.\mathrm{map}(o\mapsto\langle X,[o],[\,],dz,c\rangle)$, scalars $1$ and $k$ where
  $k\,\sqrt d^{\,|os|}=\sqrt d$; $a,b,c$ phaseless. \lean|stateCopySwapped| exchanges colors.
\item[triangle, triangle'] $[\langle T,[w],[i]\rangle,\langle Ti,[o],[w]\rangle]$, summed $[w]$, is
  $\langle I,[x],[y],dz,0\rangle$ with $\{x,y\}=\{o,i\}$ in either order; $w\ne o$, $w\ne i$;
  \lean|triangle'| has $\mathtt{Ti}$ first: $[\langle Ti,[w],[i]\rangle,\langle T,[o],[w]\rangle]$.
\item[fourier4] Four Fourier boxes in series $a\to w_1\to w_2\to w_3\to b$, summed $[w_1,w_2,w_3]$,
  are $\langle Z,[b],[a],dz,c\rangle$ with $c$ phaseless ($F^4=\mathrm{id}$), scalars $1,1$; all six
  labels distinct as listed in the hypotheses.
\end{description}

\subsection{\texttt{Bialgebra}}

Helpers \lean|tZ_two|, \lean|tZ_first|: a phaseless Z spider on $[x,y,z]$ (resp.\ $[z,x,y]$) with
$z<d$ is $[x=z\land y=z]$.

\begin{lemma}[\lean|bialgebra|]
With $d>0$, $w$ and $n_{00},n_{01},n_{10},n_{11}$ of index $dz$, $[n_{00},n_{01},n_{10},n_{11},a,b,c,e]$
and $[w,a,b,c,e]$ without duplicates, $a,b,c,e\in\mathit{iface}$ of index $dz$, and all six phases
phaseless:
\[\begin{aligned}
&[\langle X,[w],[a,b],dz,p_x\rangle,\ \langle Z,[c,e],[w],dz,p_z\rangle]\ \text{summed }[w]\\
&\quad=\sqrt d\cdot[\langle Z,[n_{00},n_{01}],[a]\rangle,\langle Z,[n_{10},n_{11}],[b]\rangle,
\langle X,[c],[n_{00},n_{10}]\rangle,\langle X,[e],[n_{01},n_{11}]\rangle]\\
&\qquad\text{summed }[n_{00},n_{01},n_{10},n_{11}].
\end{aligned}\]
\lean|bialgebraSwapped| is the same with every color exchanged. \emph{Technique:} both sides reduce
by \lean|tZ_two| to indicator sums; \lean|tX_flat| writes each X spider as a delta of sums, and the
scalar $\sqrt d$ balances $(\sqrt d)^{-3}d$ against $((\sqrt d)^{-3}d)^2$.
\end{lemma}

\subsection{\texttt{Hopf}}

Helpers: \lean|tX_perm| (permuting outputs and inputs of $\mathrm{tX}$), \lean|om_pow_zero_mod|
($s\equiv0\Rightarrow\omega^{sm}=1$), \lean|tX_antipode_in| and \lean|tX_antipode_out| (two inputs,
resp.\ outputs, carrying values with $j+j'\equiv0$ drop out of an X spider as the factor $d^{-1}$),
\lean|tZ_sum_cons| ($\sum_j\mathrm{tZ}(j{::}r)=\mathrm{tZ}(r)$), \lean|tZ_cons_cons|
($\mathrm{tZ}(j{::}j'{::}r)=[j'=j]\,\mathrm{tZ}(j{::}r)$), \lean|antipode_collapse|
($\sum_m[(m+j)\equiv0]g(m)=g((d-j)\bmod d)$), \lean|antipode_spec|.

\begin{lemma}[\lean|hopf|]
Let the Z spider $\langle Z,zo,zi,dz,p_z\rangle$ have legs $zo\mathbin{+\!\!+}zi\sim u{::}v_1{::}(zo'\mathbin{+\!\!+}zi')$,
the X spider $\langle X,xo,xi,dz,p_x\rangle$ have $xo\sim xo'$ and $xi\sim u{::}v_3{::}xi'$, and let
$\langle F,[v_2],[v_1]\rangle$, $\langle F,[v_3],[v_2]\rangle$ carry $v_1$ to $v_3$. With
$[u,v_1,v_2,v_3]$ distinct, of index $dz$, absent from the primed lists, $d>0$ and $k\,d=1$:
the four nodes summed over $[u,v_1,v_2,v_3]$ equal $k$ times
$[\langle Z,zo',zi',dz,p_z\rangle,\langle X,xo',xi',dz,p_x\rangle]$ with no summed label. The phases
are arbitrary. \lean|hopfSwapped| puts the X spider first, with its outputs carrying $u,v_1$ and the
Z spider's legs carrying $u,v_3$. \emph{Technique:} the Fourier pair is the antipode
(\lean|fou_sq|); after summing $v_2$ the X spider has two inputs with opposite values, which
\lean|tX_antipode_in| removes as $d^{-1}$, and the Z spider's two legs on $u$, $v_1$ collapse by
\lean|tZ_cons_cons| and \lean|tZ_sum_cons|.
\end{lemma}

%======================================================================
\section{W, triangle and connective rules}\label{sec:lean-core-wrules}

\subsection{\texttt{WRules}}

\lean|nz l| (an \lean|abbrev|) is \lean|l.filter (· ≠ 0)|. Facts about the W node:
\lean|nz_nil_iff| ($\mathrm{nz}\,l=[\,]\iff$ all entries are $0$), \lean|tW_eq| (the definition in
terms of $\mathrm{nz}$), \lean|tW_perm| (invariance under permuting outputs), \lean|nz_cons|,
\lean|tW_cons_zero| (a zero output drops out), \lean|tW_cons| (for $I\ne0$:
$\mathrm{tW}(j{::}A,I)=[j=0]\mathrm{tW}(A,I)+[j=I][\mathrm{nz}A=[\,]\land I<d]$),
\lean|tW_zero_in|, \lean|sum_tW_cons| (summing one output against $g$), \lean|tW_append|
($\mathrm{tW}(A\mathbin{+\!\!+}B,I)=\mathrm{tW}(A,I)\mathrm{tW}(B,0)+[I\ne0][\mathrm{nz}A=[\,]]\mathrm{tW}(B,I)$
for $I<d$).

\begin{description}[style=nextline,font=\normalfont\ttfamily]
\item[wFusion] $[\langle W,o_1,[i]\rangle,\langle W,o_2,[w]\rangle]$, summed $[w]$, with $w$ once in
  $o_1$, $w\notin o_2$, $w\ne i$, $i\in\mathit{iface}$ of index $dz$, $d>0$, is
  $\langle W,o_3,[i]\rangle$ with $o_3\sim o_1.\mathrm{erase}\,w\mathbin{+\!\!+}o_2$; scalars $1,1$.
\item[wIdentity] $\langle W,[p],[q]\rangle=\langle I,[x],[w]\rangle$ with $\{p,q\}$ as in
  \lean|identityZ|, $p,q\in\mathit{iface}$ of index $dz$.
\item[wZeroCopy] A phaseless X state $\langle X,[w],[\,],a\rangle$ into $\langle W,os,[w]\rangle$
  is $os.\mathrm{map}(o\mapsto\langle X,[o],[\,],c\rangle)$ times $k$, $k\sqrt d^{\,|os|}=\sqrt d$.
\item[triZeroState, triInvZeroState] A phaseless X state into $\langle T,[o],[w]\rangle$ (resp.\
  $\mathtt{Ti}$) is the X state $\langle X,[o],[\,],a\rangle$; scalars $1,1$.
\item[triZeroEffect] $[\langle T,[w],[i]\rangle,\langle X,[\,],[w],a\rangle]$ summed $[w]$ is
  $\sqrt d\cdot\langle Z,[\,],[i],c\rangle$; $a,c$ phaseless.
\item[wZEffect] $[\langle W,os,[i]\rangle,\langle Z,[\,],[w],a\rangle]$ summed $[w]$, with $w$ once
  in $os$ and $os'\sim os.\mathrm{erase}\,w$, $v\notin os'$, is
  $[\langle T,[v],[i]\rangle,\langle W,os',[v]\rangle]$ summed $[v]$; scalars $1,1$.
\item[wZeroEffect] $[\langle W,os,[i]\rangle,\langle X,[\,],[w],a\rangle]$ summed $[w]$ is
  $\sqrt d\cdot\langle W,os',[i]\rangle$.
\end{description}

\subsection{\texttt{WAssoc}: W associativity over groups}

\begin{definition}
\lean|sumW d M f| sums $f$ over all lists of $M$ values below $d$:
$\mathrm{sumW}(d,0,f)=f([\,])$, $\mathrm{sumW}(d,M{+}1,f)=\sum_{x<d}\mathrm{sumW}(d,M,xs\mapsto f(x{::}xs))$.
\lean|wTerm d w xs Bs| is $\prod_j w(x_j)\,\mathrm{tW}(d,B_j,x_j)$ over the zipped lists (the empty
product $1$ when either list ends).
\end{definition}

\begin{theorem}[\lean|wAssoc|, \lean|wAssoc_nil|]
For $d>0$, $w(0)=1$, $i<d$, all group lists $Bs$ and all $A$:
\[\begin{aligned}
&\mathrm{sumW}\bigl(d,|Bs|,xs\mapsto\mathrm{tW}(d,A\mathbin{+\!\!+}xs,i)\,\mathrm{wTerm}(d,w,xs,Bs)\bigr)\\
&\quad=\mathrm{tW}(A,i)\,\mathrm{tW}(\textstyle\bigsqcup Bs,0)+[i\ne0]\,[\mathrm{nz}A=[\,]]\,w(i)\,\mathrm{tW}(\bigsqcup Bs,i),
\end{aligned}\]
where $\bigsqcup Bs$ is \lean|Bs.flatten|; with $A=[\,]$ the right side is $w(i)\,\mathrm{tW}(\bigsqcup
Bs,i)$. Helpers: \lean|sumW_congr|, \lean|sumW_smul|, \lean|nz_append_single|, \lean|tW_zero_ind|.
Proved by induction on $Bs$ with \lean|tW_append|.
\end{theorem}

\subsection{\texttt{Connective}: binder and splitter}

Helpers: \lean|pair_iff| ($b<t$: $K=at+b\iff K/t=a\land K\bmod t=b$), \lean|tB_sum|
($\sum_{a<s}\sum_{b<t}\mathrm{tB}(s,t,K,a,b)=[K<st]$ for $t>0$), \lean|rt_mul|
($\sqrt{st}=\sqrt s\sqrt t$), \lean|tB_zero| ($\mathrm{tB}(s,t,K,0,0)=[K=0]$).
In each lemma $ds$, $dt$, $dk$ are dimension indices with \lean|E.dimv dk = E.dimv ds * E.dimv dt|,
$a$ has index $ds$, $b$ has index $dt$, and the product label has index $dk$.
\begin{description}[style=nextline,font=\normalfont\ttfamily]
\item[connStates] $[\langle B,[k],[a,b],ds,dt\rangle,\langle Z,[a],[\,],ds,p_1\rangle,\langle
  Z,[b],[\,],dt,p_2\rangle]$ summed $[a,b]$ is $\langle Z,[k],[\,],dk,c\rangle$ (all phaseless);
  scalars $1,1$.
\item[connEffects] $\langle S,[a,b],[k]\rangle$ with Z effects on $a$, $b$ is the Z effect on $k$.
\item[connStatesX, connEffectsX] The same with X spiders; the proof uses \lean|rt_mul| and
  \lean|tB_zero|.
\item[connInverse] $[\langle S,[a,b],[k_1],ds,dt\rangle,\langle B,[k_2],[a,b],ds,dt\rangle]$ summed
  $[a,b]$ is $\langle I,[x],[y],dk,0\rangle$ with $\{x,y\}=\{k_1,k_2\}$: $B\circ S=\mathrm{id}$.
\item[connInverseBind] $[\langle B,[k],[a_1,b_1]\rangle,\langle S,[a_2,b_2],[k]\rangle]$ summed $[k]$
  is $[\langle I,[x_1],[y_1],ds\rangle,\langle I,[x_2],[y_2],dt\rangle]$ with $\{x_1,y_1\}=\{a_1,a_2\}$
  and $\{x_2,y_2\}=\{b_1,b_2\}$: $S\circ B=\mathrm{id}\otimes\mathrm{id}$.
\end{description}

\subsection{\texttt{WZ}: the Z--W bialgebra}

Helpers: \lean|sumOver_zero|; \lean|sumOver_pin| (for distinct labels $ls$ and $f$ mapping $ls$
outside $ls$: summing $[\forall l\in ls,\tau l=\tau(f l)]\,g(\tau)$ sets each $l$ to
$\sigma(f l)$); \lean|tW_mul| ($\mathrm{tW}(os,a)\mathrm{tW}(os,b)=[a=b]\mathrm{tW}(os,a)$);
\lean|tW_prod|; \lean|zip_map_fst|.

\begin{lemma}[\lean|wzBialgebra|]
A Z spider $\langle Z,[w],zi,dz,p_z\rangle$ with $zi\ne[\,]$ inputs, whose output $w$ feeds
$\langle W,wo,[w]\rangle$, summed over $[w]$, equals the complete bipartite network
\[
 (zi\cdot rows).\mathrm{map}(\langle W,r,[z]\rangle)\ \mathbin{+\!\!+}\ (wo\cdot cols).\mathrm{map}(\langle Z,[o],col,dz,q\rangle)
\]
summed over $rows.\mathrm{flatten}$, scalars $1,1$: one W node per input $z$ of the Z spider with
outputs $r$ (its row), one phaseless Z spider per output $o$ of the W node with inputs $col$ (its
column). The certificate \lean|pin : List (ℕ × ℕ)| sends each row label to the W output it reaches:
$rows.\mathrm{flatten}=pin.\mathrm{map\,fst}$ without duplicates, every row maps under
$\mathrm{ren}\,pin$ to $wo$, the pairs $(e,o)$ over the columns are a permutation of $pin$,
$|zi|=|rows|$, $|wo|=|cols|$, and the row labels avoid $zi\mathbin{+\!\!+}wo\mathbin{+\!\!+}[w]$. Both
$p_z$ and $q$ are phaseless and $d>0$.
\end{lemma}

%======================================================================
\section{The complete axiomatisation}\label{sec:lean-core-complete}

The engine's rule library for the complete ZXW axiom sets of \citet{poor2023zxw} (cited in the code
as [P1]) and of the qufinite ZXW calculus of \citet{wang2023qufinitezxw} ([P2], building on
\citet{wang2021qufinite}) is \module{archytaszx.rewrite.rules_axioms}. Three modules state its rules:
\texttt{Completeness} for the rules that are permutations or reuse an earlier lemma, \texttt{Paper}
for the generators in the papers' own conventions, and \texttt{Weights} for everything else.
Docstrings tag each lemma with the paper and the rule's label in it, for example ``[P1] (Aso)'' or
``[P2] (2.8)''.

\subsection{\texttt{Completeness}}

\begin{description}[style=nextline,font=\normalfont\ttfamily]
\item[wPerm] ([P1] (Sym), [P2] (2.10)) For any phase indices $a,b$ and $os\sim os'$:
  $\langle W,os,[i],dz,a\rangle=\langle W,os',[i],dz,b\rangle$; scalars $1,1$. The W node ignores its
  phase index.
\item[zPerm] ([P2] (2.9)) For $os\mathbin{+\!\!+}is\sim os'\mathbin{+\!\!+}is'$:
  $\langle Z,os,is,dz,a\rangle=\langle Z,os',is',dz,a\rangle$: legs move freely between inputs and
  outputs.
\item[wAso] ([P1] (Aso)) Two nestings $[\langle W,o_1,[i]\rangle,\langle W,o_2,[w]\rangle]$ (summed
  $[w]$) and $[\langle W,p_1,[i]\rangle,\langle W,p_2,[w']\rangle]$ (summed $[w']$) with
  $o_1.\mathrm{erase}\,w\mathbin{+\!\!+}o_2\sim o_3\sim p_1.\mathrm{erase}\,w'\mathbin{+\!\!+}p_2$ are
  equal; the proof chains \lean|wFusion| twice.
\item[tX\_kj] For phases $K_j$, $\mathrm{ex}(\varphi(k))=\omega_d^{kj}$:
  $\mathrm{tX}(d,\varphi,os;is)=(\sqrt d)^{-(|os|+|is|)}\,d\,[\sum os+j\equiv\sum is\pmod d]$.
\item[prod\_ite\_val, cancel\_iff] $\prod_{o\in l}[\tau o=v]c=[\forall o,\tau o=v]c^{|l|}$;
  for $w<d$: $(w+j)\bmod d=0\iff w=(d-j\bmod d)\bmod d$.
\item[kCopy] ([P1] (K0)) A pink $K_j$ state $\langle X,[w],[\,],a\rangle$ ($\mathrm{ex}(\varphi_a(q))=
  \omega^{qj}$) into a phaseless $\langle Z,os,[w],b\rangle$ is
  $os.\mathrm{map}(o\mapsto\langle X,[o],[\,],c\rangle)$ with $c$ also of phases $K_j$, times $k$ where
  $k\sqrt d^{\,|os|}=\sqrt d$.
\end{description}

\subsection{\texttt{Paper}: the generators in the papers' conventions}

\texttt{Paper} restates each generator as [P1] or [P2] defines it and proves it equal to the
ArchytasZX tensor it is translated to, for every $d$. Leg values are read below $d$.
\begin{center}\small
\begin{tabularx}{\textwidth}{@{}l>{\raggedright\arraybackslash}X>{\raggedright\arraybackslash}p{4.6cm}@{}}
\toprule
Definition & Paper entry & Equality \\
\midrule
\lean|paperH d j k| & $\frac{1}{\sqrt d}\,\omega^{jk}$ ([P1] Hadamard box) & \lean|paperH_eq|: $=F_d[j,k]$ \\
\lean|paperZ d| $\varphi$ \lean|vs| & $\sum_{j<d}\mathrm{ex}(\varphi(j))\prod_v[v=j]$ & \lean|paperZ_eq|: $=\mathrm{tZ}$ \\
\lean|paperZbox d a vs| & $\sum_{j<d}(j=0\,?\,1:a_j)\prod_v[v=j]$ ([P1] Z box) & \lean|paperZbox_eq|: $=\mathrm{tZw}$ with $c(0)=1$, $c(j)=a_j$ \\
\lean|paperW d o₁ o₂ i| & $\ket{00}\bra0+\sum_{i=1}^{d-1}(\ket{0i}+\ket{i0})\bra i$ & \lean|paperW_eq|: $=\mathrm{tW}(d,[o_1,o_2],i)$ \\
\lean|paperT o i| & $[o=i]+[o=0\land1\le i]$ ([P1] (YT)) & \lean|paperT_eq|: $=\mathrm{tT}$ \\
\lean|paperTi o i| & $[o=i]-[o=0\land1\le i]$ & \lean|paperTi_eq|: $=\mathrm{tTi}$ \\
\lean|paperSplit d e k i j| & $[i<d][j<e][k=ie+j]$ ([P2] splitter) & \lean|paperSplit_eq|: $=\mathrm{tB}$ (by \lean|rfl|) \\
\lean|paperPink d j os is| & $[\sum os+j\equiv\sum is\pmod d]$ ([P1] (HZ)) & \lean|pink_eq| \\
\lean|paperD d o i| & $[(o+i)\bmod d=0]$ ([P1] (Du) dualiser) & \lean|paperD_eq| \\
\bottomrule
\end{tabularx}
\end{center}
\lean|pink_eq| ($d>0$): $\sqrt d^{\,|os|+|is|}\cdot\mathrm{tX}(d,k\mapsto jk/d,os;is)=d\cdot
\mathrm{paperPink}(d,j,os,is)$, so the pink $K_j$ spider is $d^{(m+n)/2-1}$ times the X spider with
phases $jk/d$ turns. \lean|paperD_eq| ($d>0$): $\sum_{k<d}F[o,k]F[k,i]=\mathrm{paperD}(d,o,i)$.
Helpers: \lean|prod_indicator|, \lean|prod_fou_out|, \lean|prod_fou_in| (the products of $F$ entries
divided, not multiplied, by powers of $\sqrt d$). The docstring of \texttt{Paper} refers the scalars
of this translation to section 2 of the engine's completeness document.

\subsection{\texttt{Weights}: rules through coefficient vectors}\label{sec:lean-core-weights}

\texttt{Weights} (9{,}479 lines, 213 theorems, 64 definitions) states every remaining rule of the
complete axiom sets. Its lemmas are stated for Z and weighted Z spiders alike through their
coefficient vectors, and for X and weighted X spiders through theirs:
\[
\mathrm{zc}(E,n,k)=\begin{cases}1\cdot\mathrm{ex}(\varphi_{n.ph}(k))&n.\mathrm{kind}=\mathtt Z\\
c_{n.ph}(k)&n.\mathrm{kind}=\mathtt{ZW}\\0&\text{otherwise,}\end{cases}\qquad
\mathrm{xc}(E,n,k)=\begin{cases}\mathrm{ex}(\varphi_{n.ph}(k))&\mathtt X\\c_{n.ph}(k)&\mathtt{XW}\\0&\text{otherwise,}\end{cases}
\]
with predicates \lean|isZc n| (\lean|n.kind == .Z || n.kind == .ZW|) and \lean|isXc n|. The bridge
lemmas \lean|val_zc| ($\mathrm{val}=\mathrm{tZw}(d,\mathrm{zc}(E,n),\ldots)$ on
$(outs\mathbin{+\!\!+}ins).\mathrm{map}\,\sigma$) and \lean|val_xc| let one lemma cover both kinds; a
rule hypothesis is then an equation between coefficient vectors at every $k<d$, for example
$\mathrm{zc}(m,k)=\mathrm{zc}(n_a,k)\,\mathrm{zc}(n_b,k)$.

\paragraph{Organization.} The file has no sections; it is a sequence of groups, each group
consisting of value-level helper lemmas, then the definitions that name the rule's node lists and
summed labels, then one or more \lean|LocalEq| rule lemmas. Its construction scheme is uniform:
\begin{enumerate}
\item A rule's sides are given as node lists built by small definitions over records of labels
  (structures \lean|HdC|, \lean|ZhdC|, \lean|WwG|, \lean|WwC|, \lean|QwC|, \lean|MixL|,
  \lean|DzC|, \lean|TaC|; list builders \lean|chainNodes|, \lean|chain2Nodes|, \lean|multNodes|,
  \lean|vNodes|, \lean|vaNodes|, \lean|kzNodes|, \lean|hdNodes|, \lean|zhdNodes|, \lean|wwLeft|,
  \lean|wwRight|, \lean|qwRight|, \lean|MixL.nodes|, \lean|mixNodes|, \lean|mix3Nodes|,
  \lean|dzNodes|, \lean|taNodes|, \lean|qddP|) and summed-label builders (\lean|chainLabs|,
  \lean|chain2Labs|, \lean|vLabs|, \lean|hdInt|, \lean|zhdInt|, \lean|mixLabs|, \lean|mix3Labs|,
  \lean|dzInLabs|, \lean|dzOutLabs|, \lean|dzLabs|, \lean|taLabs|), so one lemma covers every
  arity: a list of chains, multipliers or legs of any length.
\item A \emph{value lemma} evaluates each gadget's internal sum in closed form (for example
  \lean|multiplier_val|, \lean|vBoxVal|, \lean|kzChainVal|, \lean|hdChainVal|, \lean|dz_chain|,
  \lean|taLegVal|, \lean|mixSum|), and \emph{sum lemmas} split a sum over the disjoint label groups
  of independent gadgets into a product (\lean|sumOver_groups|, \lean|sumOver_prodList|,
  \lean|sumOver_mul_split|, \lean|sumOver_split_mul|) or pin labels to values
  (\lean|sumOver_pinK|, \lean|sumOver_pinV|, \lean|sumOver_pin2|, \lean|sumOver_pinAll|,
  \lean|sumOver_allEq|, \lean|sumOver_W|, \lean|sumOver_cols_zero|, \lean|sumOver_cols_hot|).
\item The rule lemma compares the two closed forms; character sums over $F$ entries close by
  \lean|fou_unitary|, \lean|char_sum| and the triangle-number identity \lean|om_tri|
  ($\omega_d^{d(d-1)/2}=(-1)^{d+1}$).
\end{enumerate}
Four combinators on local equations serve the exporter:
\lean|LocalEq.trans'| (chaining two equations with scalars $1,1$), \lean|LocalEq.symm'| (swapping
sides with their scalars), \lean|LocalEq.flip'| (from $L=s\cdot R$ with $s\ne0$ and $t=s^{-1}$,
$R=t\cdot L$), and \lean|LocalEq.idR|, \lean|LocalEq.idL|: from an equation with scalars $1,1$, an
interface label $x$ not summed, a label $a\ne x$ summed on neither side and absent from the right
side's labels, both of index $dz$, and a phaseless $ph$, the right side $R$ with $x$ renamed to $a$
followed by the identity spider $\langle Z,[p],[q],dz,ph\rangle$ ($\{p,q\}=\{x,a\}$), summed over
$iR\mathbin{+\!\!+}[a]$, is still equal to the left side (\lean|idL| does the same on the left).

\paragraph{Gadgets.} The recurring sub-networks, with $d=\mathrm{dimv}\,dz$:
\begin{itemize}
\item \emph{Multiplier} $c=(a,ws,b)$: \lean|multNodes dz pz px c = [⟨Z, ws, [a], dz, pz⟩, ⟨X, [b], ws, dz, px⟩]|,
  a plain green dot from $a$ to the wires $ws$ and a phaseless pink dot from $ws$ to $b$.
  \lean|multiplier_val| and \lean|multVal|: summed over $ws$ it is
  $(\sqrt d)^{-(1+|ws|)}\,d\,[b\equiv|ws|\cdot a\pmod d]$, the map $\ket a\mapsto\ket{ea}$ with
  $e=|ws|$. \lean|multiplierUp_val| and \lean|multT_val| give the transposed multiplier.
\item \emph{Fourier chains}: \lean|chainNodes dz c| is $F;F;F$ from $c.1$ through $c.2.1.1$,
  $c.2.1.2$ to $c.2.2$, and \lean|chainVal| its entry; \lean|fou_cube|: $F^3=\overline F$.
  \lean|chain2Nodes| is $F;F$ (the antipode) with \lean|chain2Val|.
\item \emph{V box} with front label $u$ and multipliers $ms$: \lean|vNodes dz u o pq pz px m ms| is
  the Z box $m$ (from $i$ to $u$), the W node $u\to ms.\mathrm{map\,fst}$, the pink dot
  $ms.\mathrm{map}(c\mapsto c.2.2)\to o$, and the multipliers; \lean|vLabs u ms| lists its summed
  labels; \lean|vConst d ms| $=\prod_{c}(\sqrt d)^{-(1+|c.2.1|)}d\cdot(\sqrt d)^{-|ms|}\sqrt d$;
  \lean|vBoxVal|: its entry is $\mathrm{zc}(m,\sigma i)\cdot\mathrm{vConst}\cdot
  \bigl([\sigma i=0][\sigma o=0]+[\sigma i\ne0]\sum_c[\sigma o\equiv|c.2.1|\,\sigma i]\bigr)$;
  \lean|vF| is that entry over its constant.
\item \emph{Mixed-dimensional Z box} ([P2] 2.3.3): each leg is a record
  \lean|MixL| with fields \lean|a| (pink label), \lean|K| (product label), \lean|v| (leg label),
  \lean|dr|, \lean|p| (pink dot's dimension and phase index), \lean|dt| (leg dimension index),
  \lean|out| (output leg through a splitter, else input leg through a binder);
  \lean|MixL.nodes l| is \lean|[⟨X, [], [a], dr, p⟩, ⟨S, [a, v], [K], dr, dt⟩]| for an output and
  \lean|[⟨X, [a], [], dr, p⟩, ⟨B, [K], [a, v], dr, dt⟩]| for an input. \lean|mixZ|, \lean|mixZc|,
  \lean|mix_ok|: the legs and a (weighted) Z box on the product dimension, summed over every
  $[l.a,l.K]$, equal $\prod_l\sqrt{\mathrm{dimv}\,l.dr}\cdot\mathrm{tZw}(\mathrm{zc}(mb),
  (\sigma\,l.v)_l)$; \lean|MixOK| and \lean|Mix3OK| collect the side conditions (distinct labels,
  positive phaseless pink dots, product dimension $=\prod$ leg dimensions for three legs, leg
  dimension $\le$ box dimension).
\end{itemize}

\paragraph{Rule lemmas.} \Cref{tab:lean-core-weights} lists every \lean|LocalEq| lemma of
\texttt{Weights} with the rule it states. Each conclusion is \lean|LocalEq E L iL R iR iface 1 s|
(or $s,1$ where noted) with the scalar $s$ given by a hypothesis.

{\small
\begin{xltabular}{\textwidth}{@{}>{\raggedright\arraybackslash}p{2.6cm}>{\raggedright\arraybackslash}X@{}}
\caption{The rule lemmas of \texttt{Weights}.}\label{tab:lean-core-weights}\\
\toprule
Lemma & Statement \\
\midrule
\endfirsthead
\toprule
Lemma & Statement \\
\midrule
\endhead
\bottomrule
\endfoot
\lean|zcCongr| & [P1] green spider: a Z/ZW node equals any Z/ZW node with the same dimension, permuted legs and equal coefficients below $d$ (engine \code|phase_to_weight|, \code|weight_to_phase|). \\
\lean|zcFusion| & [P1] (S1): two Z/ZW spiders joined by $w$ fuse, $\mathrm{zc}(m,k)=\mathrm{zc}(n_a,k)\mathrm{zc}(n_b,k)$. \\
\lean|zcFusionN| & (S1) along a list $ws$ of distinct shared wires ($ws\ne[\,]$); the result's legs are the others. \\
\lean|zcScalar| & A leg-free Z/ZW node is the scalar $s=\sum_{k<d}\mathrm{zc}(n,k)$. \\
\lean|zcEpt| & [P1] (Ept): the phaseless pink state into a Z effect of coefficient $1$ at $0$ (\lean|tX_zero_state|) is $\sqrt d$. \\
\lean|zcZer| & [P1] (Zer): the Z state of coefficients $[k=0]$ is the phaseless pink state times $(\sqrt d)^{-1}$. \\
\lean|zcAd| & [P1] (AD): Z effects $a$, $b$ on the two outputs of a W node are the Z effect with $m_0=a_0b_0$, $m_k=a_kb_0+a_0b_k$. \\
\lean|zcK2| & [P1] (K2): pink $K_j$ (1$\to$1) into the spider $a$ is the spider $b$ into pink $K_j$, times $s$, when $a_k=s\,b_{(k+j)\bmod d}$ (\lean|tX_shift|: the pink $K_j$ 1$\to$1 entry is $[a+j\equiv b]$). \\
\lean|zcTi|, \lean|zcT| & [P1] (YT): a W node with a Z effect of coefficients $1,-1,\dots,-1$ (resp.\ $1,\dots,1$) on its first output is $\mathtt{Ti}$ (resp.\ $\mathtt T$). \\
\lean|zcPcy| & [P1] (Pcy): a Z box $n$ ($n_0=1$) on the input of a W node is copies of it on both outputs. \\
\lean|fourierSq|, \lean|fourierSq'| & [P1] (P1): $F;F$ is the phaseless pink state with two outputs, the first bent onto the input by a phaseless Z spider; or the pink state on the two outer legs. \\
\lean|identityXZ| & [P1] (S2): the phaseless pink spider 1$\to$1 is the phaseless green one. \\
\lean|zxwTA| & [P1] (TA): the pink 1$\to$2 under a W node, crossed into two green dots, is a W node in its place, times $(\sqrt d)^{-1}$ (\lean|collapse4|, \lean|ta_arith|). \\
\lean|zxwBsj| & [P1] (Bsj): pink $K_j$ into a W node whose first output ends on a green effect is the Z state of coefficients $[k=0\lor k=(d-j)\bmod d]$, times $\sqrt d$. \\
\lean|kOne| & [P1] (K1): pink $K_j$ (1$\to$1) before a green dot with outputs $os$ is the green dot followed by pink $K_j$ on each output (fresh labels $us$). \\
\lean|xDef|, \lean|xDefT|, \lean|xDefW| & X spider definition: the X spider is the Z spider (same phase) with $F$ on each output and $F;F;F$ on each input (\lean|xDef|); or the Z spider of the negated phase $\varphi_b(k)=\varphi_a((d-k)\bmod d)$ with $F$ into each input and $F;F;F$ out of each output (\lean|xDefT|); \lean|xDefW| for XW/ZW. \lean|kjNeg|: phases $K_{j'}$, $j'+j\equiv0$, are $K_j$ at $-k$. \\
\lean|fourierMult| & [P1] (P1): $F;F$ is the multiplier $e=|ws|$ when $e+1\equiv0$, times $s=\sqrt d^{\,e-1}$. \\
\lean|qfMW| & [P2] (MW): a box $m$ (zero except at $0$ and $d-1$) and the multiplier $e\not\equiv0$ before a W node are the box, the W node and the multiplier $e$ on each output; $s=\sqrt d^{\,e-1}$. \\
\lean|bindDef|, \lean|bindDef0|, \lean|bindDefP| & [P2] (2.8): the binder $\langle B,[c],[a,b]\rangle$ is the splitter with its three legs bent by phaseless Z spiders; read the other way $\langle S,[a,b],[c]\rangle$ (\lean|bindDef0|, no conditions); or with only the legs whose flag is set bent (\lean|bindDefP|, \lean|bendLeg f l x = if f then x else l|). \\
\lean|zxwD1|, \lean|qfD1| & [P1]/[P2] (D1): a spider $m$ followed by $F;F$ is $F;F$ followed by the spider $m'$ of reversed coefficients $\mathrm{zc}(m',k)=\mathrm{zc}(m,(d-k)\bmod d)$; \lean|qfD1| for $F;F$ on every input versus every output (\lean|spiderF2|, \lean|sumOver_chains2|). \\
\lean|zxwVW|, \lean|zxwVW'| & [P1] (VW): a V box into a W node is the W node with a V box on each output (scalar on the right, resp.\ left). \\
\lean|zxwZV| & [P1] (ZV): a Z state $a$ into a V box is the Z state of coefficients $1$ at $0$ and $a_{d-1}$ elsewhere. \\
\lean|zxwVA| & [P1] (VA): boxes $b_c$ on a W node's outputs through multipliers into a pink dot are a V box followed by the Z box $A$ (\lean|vaVal|, \lean|vaClosed|). \\
\lean|qfVA| & [P2] (VA): the same left side is the box $b$, pink $K_{-1}$, a W transpose on a plain green state, pink $K_1$ and the box $A$ (\lean|qfVaR|). \\
\lean|zxwKZ| & [P1] (KZ): a pink $K_1$ state on a W node whose other outputs meet a pink dot through transposed multipliers $e_c$ is a pink $K_1$ state and pink dot states, times $s$ (\lean|kzVal|, \lean|wKz|, \lean|kz_const|). \\
\lean|qfHD| & [P2] (HD): the Hadamard decomposition through chains \lean|HdC| (top value $t$, multiplier wires $ws$, green dot $m,w\to y$, $K_{-1}$, W transpose on a green state, $K_c$) is $s\cdot\langle F,[o],[i]\rangle$ (\lean|hdVal|, \lean|hdClosed|, \lean|chainsVal|, \lean|chainsClosed|). \\
\lean|zxwHD| & [P1] (HD): the same through M boxes (chains \lean|ZhdC|) is $s\cdot\langle F,[o],[i]\rangle$. \\
\lean|hopfW|, \lean|hopfSwappedW|, \lean|hopfXW|, \lean|hopfSwappedXW|, \lean|hopfWXW|, \lean|hopfSwappedWXW| & (Hopf) with a weighted Z spider, a weighted X spider, or both (helpers \lean|tZw_sum_cons|, \lean|tZw_cons_cons|, \lean|tXw_perm|, \lean|tXw_antipode_in|, \lean|tXw_antipode_out|). \\
\lean|xcFusion| & Fusion of two X/XW spiders along $w$, coefficients multiplying. \\
\lean|qfE| & [P2] (E): the pink state of dimension $n$ through the multiplier $e$ and the plain green dot $n\to t$ (binder, weighted Z box, splitter) is $s$ times the pink state of dimension $t$, $s\sqrt n^{\,e}=\sqrt n^{\,3}$. \\
\lean|qfK0| & [P2] (K0): pink $K_j$ of dimension $m$ into the plain green dot with legs $m,n_1,n_2$ is $s$ times pink $K$ states of dimensions $n_1$, $n_2$. \\
\lean|zxwWW| & [P1] (WW): two W nodes on a summed $z$ (outputs $xs$ and the $w$ legs of green dots $p,w\to y$), with a W node $u\to ps$, are the crossed W nodes $x_k\to\mathrm{row}_k$, column W nodes $q_c\to\mathrm{col}_c$ and green dots $p,q\to y$ (\lean|wwLeft|, \lean|wwRight|, \lean|wwRows|, \lean|wwB_zero|, \lean|wwB_hot|). \\
\lean|qfWW| & [P2] (WW): two W nodes on $z$ with outputs $xs$ and $ys$ are the crossed W nodes and column W nodes projected through green dots $q_c\to y_c,s_c$ and a W node $e\to ss$ on a green effect (\lean|qwRight|). \\
\lean|qfDD| & [P2] (DD): the splitter $k\to(a,b)$ of $mn\to(m,n)$ is the pink dot $k\to(x,y)$ of dimension $mn$, $x$ through the multiplier $n$ into the mixed Z box $z\to a$, and $y$ into the mixed Z box $y\to b$ (\lean|mix2|, \lean|qddP|). \\
\lean|qfBZW| & [P2] (BZW): W nodes of dimensions $m$ and $n$ into mixed Z boxes $x_c,y_c\to o_c$ are the mixed Z box $u,v\to w$ into the W node of dimension $k$ (\lean|mix3|, \lean|bzwVal|, \lean|sum4_pin|). \\
\lean|qfDZX| & [P2] (DZX): a pink dot of dimension $mn$ whose legs pass through multipliers $n$ and two-legged mixed Z boxes to dimension $m$ is the pink dot of dimension $m$ (\lean|dzInVal|, \lean|dzOutVal|, \lean|dz_mod|). \\
\lean|qfTA| & [P2] (TA): a W node of dimension $m$ and a pink dot of dimension $n$ whose legs pass through multipliers $e_c$ into mixed Z boxes are, up to $s$, the same with a W node of dimension $n$ for the pink dot (\lean|taNet|, \lean|ta_id|). \\
\lean|qfS1|, \lean|qfS1c| & [P2] (S1): mixed-dimensional Z boxes joined along legs $W$ fuse into one box on the remaining legs, entries multiplying (\lean|zz_join|); with every leg joined, the scalar $\sum_ja_jb_j$ times the pink dots' square roots. \\
\end{xltabular}
}

\paragraph{Other value lemmas of \texttt{Weights}.} \lean|tZw_at|, \lean|tZw_perm| (the weighted
analogues of \lean|tZ_at|, \lean|tZ_perm|); \lean|tW_pair| (a two-output W node at $I<d$ is
$[j_1=I\land j_2=0]+[j_1=0\land j_2=I\land I\ne0]$); \lean|dsum_pick|; \lean|tX_shift|;
\lean|tX_zero_state|, \lean|tX_zero_effect| (phaseless pink state, effect at $j$: $\sqrt d$ if
$j=0$); \lean|tX_kj_state|; \lean|tB_zero_left|; \lean|tZ_pair|, \lean|tZ_pair'|, \lean|tZ_three|;
\lean|val_zc_two|; \lean|fou_neg| ($F[(d-k)\bmod d,x]=\overline{F[k,x]}$); \lean|neg_neg_mod|;
\lean|mult_neg_iff|; \lean|shift_iff|; \lean|sum_two_points| ($\sum_{j<d}f(j)=f(0)+f(d-1)$ for
$d>1$ when $f$ vanishes at every other $j<d$); \lean|sum_pin_shift|; \lean|wGreen|, \lean|wGreen'|
($\sum_g\mathrm{tW}([a,g],b)=[a=0]+[a\ne0][a=b]$); \lean|w_pair_nonzero|; \lean|wsum2|;
\lean|tW_support| (a W node at $u$ is nonzero only on all-zero outputs with $u=0$, or one nonzero
output $u<d$); \lean|tW_cons'|; \lean|wPairSum|, \lean|wPairSum_zero|, \lean|wPairSum_hot|;
\lean|tW_replicateZero|, \lean|tW_single|, \lean|prod_tW_zeros|, \lean|prod_tW_zero_in|,
\lean|tW_allZero|, \lean|tW_two_nz|; \lean|hot_cons|, \lean|hot_split|, \lean|sum_hot_rows|,
\lean|prod_isZero|, \lean|prod_one_hot|; \lean|mixLegIn|, \lean|mixLegOut| (a binder, splitter
with a phaseless pink state, effect summed over $a$: $\sqrt r\,[K=V]$); \lean|wPinkSum|,
\lean|wMultPinkVal| (W node into gadgets into a pink dot, as a character sum); \lean|prod_om_pow|,
\lean|prod_vConst|, \lean|vConst_pow| ($\mathrm{vConst}\cdot\sqrt d^{\sum|c.2.1|}=\sqrt d$),
\lean|vConst_congr|, \lean|vF_congr|; list facts \lean|zip_map_self|, \lean|prod_mapIdx_const|,
\lean|sum_map_range|, \lean|flatMap_single_snd|, \lean|perm_map_flatMap|, \lean|prod_flatMap'|,
\lean|prod_ind|, \lean|prod_ite_all|, \lean|prod_map_mul3|, \lean|sum_ind_count|,
\lean|sum_list_comm|, \lean|filter_single_sum|, \lean|length_filter_fst|, \lean|sum_map_mod|,
\lean|reads_prod|, \lean|prodVal_flatMap'|, \lean|prodVal_mixL|, \lean|sumOver_fsum|,
\lean|sumOver_finsum|, \lean|sumOver_chains|, \lean|sumOver_chains'|, \lean|sumOver_mults|;
\lean|pinAll cs u g σ| ($\sigma$ with each $u(b)$ set to $g(b)$) with \lean|pinAll_update|,
\lean|pinAll_at|, \lean|pinAll_not|; label-coverage lemmas \lean|vNodes_labels|,
\lean|vaNodes_labels|, \lean|kzNodes_labels|, \lean|hdNodes_labels|, \lean|zhdNodes_labels|;
\lean|hdGroup|, \lean|zhdGroup|, \lean|hdAt0|, \lean|hdAtTop|, \lean|zhdAt0|, \lean|zhdAtTop| (one
chain's contribution at $0$ and at $d-1$); \lean|wwLeftVal|, \lean|wwRightVal|, \lean|wwRows_val|,
\lean|wwRows_reads|, \lean|qwRightVal|, \lean|sum_tW_zero|, \lean|sum_tW_hot|.

\paragraph{Side conditions.} Every rule lemma of \texttt{Weights} carries, besides the coefficient
hypotheses, the structural hypotheses of the lemmas of \cref{sec:lean-core-zxrules}: the legs of
each node as explicit lists (\lean|hl : m.outs ++ m.ins = [o]| and the like), the dimension index
of each label (\lean|labIdx E.labs x = dz|), the interface labels in \lean|iface|, the distinctness
of the named labels (\lean|List.Nodup| of a listed tuple), $d>0$, and the phaselessness of every
plain dot. The exporter proves these hypotheses with \lean|rfl|, \lean|decide| or \lean|simp|.

%======================================================================
\section{Rules with bundled ports (\texttt{Unfuse}, \texttt{FuseB}, \texttt{WFuse})}\label{sec:lean-core-bundled}

A \bbox{} family is exported as a network whose ports bundle the copies of a port
(\cref{sec:lean-core-net-digits}); \cref{ch:lean-boxes} specifies that translation (module
\texttt{Bundle} and its successors). The three modules here state rule lemmas over bundled ports.
They use the kind builders of \texttt{Bundle}: \lean|zk cs| is \lean|Z| when $cs=[\,]$ and
\lean|ZB cs| otherwise; \lean|xk co ci| is \lean|X| when both are empty and \lean|XB co ci|
otherwise.

\subsection{\texttt{Unfuse}: port-box unfusion}

Value lemmas:
\begin{itemize}
\item \lean|tZ_split|: $\mathrm{tZ}(d,\varphi,A\mathbin{+\!\!+}B)=\sum_{j<d}\mathrm{tZ}(d,\varphi,A\mathbin{+\!\!+}[j])\cdot
  \mathrm{tZ}(d,0,j{::}B)$, a Z spider split into two along one shared value, the second phaseless.
\item \lean|tZ_phaseless|, \lean|tX_phaseless|: a phaseless spider equals the spider of the zero phase
  function.
\item \lean|bundled_split|: for $|tl|=|cs|$,
  $\mathrm{bundled}(\mu,d,(hd\mathbin{+\!\!+}tl).\mathrm{map}\,\sigma,cs)=hd.\mathrm{map}\,\sigma\mathbin{+\!\!+}
  \mathrm{tailVals}(\mu,d,tl.\mathrm{map}\,\sigma,cs)$; \lean|bundled_one|:
  $\mathrm{bundled}(\mu,d,[v],[c])=\mathrm{digsC}(\mu,d,c,v)$; \lean|tailVals_absent|.
\item \lean|fou_unitary'| (the conjugate-first form of \lean|fou_unitary|), \lean|tX_split_out|,
  \lean|tX_split_in|: an X spider split into two along one output, resp.\ input, the second spider
  phaseless with that one value as its single input, resp.\ output ($d>0$).
\end{itemize}

\begin{lemma}[\lean|unfuseZB|, \lean|unfuseXB|, \lean|unfuseXBi|]
\begin{itemize}
\item \lean|unfuseZB|: with $|tl|=|cs|$, $y$ of index $dz$ and absent from
  $os\mathbin{+\!\!+}hd\mathbin{+\!\!+}tl\mathbin{+\!\!+}[b]$, $o_2\mathbin{+\!\!+}i_2\sim os\mathbin{+\!\!+}hd\mathbin{+\!\!+}[y]$, and
  $c_{ph}$ phaseless:
  \[\begin{aligned}&\langle \mathtt{ZB}(cs\mathbin{+\!\!+}[c]),os,hd\mathbin{+\!\!+}tl\mathbin{+\!\!+}[b],dz,a\rangle\\
  &\quad=\bigl[\langle \mathtt{zk}\,cs,o_2,i_2\mathbin{+\!\!+}tl,dz,a\rangle,\ \langle\mathtt{ZB}[c],[\,],[y,b],dz,c_{ph}\rangle\bigr]
  \text{ summed }[y],\end{aligned}\]
  scalars $1,1$: the last bundled input $b$ moves onto a new phaseless spider joined to the rest by
  the plain wire $y$.
\item \lean|unfuseXB| (last bundled output, $|tl|=|co|$, $o_2\sim hd\mathbin{+\!\!+}[y]$, $d>0$):
  $\langle\mathtt{XB}(co\mathbin{+\!\!+}[c])\,ci,hd\mathbin{+\!\!+}tl\mathbin{+\!\!+}[b],is\rangle=
  [\langle\mathtt{xk}\,co\,ci,o_2\mathbin{+\!\!+}tl,is\rangle,\langle\mathtt{XB}[c][\,],[b],[y],c_{ph}\rangle]$
  summed $[y]$.
\item \lean|unfuseXBi| (last bundled input, $|tl|=|ci|$, $i_2\sim hd\mathbin{+\!\!+}[y]$, $d>0$):
  $\langle\mathtt{XB}\,co\,(ci\mathbin{+\!\!+}[c]),os,hd\mathbin{+\!\!+}tl\mathbin{+\!\!+}[b]\rangle=
  [\langle\mathtt{xk}\,co\,ci,os,i_2\mathbin{+\!\!+}tl\rangle,\langle\mathtt{XB}[\,][c],[y],[b],c_{ph}\rangle]$
  summed $[y]$.
\end{itemize}
These are the lemmas of the engine rule \code|port_box_unfusion|.
\end{lemma}

\subsection{\texttt{FuseB}: fusion with bundled ports}

Helpers: \lean|perm_update_once|; \lean|tZ_fuse|
($\sum_j\mathrm{tZ}(\varphi_a,j{::}P)\,\mathrm{tZ}(\varphi_b,j{::}Q)=\mathrm{tZ}(\varphi_c,P\mathbin{+\!\!+}Q)$ when phases
multiply); \lean|tX_fuse| ($d>0$: $\sum_j\mathrm{tX}(\varphi_a,j{::}P_1;P_2)\,\mathrm{tX}(\varphi_b,Q_1;j{::}Q_2)=
\mathrm{tX}(\varphi_c,P_1\mathbin{+\!\!+}Q_1;P_2\mathbin{+\!\!+}Q_2)$); \lean|labDim_w|; \lean|tailVals_append_map|;
\lean|fusionZ_gen|, \lean|fusionX_core|, \lean|fusionX_gen| (the fusions with extra leg values
$D_1,\dots,D_4$ appended to each spider's lists).

\begin{lemma}[\lean|fusionZG|]
With $|t_1|=|cs_1|$, $|t_2|=|cs_2|$, $w$ occurring once in $o_1\mathbin{+\!\!+}i_1$ and once in
$o_2\mathbin{+\!\!+}i_2$, $w\notin t_1\mathbin{+\!\!+}t_2$, $o_3\mathbin{+\!\!+}i_3\sim(o_1\mathbin{+\!\!+}i_1).\mathrm{erase}\,w\mathbin{+\!\!+}
(o_2\mathbin{+\!\!+}i_2).\mathrm{erase}\,w$, and multiplying phases:
\[[\langle\mathtt{zk}\,cs_1,o_1,i_1\mathbin{+\!\!+}t_1,a\rangle,\langle\mathtt{zk}\,cs_2,o_2,i_2\mathbin{+\!\!+}t_2,b\rangle]
\text{ summed }[w]=\langle\mathtt{zk}(cs_1\mathbin{+\!\!+}cs_2),o_3,i_3\mathbin{+\!\!+}(t_1\mathbin{+\!\!+}t_2),c\rangle.\]
\lean|fusionXG| is the X form: $\mathtt{xk}\,co_1\,ci_1$ with outputs $o_1\mathbin{+\!\!+}to_1$ and inputs
$i_1\mathbin{+\!\!+}ti_1$, $w$ once in $o_1$ and once in $i_2$, fuse into
$\mathtt{xk}(co_1\mathbin{+\!\!+}co_2)(ci_1\mathbin{+\!\!+}ci_2)$ with outputs $o_3\mathbin{+\!\!+}(to_1\mathbin{+\!\!+}to_2)$ and inputs
$i_3\mathbin{+\!\!+}(ti_1\mathbin{+\!\!+}ti_2)$.
\end{lemma}

\paragraph{Copied bare wires.} \lean|digits_eq|: two values below $d^k$ with equal digits $0..k-1$ are
equal. \lean|Kind.tower ms| is \lean|I| for $ms=[\,]$ and \lean|Bun m (Kind.tower ms')| for
$ms=m{::}ms'$; \lean|Kind.tower_cs|: its one output and one input have the chain $ms$.
\lean|tensor_tower|: for values below $d^{\prod\mu(ms)}$ the tower's tensor is
$\mathrm{tI}(d^{\prod\mu(ms)},v_x,v_w)$. \lean|bunI|: with
\lean|E.dimv bz = E.dimv dz ^ (ms.map E.mult).prod|, $d>0$, and $x,w\in\mathit{iface}$ of index $bz$,
$\langle\mathtt{tower}\,ms,[x],[w],dz,0\rangle=\langle I,[x],[w],bz,0\rangle$: copies of a bare wire
along $ms$ are one bare wire on the bundled dimension.

\subsection{\texttt{WFuse}: W fusion over a bundled port}

Helpers: \lean|digs_cons_low| ($\mathrm{digs}(d,M{+}1,x+dy)=x{::}\mathrm{digs}(d,M,y)$ for $x<d$);
\lean|sum_digs| ($\sum_{v<d^M}f(\mathrm{digs}(d,M,v))=\mathrm{sumW}(d,M,f)$: a bundled value is
summed as its digits); \lean|wTerm_map|; \lean|tailVals_reps|, \lean|bundled_reps| (bundling every
value along $[m]$ is concatenating their \lean|digs d (μ m)|); \lean|copyLegs_plain|; \lean|tZw_pair|;
\lean|tW_one|.

\begin{lemma}[\lean|wFuseB|]
Let $b\in\{\mathtt Z,\mathtt{ZW},\mathtt W\}$ and $w:\N\to\C$ with $w(0)=1$, and let every copy of $b$
at input $x<d$ satisfy $\mathrm{tensor}(E,b,d,p_b,vs,[x])=w(x)\,\mathrm{tW}(d,vs,x)$ for
$|vs|=|ys|$. With \lean|E.dimv bz = E.dimv dz ^ ([m].map E.mult).prod|, $t$ of index $bz$, $a,a'$ of
index $dz$, $t\ne a$, $t\notin ys$, $a'\ne a$, $a'\notin ys$, $d>0$ and
\lean|E.coef pz k = w k| for $k<d$:
\[\begin{aligned}&\bigl[\langle\mathtt{WB}\,[[m]],[t],[a],dz,p_w\rangle,\langle\mathtt{Bun}\,m\,b,ys,[t],dz,p_b\rangle\bigr]\text{ summed }[t]\\
&\quad=\bigl[\langle\mathtt{ZW},[a^{\prime}],[a],dz,p_z\rangle,\langle\mathtt{WB}\,(ys.\mathrm{map}\,(\_\mapsto[m])),ys,[a^{\prime}],dz,p^{\prime}_w\rangle\bigr]
\text{ summed }[a^{\prime}],\end{aligned}\]
with interface $a{::}ys$ and scalars $1,1$: a W node whose one bundled output fans into $\mu(m)$
copies of $b$ is a weighted Z spider before one W node on all the copies' outputs.
\emph{Technique:} \lean|sum_digs| turns the bundled sum into \lean|sumW|, \lean|wTerm_map| turns the
copies' product into \lean|wTerm|, and \lean|wAssoc_nil| closes it.
\end{lemma}
\lean|wFuseZ| specializes \lean|wFuseB| to a copy that is a Z or weighted Z spider with one output
$y$, its coefficients those of the weighted Z spider
(\lean|zc E ⟨b, [y], [t], dz, pb⟩ k = zc E ⟨.ZW, [a'], [a], dz, pz⟩ k| and value $1$ at $k=0$); the
exporter invokes it for W fusion in families.

%======================================================================
\section{Mixed networks (\texttt{Mixed})}\label{sec:lean-core-mixed}

A diagram with discards $G$ denotes a completely positive map~\cite{selinger2007cpm}; the library
holds it as its purification, as in the mixed-state completeness of
\citet{carette2019mixed}.

\begin{definition}
\begin{lstlisting}[style=leancore]
structure MNet where
  net : Net          -- the purification
  env : List ℕ       -- the boundary labels its discards end
  deriving DecidableEq, Repr
\end{lstlisting}
\lean|envSum D ls f σ σ'| sums over a value $j<D(l)$ for each $l\in ls$, written into both
assignments:
$\mathrm{envSum}(D,[\,],f,\sigma,\sigma')=f(\sigma,\sigma')$ and
$\mathrm{envSum}(D,l{::}ls,f,\sigma,\sigma')=\sum_{j<D(l)}\mathrm{envSum}(D,ls,f,\sigma[l\mapsto
j],\sigma'[l\mapsto j])$. The CPM entry at plain boundary values $\sigma$ and conjugate values
$\sigma'$ is
\[\begin{aligned}
 M.\mathrm{cpm}(E,\sigma,\sigma')=\mathrm{envSum}\bigl(&\mathrm{labDim}\,E,\ M.\mathrm{env},\\
 &(a,b)\mapsto\mathrm{eval}(E,M.\mathrm{net},a)\,\overline{\mathrm{eval}(E,M.\mathrm{net},b)},\ \sigma,\sigma'\bigr).
\end{aligned}\]
\end{definition}

\begin{theorem}[Doubling facts]
\begin{itemize}
\item \lean|envSum_congr|, \lean|envSum_mul|, \lean|envSum_swap| (congruence under a predicate
  preserved by in-range updates, scaling, swapping the two assignments); \lean|InRange.update|.
\item \lean|cpm_pure| (F3): $\langle N,[\,]\rangle.\mathrm{cpm}(\sigma,\sigma')=\mathrm{eval}(N,\sigma)\,
  \overline{\mathrm{eval}(N,\sigma')}$, by \lean|rfl|.
\item \lean|Equal.cpm| (F4): from \lean|Equal E s₁ N₁ s₂ N₂| and $\sigma,\sigma'$ in range on
  $N_1.\mathrm{bnd}$, for every $env$:
  \[s_1\overline{s_1}\,\langle N_1,env\rangle.\mathrm{cpm}(\sigma,\sigma')=
  s_2\overline{s_2}\,\langle N_2,env\rangle.\mathrm{cpm}(\sigma,\sigma'):\]
  a pure equation lifts to every mixed context, its scalar by its squared modulus.
\item \lean|cpm_phase| (F5): from \lean|Equal E 1 N₁ s N₂| and $s\overline s=1$, the CPMs agree.
\item \lean|cpm_conj| (F1): if $\mathrm{eval}(N_2,a)=\overline{\mathrm{eval}(N_1,a)}$ for all $a$, then
  $\langle N_2,env\rangle.\mathrm{cpm}(\sigma,\sigma')=\langle N_1,env\rangle.\mathrm{cpm}(\sigma',\sigma)$.
\item \lean|cpm_causal| (F6): if the CPM is $[\forall l\in keep,\sigma l=\sigma' l]$ everywhere, it is
  $1$ wherever the plain and conjugate values agree on $keep$.
\end{itemize}
\end{theorem}

The discard rules of \module{archytaszx.rewrite.rules_discard} hold as tensor identities:
\begin{center}\small
\begin{tabular}{@{}lll@{}}
\toprule
Lemma & Identity & Hypotheses \\
\midrule
\lean|discard_fourier| & $\sum_{o<d}F[o,i]\,\overline{F[o,i']}=[i=i']$ & $d>0$, $i,i'<d$ \\
\lean|discard_phase| & $\sum_{o<d}\mathrm{tZ}(\varphi,[o,i])\,\overline{\mathrm{tZ}(\varphi,[o,i'])}=[i=i']$ & $i,i'<d$ \\
\lean|discard_zero| & $\sum_{o<d}[o=0]\,\overline{[o=0]}=1$ & $d>0$ \\
\lean|discard_loop| & $\sum_{a<d}\sum_{b<d}\mathrm{tI}(a,b)\,\overline{\mathrm{tI}(a,b)}=d$ & none \\
\lean|discard_binder| & $\sum_{k<st}\mathrm{tB}(s,t,k,a,b)\,\overline{\mathrm{tB}(s,t,k,a',b')}=[a=a'\land b=b']$ & $a,a'<s$, $b,b'<t$ \\
\bottomrule
\end{tabular}
\end{center}
with the helpers \lean|ex_mul_conj| ($\mathrm{ex}(\theta)\overline{\mathrm{ex}(\theta)}=1$) and
\lean|mixed_tZ_leg| ($\mathrm{tZ}(\varphi,[o,i])=[o=i]\,\mathrm{ex}(\varphi(i))$ for $i<d$).

%======================================================================
\section{Exact arithmetic}\label{sec:lean-core-arith}

The exporters prove scalar identities and refutations by exact computation in cyclotomic fields,
by enclosures of real numbers, and by one real root of an integer polynomial. Four modules supply
these tools.

\subsection{\texttt{Cyclo}: coefficient vectors in \texorpdfstring{$\Q(\zeta)$}{Q(ζ)}}\label{sec:lean-core-cyclo}

An integer list $l=(l_0,\dots,l_{n-1})$ stands for $\sum_il_i z^i$:
$\mathrm{cvz}(z,[\,])=0$, $\mathrm{cvz}(z,a{::}l)=a+z\,\mathrm{cvz}(z,l)$ (Horner's rule). The
operations:
\begin{itemize}
\item \lean|addAt l n c|: $c$ added to entry $n$ (unchanged when $n\ge|l|$); \lean|addAt_length|,
  \lean|cvz_addAt| ($n<|l|$: $\mathrm{cvz}(\mathrm{addAt}(l,n,c))=\mathrm{cvz}(l)+cz^n$),
  \lean|cvz_replicate|.
\item \lean|kadd a b|: entrywise sum, the shorter list padded with zeros; \lean|pscale c a|:
  entrywise $c\cdot$; \lean|pmul a b|: $\mathrm{pmul}([\,],b)=[\,]$,
  $\mathrm{pmul}(a{::}as,b)=\mathrm{kadd}(\mathrm{pscale}(a,b),0{::}\mathrm{pmul}(as,b))$;
  \lean|ktrim|: trailing zeros removed.
\item \lean|pmodStep phi a|: if $|a|<|\phi|$ then $a$; else with $top$ the last entry of $a$,
  $\mathrm{ktrim}(\mathrm{kadd}(a,\ 0^{|a|-|\phi|}\mathbin{+\!\!+}\mathrm{pscale}(-top,\phi)))$, one step reducing
  the top coefficient by the monic $\phi$. \lean|pmod phi n a| applies $\mathrm{pmodStep}\circ\mathrm{ktrim}$
  $n$ times.
\item Value lemmas: \lean|cvz_kadd|, \lean|cvz_pscale|, \lean|cvz_pmul|, \lean|cvz_ktrim|,
  \lean|cvz_append_replicate| ($\mathrm{cvz}(0^k\mathbin{+\!\!+}a)=z^k\mathrm{cvz}(a)$),
  \lean|cvz_pmodStep|, \lean|cvz_pmod| (reduction by a $\phi$ vanishing at $z$ keeps the value),
  \lean|cvz_zeros|, and \lean|cvz_eq_zero_of_pmod|: if $\mathrm{cvz}(z,\phi)=0$ and
  $\mathrm{pmod}(\phi,n,a)$ is all zeros, then $\mathrm{cvz}(z,a)=0$.
\item Polynomials: \lean|pz l| $=\sum_iC(l_i)X^i\in\Q[X]$ with \lean|aeval_pz|, \lean|coeff_pz|,
  \lean|pz_ne| (nonzero when some entry is nonzero), \lean|degree_pz| (degree $<|l|$).
\end{itemize}

\begin{theorem}[\lean|cyc_ne|]
For $M>0$ and an integer list $r$ with $|r|\le\varphi(M)$ (Euler's totient, \lean|M.totient|) and some
nonzero entry, $\mathrm{cvz}(\omega_M,r)\ne0$. \emph{Technique:} $\mathrm{pz}(r)$ is a nonzero rational
polynomial of degree $<\varphi(M)$, and the $M$-th cyclotomic polynomial, irreducible over $\Q$ of
degree $\varphi(M)$, is the minimal polynomial of $\omega_M$ (Mathlib,
\nolinkurl{Mathlib.RingTheory.Polynomial.Cyclotomic.Roots}).
\end{theorem}

Scalar lemmas writing constants as powers of $\omega_M$:
\begin{itemize}
\item \lean|fsum_list| (a \lean|Finset| sum over $\mathrm{range}\,d$ as a list sum),
  \lean|pow_mod_of| ($z^M=1\Rightarrow z^n=z^{n\bmod M}$).
\item \lean|om_sub|: $M=mq\Rightarrow\omega_m=\omega_M^q$; \lean|ex_of|: $\theta=p/M+z$ with $z\in\Z\Rightarrow
  \mathrm{ex}(\theta)=\omega_M^p$; \lean|I_sub|: $M=4q\Rightarrow i=\omega_M^q$;
  \lean|rt_sub2|: $M=8q\Rightarrow\sqrt2=\omega_M^q+\omega_M^{7q}$; \lean|rt_sub3|: $M=12q\Rightarrow
  \sqrt3=\omega_M^q+\omega_M^{11q}$; \lean|rt_sub5|: $M=5q\Rightarrow\sqrt5=1+2\omega_M^q+2\omega_M^{4q}$;
  \lean|rt_split|: $x=k^2s\Rightarrow\sqrt x=k\sqrt s$; \lean|rt_one|; \lean|rt_prod|:
  $x=ab\Rightarrow\sqrt x=\sqrt a\sqrt b$; \lean|conj_om_pow|: $\overline{\omega_d}=\omega_d^{d-1}$.
\item \lean|tWZ d os i : ℤ| is the W entry as an integer, \lean|tWZ_eq| its cast to $\mathrm{tW}$.
\item Length lemmas: \lean|digs_length|, \lean|digsC_length| ($|\mathrm{digsC}(\mu,d,c,v)|=\prod_c\mu$),
  \lean|tailVals_length|, \lean|bundled_length|, \lean|outCs_length|, \lean|inCs_length|,
  \lean|copyOuts_length|, \lean|copyIns_length|; \lean|prod_pow_mul|.
\item Monomials in atoms: \lean|mono xs ks| $=\prod_ix_i^{k_i}$ (integer exponents, zipped);
  \lean|mono_nil|, \lean|mono_kadd| (for nonzero atoms), \lean|mono_ktrim|; \lean|ex_zmul|
  ($\mathrm{ex}(kg)=\mathrm{ex}(g)^k$), \lean|ex_lin|, \lean|ex_atoms|: if
  $\theta=p/M+z+\sum_ik_ig_i$ then $\mathrm{ex}(\theta)=\omega_M^p\cdot\mathrm{mono}((\mathrm{ex}\,g_i)_i,k)$;
  \lean|exp_eq_ex|, \lean|ex_atoms_ne|.
\end{itemize}

\subsection{\texttt{GaussSign}: the sign of the quadratic Gauss sum}

For odd $p$ with $h=(p-1)/2$, the quadratic Gauss sum $G=\sum_{k<p}\left(\frac kp\right)\zeta_p^k$
satisfies $G^2=(-1)^hp$ and the classical product formula
$G=\prod_{j<h}(\zeta^{2j+1}-\zeta^{-(2j+1)})$~\cite{berndt1998gauss,landsberg1893gauss}. The module
derives the sign from these two facts.
\begin{itemize}
\item \lean|gneg p| $=|\{j<(p-1)/2:p<2(2j+1)\}|$.
\item \lean|gfac p a|: the vector of $x^a-x^{p-a}$ of length $p+1$; \lean|gprod p n|:
  $\prod_{j<n}(x^{2j+1}-x^{p-(2j+1)})$ by \lean|pmul|; \lean|lgVec p|: the Legendre symbols
  \lean|QSem.lg p k| for $k<p$ (\lean|QSem.lg p a| is $0$ if $a\equiv0$, $1$ if $a^{\lfloor p/2\rfloor}=1$ in
  $\Zd p$, else $-1$; defined in \texttt{PathSumQudit}).
\item \lean|cvz_map_range|, \lean|cvz_gfac|, \lean|cvz_gprod| ($2n\le p+1$) evaluate these vectors;
  \lean|om_factor|: $\omega_p^a-\omega_p^{p-a}=2i\sin(2\pi a/p)$; \lean|sin_sign|: for $0<a<p$,
  $2a\ne p$, $\sin(2\pi a/p)=(p<2a\,?\,-1:1)\,|\sin(2\pi a/p)|$; \lean|prod_sign|.
\end{itemize}
\begin{theorem}[\lean|gauss_sign|, \lean|gauss_rt|]
\lean|gauss_sign|: for odd $p>0$, if $G=\prod_{j<h}(\omega_p^{2j+1}-\omega_p^{p-(2j+1)})$ and
$G^2=(-1)^hp$, then $G=i^h(-1)^{\mathrm{gneg}(p)}\sqrt p$. \lean|gauss_rt|: for odd $p>1$ and a step
count $N$, if
\[\begin{gathered}\mathrm{pmod}(1^p,N,\mathrm{kadd}(\mathrm{lgVec}\,p,\mathrm{pscale}(-1,\mathrm{gprod}(p,h))))\quad\text{and}\\
\mathrm{pmod}(1^p,N,\mathrm{kadd}(\mathrm{pmul}(\mathrm{lgVec}\,p,\mathrm{lgVec}\,p),[-(-1)^hp]))\end{gathered}\]
are all zeros (reductions by $1+x+\dots+x^{p-1}$, written \lean|List.replicate p 1|, checked by
\lean|decide|), then $\sum_{k<p}\left(\frac kp\right)\omega_p^k=i^h(-1)^{\mathrm{gneg}(p)}\sqrt p$.
\end{theorem}

\subsection{\texttt{Ball}: ball arithmetic in \texorpdfstring{$\C$}{C}}

Enclosures in the style of ball arithmetic~\cite{vanderhoeven2009ball}, with the $\ell^1$ distance
of real and imaginary parts.
\begin{definition}
\lean|Near z x y r| $:\iff|\mathrm{Re}\,z-x|+|\mathrm{Im}\,z-y|\le r$; \leanLLone|NearR a m δ| $:\iff|a-m|\le\delta$.
\end{definition}
\begin{itemize}
\item Complex balls: \lean|near_round| (enlarge to a ball $(x',y',r')$ with
  $|x-x'|+|y-y'|+r\le r'$), \lean|near_const|, \lean|near_I| ($i$ in $(0,1,0)$), \lean|near_real|,
  \lean|near_add| (radii add), \lean|near_neg|, \lean|near_sub|, \lean|near_mul|:
  \[\mathrm{Near}(a,x,y,r),\mathrm{Near}(b,u,v,s)\Rightarrow
  \mathrm{Near}(ab,\ xu-yv,\ xv+yu,\ (|x|+|y|)s+(|u|+|v|)r+rs),\]
  using \lean|abs_re_add_abs_im_mul|; \lean|ne_zero_of_near| ($r<|x|$ or $r<|y|$ gives $z\ne0$);
  \lean|re_pos_of_near| ($r<x$ gives $\mathrm{Re}\,z>0$); \lean|near_ofReal|, \lean|near_rat|,
  \lean|near_int|.
\item Exponentials: \lean|cosT A n| $=\sum_{k<n}(-1)^kA^{2k}/(2k)!$ and
  \lean|sinT A n| $=\sum_{k<n}(-1)^kA^{2k+1}/(2k+1)!$; \lean|partial_eq|, \lean|partial_re_im| (the
  partial exponential series of $Ai$ up to degree $2n$ is $\mathrm{cosT}+i\,\mathrm{sinT}$);
  \lean|near_exp|: for $n>0$, $|a-A|\le\varepsilon$ and $|A|\le1$,
  \[\mathrm{Near}\Bigl(e^{ai},\ \mathrm{cosT}(A,n),\ \mathrm{sinT}(A,n),\
  2\Bigl(\varepsilon+|A|^{2n}\frac{2n+1}{(2n)!\,2n}\Bigr)\Bigr).\]
\item Phase bridges: \lean|ex_eq_exp| ($\mathrm{ex}(g)=e^{(2\pi g)i}$), \lean|ex_pow_eight|
  ($\mathrm{ex}(g)=\mathrm{ex}(g/8)^8$), \lean|om_eq_ex| ($\omega_m^e=\mathrm{ex}(e/m)$), \lean|ex_inv|,
  \lean|ex_zpow|, \lean|ex_npow|, \lean|rt_inv_ofReal|.
\item Real balls: \lean|sqrt_near|, \lean|nearR_sqrt| ($0\le m-\delta$, $(m-\delta)^2\le p\le(m+\delta)^2$
  give $|\sqrt p-m|\le\delta$); \lean|inv_near|, \lean|nearR_inv| ($0<m-\delta$, $1/(m-\delta)\le m'+\delta'$,
  $m'-\delta'\le1/(m+\delta)$ give $|x^{-1}-m'|\le\delta'$); \lean|pi_near|, \lean|nearR_pi|:
  $|\pi-3.14159265358979323846|\le10^{-20}$ (from Mathlib's \nolinkurl{Mathlib.Analysis.Real.Pi.Bounds});
  \lean|nearR_const|, \lean|nearR_round|, \lean|nearR_add|, \lean|nearR_neg|, \lean|nearR_mul|
  (radius $|m|\varepsilon+|n|\delta+\delta\varepsilon$).
\item \lean|sqrt_of_sq|: $g^2=p$ and $\mathrm{Re}\,g>0$ give $g=\sqrt p$.
\end{itemize}

\subsection{\texttt{NumField}: one real root and coefficient congruence}

\begin{itemize}
\item \lean|pev cs x| $=\sum_kcs_kx^k$ on $\R$ (Horner, constant term first), \lean|pdv cs x| its
  derivative ($\mathrm{pdv}([\,])=0$, $\mathrm{pdv}(c{::}cs,x)=\mathrm{pev}(cs,x)+x\,\mathrm{pdv}(cs,x)$),
  \lean|pevC| the same polynomial on $\C$; \lean|hasDerivAt_pev|, \lean|pevC_ofReal|,
  \lean|pevC_root|.
\item \lean|root_unique|: if $lo<hi$, $\mathrm{pev}(lo)<0<\mathrm{pev}(hi)$ and $\mathrm{pdv}>0$ on
  $(lo,hi)$, then exactly one $x\in(lo,hi)$ has $\mathrm{pev}(cs,x)=0$ (\lean|∃!|). The exporter uses it
  to name a real algebraic phase or coefficient $\theta$ by its polynomial and isolating interval.
\item \lean|valCoefZ|, \lean|valCoefX|: two Z/ZW (resp.\ X/XW) nodes with equal legs and dimension and
  equal \lean|zc| (resp.\ \lean|xc|) below $d$ have equal values. \lean|netCoef|: networks with
  equal \lean|int| and \lean|bnd| whose nodes have pairwise equal values (\lean|List.Forall₂|) are
  \lean|Equal E 1 A 1 B|.
\end{itemize}

%======================================================================
\section{Exact evaluation and refutation}\label{sec:lean-core-refute}

\subsection{\texttt{Refute}: formal sums and evaluation at a point}

\begin{definition}[Formal sums]
\lean|Terms| is \lean|List (ℤ × ℕ × List ℤ)|; a term $(c,e,k)$ stands for $c\,\zeta^e\,\mathrm{mono}(xs,k)$, and
\lean|tv z xs t| is the sum of the terms at $\zeta=z$ and atoms $xs$. \lean|tmul a b| multiplies term
by term ($(c c',e+e',\mathrm{kadd}(k,k'))$); \lean|tc c| is the constant $[(c,0,[\,])]$. Lemmas:
\lean|tv_cons|, \lean|tv_append|, \lean|tv_tc|, \lean|tv_flatMap|, \lean|tv_mul| (for nonzero
atoms), \lean|tv_foldr_tmul|.
\end{definition}

\begin{definition}[Groups]
\lean|Grp| is \lean|List (List ℤ × List ℤ)|: per atom monomial $k$ one coefficient vector of
powers of $\zeta$; \lean|gval z xs g| $=\sum_{(k,l)}\mathrm{mono}(xs,k)\,\mathrm{cvz}(z,l)$. \lean|klt| is the
lexicographic order on monomials ($[\,]<$ any nonempty list; heads compared first).
\lean|gadd M k e c g| adds $c\zeta^e\,x^k$: into the group of monomial $k$ at entry $e\bmod M$ if
present, else as a new group $(k,\mathrm{addAt}(0^M,e\bmod M,c))$ inserted before the first group whose
monomial is greater, so groups stay sorted. \lean|gnorm M t| folds \lean|gadd| over $t$ (right to
left, each monomial trimmed by \lean|ktrim|) and drops groups whose vector is all zeros. Lemmas:
\lean|gadd_len|, \lean|gval_gadd| ($\zeta^M=1$), \lean|gval_filter|, \lean|gval_gnorm|
($\mathrm{gval}(\mathrm{gnorm}(t))=\mathrm{tv}(t)$ when $\zeta^M=1$, $M>0$).
\lean|ofCoeffs k l i| turns a vector back into terms (nonzero entries only, exponents from $i$),
\lean|gterms g| the groups into terms, and \lean|compact M t| $=\mathrm{gterms}(\mathrm{gnorm}(M,t))$
is a formal sum with at most one term per monomial and exponent below $M$; \lean|tv_ofCoeffs|,
\lean|tv_gterms|, \lean|tv_compact| ($\mathrm{tv}(\omega_M,\cdot)$ is preserved).
\end{definition}

\begin{definition}[Formal entries of generators]
A phase table $P:\N\to\N\to\N\times\mathrm{List}\,\Z$ gives for phase index $p$ and $k$ the exponent of
$\zeta$ and the atom monomial of $\mathrm{ex}(\varphi_p(k))$. With $M$ a multiple of $d$:
\begin{itemize}
\item \lean|tZT P d vs|: for each $k<d$ with every value $=k$, the term $(1,P(k))$.
\item \lean|tXT M P d os is|: for each $k<d$ the term
  $\bigl(1,\ P(k).1+\tfrac Md\,k\textstyle\sum os+\tfrac Md\,(d-1)\,k\sum is,\ P(k).2\bigr)$, the X entry
  without its $(\sqrt d)^{-1}$ factors ($\overline\omega=\omega^{d-1}$).
\item \lean|tensorT S M P|: the formal entry of every kind as in \cref{def:lean-core-tensor}: $F$ is
  $[(1,\tfrac Md\,oi,[\,])]$; $T$, $Ti$, $W$, $I$ are integer constants; bundled kinds bundle first;
  \lean|Bun m b| folds \lean|tmul| over the copies; \lean|ZP dx dy c| on legs $[x,y]$ is
  \lean|tPT M d c| $=[(1,\ \tfrac Md\,((c\,x\,y)\bmod d),\ [\,])]$ (\lean|tP_cyc|: its value at
  $\omega_M$ is the pair phase for $d\mid M$, $M>0$); \lean|ZW|, \lean|XW| and all other cases give
  $[\,]$.
\item \lean|sqc S k d no ni| counts the $(\sqrt d)^{-1}$ factors: $no+ni$ for \lean|X|, $1$ for
  \lean|F|, the bundled lengths for \lean|XB|, $\mu(m)\cdot\mathrm{sqc}(b)$ for \lean|Bun m b|, else $0$.
\item \lean|Kind.noZW|: no \lean|ZW| or \lean|XW|, also under \lean|Bun|.
\item \lean|isBS n| (binder or splitter), \lean|valBS| (their entries as constants), \lean|valT| (a
  node's formal entry), \lean|sqN| (its factor count, $0$ for $B$ and $S$).
\end{itemize}
A \lean|Shape| is an environment's \lean|dimv|, \lean|labs| and \lean|mult| (\lean|Env.shape|), with
\lean|Shape.dim S l| the label's dimension.
\end{definition}

\begin{theorem}[Correctness of formal entries]
\lean|tZ_cyc|, \lean|tX_cyc|, \lean|tensor_cyc|: if every phase satisfies
$\mathrm{ex}(\varphi_p(j))=\omega_M^{P(p,j).1}\mathrm{mono}(xs,P(p,j).2)$ for $j<d$, $d\mid M$, $M>0$, the atoms
are nonzero, and the kind is \lean|noZW|, then
$\mathrm{tensor}(E,k,d,p,os,is)=(\sqrt d)^{-\mathrm{sqc}}\cdot\mathrm{tv}(\omega_M,xs,\mathrm{tensorT}(\ldots))$.
\lean|nodeOK S M keys n| checks that $n$ is a binder or splitter, or that it is \lean|noZW|, $M\bmod
d=0$, and $(n.ph,j)\in keys$ for every $j<d$; \lean|PhOK E M gs P keys| asserts the phase equation for
every listed key with atoms $\mathrm{ex}(g)$, $g\in gs$ (\lean|PhOK.get| extracts one).
\lean|val_cyc| and \lean|prodVal_cyc| lift the equation to nodes and products (\lean|prodT| folds
\lean|tmul| and \lean|compact|).
\end{theorem}

\begin{definition}[Network entry as groups]
\lean|sumT M D ls f σ| sums formal sums over the labels, \lean|compact|ing after each label;
\lean|sumOver_cyc| relates it to \lean|sumOver|. \lean|evalG S M P N σ| $=\mathrm{gnorm}(M,\mathrm{sumT}(M,S.\mathrm{dim},N.\mathrm{int},\mathrm{prodT}(S,M,P,N.\mathrm{nodes}),\sigma))$;
\lean|sqL S N| lists each node's dimension and factor count and \lean|sqK S N|
$=\prod(\sqrt{d_n})^{-c_n}$.
\end{definition}

\begin{theorem}[\lean|eval_cyc|]
If every node is \lean|nodeOK| and \lean|PhOK| holds, then
$\mathrm{eval}(E,N,\sigma)=\mathrm{sqK}(E.\mathrm{shape},N)\cdot\mathrm{gval}(\omega_M,(\mathrm{ex}\,g)_{g\in gs},
\mathrm{evalG}(E.\mathrm{shape},M,P,N,\sigma))$: a network's entry at a concrete environment, computed
exactly.
\end{theorem}

\begin{theorem}[Refutation]
\lean|inRange_of|: a \lean|Bool| check that every label is below its dimension gives
\lean|InRange|. \lean|refute|: if $\sigma$ is in range on $N_1.\mathrm{bnd}$ and $s_1\mathrm{eval}(N_1,\sigma)\ne
s_2\mathrm{eval}(N_2,\sigma)$, then $\lnot\,$\lean|Equal E s₁ N₁ s₂ N₂|. \lean|refute_bnd|: boundaries that
are not permutations of each other refute equality. \lean|assigns D ls| lists every assignment of
$ls$ below $D$ (other labels $0$); \lean|assigns_cover|: every in-range assignment agrees on $ls$
with one of them.
\end{theorem}
A refutation proof therefore evaluates both sides at one boundary assignment with
\lean|eval_cyc| (or \lean|eval_ve_at| below), reduces the difference of the group values with
\lean|pmod|, and concludes it is nonzero with \lean|cyc_ne| or a \lean|Near| enclosure with
\lean|ne_zero_of_near|.

\subsection{\texttt{Contract}: evaluation by checked contraction}

Summing every internal label of a large network one assignment at a time is exponential in the number
of labels. \texttt{Contract} certifies a variable-elimination order
instead~\cite{dechter1999bucket}: factors are tables over small scopes, and each step sums one
label out of the factors containing it.
\begin{definition}[Factor tables]
\begin{lstlisting}[style=leancore]
inductive FT where
  | leaf : Terms → FT          -- a formal sum
  | nil  : FT                  -- end of a level
  | cons : FT → FT → FT        -- the subtable of the next value, then the rest
  deriving DecidableEq, Repr
structure Fac where
  scope : List ℕ               -- one level per label
  table : FT
  deriving DecidableEq, Repr
\end{lstlisting}
\lean|FT.node cs| builds a level from the subtables $cs$; \lean|FT.child T k| is the $k$-th subtable
(\lean|leaf []| when absent); \lean|FT.get σ ls T| reads the levels at the labels $ls$ in order and
returns the leaf's terms ($[\,]$ when the shape does not match). \lean|Fac.at D F σ| is
\lean|F.table.get σ F.scope|; \lean|fv z xs D F σ| its value; \lean|fprod| the product of factor
values; \lean|ffold| the product of factor entries as one formal sum (\lean|tv_ffold|). Lemmas:
\lean|FT.get_congr|, \lean|fv_reads| (a factor reads its scope), \lean|fprod_reads|,
\lean|fprod_filter|.
\end{definition}

\begin{definition}[Checks]
\begin{itemize}
\item \lean|leafOK S M P n F|: every label of $n$ is in $F$'s scope, and for every assignment of the
  scope (\lean|assigns|) the factor's entry equals \leanLLone|compact M (valT S M P τ n)|.
  \lean|valT_congr|, \lean|leaf_val|.
\item \lean|elimOK M D l Fs Fs'|: with $ins$ the factors of $Fs$ containing $l$ and $outs$ the
  others, $Fs'$ is $outs$ followed by one new factor whose scope excludes $l$, is covered by the
  scopes of $ins$, covers every label of $ins$ other than $l$, and whose entry at every assignment
  of its scope is \lean|compact M| of $\bigcup_{j<D(l)}\mathrm{ffold}(D,ins,\tau[l\mapsto j])$.
  \lean|elim_val|: then $\sum_{j<D(l)}\mathrm{fprod}(Fs,\sigma[l\mapsto j])=\mathrm{fprod}(Fs',\sigma)$ for
  $\sigma$ in range on the scopes of $Fs'$.
\item \lean|chainOK M D Fs steps final|: the steps $(l,Fs')$ each satisfy \lean|elimOK| in turn and
  the last list is $final$. \lean|scopesOK U Fs steps|: every label of every factor lies in $U$, is
  the current step's label, or is eliminated later. \lean|chain_val|:
  $\mathrm{sumOver}(D,(\text{step labels}).\mathrm{reverse},\mathrm{fprod}(Fs),\sigma)=\mathrm{fprod}(final,\sigma)$.
\item \lean|prodVal_leaves|, \lean|scopesOK_leaves| connect node products to leaf factors.
\end{itemize}
\end{definition}

\begin{theorem}[\lean|eval_ve| and variants]
If the nodes are \lean|nodeOK|, \lean|PhOK| holds, one leaf factor per node passes \lean|leafOK|,
$N.\mathrm{int}$ is a permutation of the reversed step labels, the chain passes \lean|chainOK| and
\lean|scopesOK N.bnd|, and $\sigma$ is in range on the boundary, then
$\mathrm{eval}(E,N,\sigma)=\mathrm{sqK}(E.\mathrm{shape},N)\cdot\mathrm{fprod}(\omega_M,(\mathrm{ex}\,g),\ldots,final,\sigma)$.
\begin{itemize}
\item \lean|eval_ve_new|: the steps give only their new factors $(l,new)$: \lean|elimNew| checks one
  step, \lean|elimStep l Fs new| $=Fs$ without the factors containing $l$, then $new$;
  \lean|expand| rebuilds the factor lists; \lean|chainNew| checks the chain; \lean|chainNew_nil|,
  \lean|chainNew_step|, \lean|elimOK_step|, \lean|chainOK_expand|, \lean|expand_fst| relate the two
  forms.
\item \lean|eval_ve_at|: leaves may read the boundary labels from $\sigma$ (\leanLLone|pin U σ τ x| is
  $\sigma(x)$ for $x\in U$, else $\tau(x)$; \lean|leafAt| replaces \lean|leafOK|; \lean|leaf_val_at|,
  \lean|prodVal_leaves_at|), provided no boundary label is a step label.
\item \lean|eval_ve_groups|: $\mathrm{fprod}(final,\sigma)=\mathrm{gval}(\mathrm{gnorm}(M,\mathrm{ffold}(final,\sigma)))$.
\end{itemize}
\end{theorem}

%======================================================================
\subsection{\texttt{RefuteMod}: refutation by residues modulo a prime}\label{sec:lean-core-refutemod}

\file{RefuteMod.lean} (imports \file{Contract}) evaluates a network's formal entry modulo a prime $p$ in which an element $w$ of order $M$ stands for $\zeta=\omega_M$, by a checked contraction of residue tables.

\begin{definition}[Weighted formal entries]\label{def:lean-core-tensorW}
For coefficient tables $C:\N\to\N\to\mathrm{Terms}$ (the coefficient of entry $k$ of phase row $p$ as a formal sum):
\begin{itemize}
\item \lean|tZWT C d vs| $=\mathrm{flatMap}_{k<d}\,[\,C(k)$ if every value of $vs$ is $k$, else $[\,]\,]$;
\item \lean|tshift e t| adds $e$ to every exponent of $t$, and
\[
\mathrm{tXWT}(M,C,d,os,is)=\mathrm{flatMap}_{k<d}\ \mathrm{tshift}\bigl(\tfrac Md\,k\textstyle\sum os+\tfrac Md\,(d-1)\,k\sum is,\ C(k)\bigr);
\]
\item \lean|tensorW S M P C| is \lean|tZWT| for \lean|ZW| (over outputs then inputs), \lean|tXWT| for \lean|XW|, the \lean|tmul| fold over the $\mu(m)$ copies (legs by \lean|copyLegs|) for \lean|Bun m b|, and \lean|tensorT| for every other kind;
\item \lean|sqcW| counts $(\sqrt d)^{-1}$ factors ($no+ni$ for \lean|XW|, $\mu(m)\cdot\mathrm{sqcW}(b)$ for \lean|Bun m b|, else \lean|sqc|); \lean|qcW| counts $(Q\,ph)^{-1}$ factors ($1$ for \lean|ZW| and \lean|XW|, $\mu(m)\cdot\mathrm{qcW}(b)$ for \lean|Bun m b|, else $0$); \lean|Kind.phased| holds for \lean|Z|, \lean|X|, \lean|ZB|, \lean|XB|, \lean|ZW|, \lean|XW|, and bundles of them.
\end{itemize}
\lean|valW|, \lean|sqNW|, \lean|qNW| are the node forms (binders and splitters by \lean|valBS|, with no factors). \lean|nodeOKW S M keys ckeys fs n| holds for a binder or splitter, or when $d_n\mid M$ and either the kind is not phased, or it has no \lean|ZW|/\lean|XW| and its phase row is in $fs$, or every entry $j<d_n$ is listed in $keys$ (without weights) or in $ckeys$ (with weights). \lean|CoefOK E M xs C Q keys| asserts $Q(e.1)\ne 0$ and $Q(e.1)\cdot E.\mathrm{coef}(e.1,e.2)=\mathrm{tv}(\omega_M,xs,C(e.1,e.2))$ for each listed key $e$. \lean|prodW| folds \lean|tmul| over the nodes, \lean|sumLT D ls f σ| sums formal sums over the labels by \lean|flatMap|, and \lean|KW S Q N| $=\prod_n(\sqrt{d_n})^{-\mathrm{sqNW}(n)}\,Q(n.ph)^{-\mathrm{qNW}(n)}$.
\end{definition}

\begin{theorem}[Weighted formal entries]\label{thm:lean-core-eval-W}
\lean|tZw_cyc|, \lean|tXw_cyc|, \lean|tensorW_cyc|, \lean|val_cycW|, \lean|prodVal_W|, and
\begin{lstlisting}[style=lean]
theorem eval_W (E : Env) (M : ℕ) (gs : List ℝ) (P : ℕ → ℕ → ℕ × List ℤ) (C : ℕ → ℕ → Terms)
    (Q : ℕ → ℕ) (keys ckeys : List (ℕ × ℕ)) (fs : List ℕ) (N : Net) (σ : ℕ → ℕ) (hM : 0 < M)
    (hok : N.nodes.all (nodeOKW E.shape M keys ckeys fs) = true) (hph : PhOK E M gs P keys)
    (hco : CoefOK E M (gs.map ex) C Q ckeys)
    (hf : ∀ i ∈ fs, ∀ k, ex (E.phase i k) = om M ^ (P i k).1 * mono (gs.map ex) (P i k).2) :
    eval E N σ = KW E.shape Q N * tv (om M) (gs.map ex) (sumLT E.shape.dim N.int (prodW E.shape M P C N.nodes) σ)
\end{lstlisting}
\lean|ex_tab|: for $M=\mathit{den}\cdot q$, $\mathrm{ex}(n/\mathit{den})=\omega_M^{\,((nq)\bmod M)}$.
\end{theorem}

\begin{definition}[Residues]\label{def:lean-core-tmv}
\lean|tmv p w t| $=\sum_{(c,e,k)\in t}c\,w^e\in\Z/p$ (atoms read as $1$); \lean|tpoly t| $=\sum c\,X^e\in\Z[X]$. \lean|pm p k b e| is $b^e\bmod p$ by $k$ rounds of square-and-multiply (\lean|pm_eq|: exact for $e<2^k$). \lean|tmN p w M t| folds $\bigl((c\bmod p)\cdot\mathrm{pm}(p,M,w,e\bmod M)+\mathit{acc}\bigr)\bmod p$ over the terms; \lean|tmN_cast| ($p\ne0$, $w^M=1$ in $\Z/p$, $0<M\le 2^{64}$) and \lean|tmN_castW| (any $M$) equate its cast with \lean|tmv|.
\end{definition}

\begin{theorem}[Residues of vanishing sums]\label{thm:lean-core-tmv-zero}
\lean|tmv_eq_zero|: for $p$ prime, $0<M$, $w$ a primitive $M$-th root of unity in $\Z/p$, and $\mathrm{tv}(\omega_M,[\,],t)=0$, $\mathrm{tmv}(p,w,t)=0$. The proof passes through \lean|tpoly| (\lean|aeval_tpoly|, \lean|eval₂_tpoly|) and the cyclotomic polynomial $\Phi_M$, which divides every integer polynomial vanishing at $\omega_M$ and vanishes at $w$. \lean|prim_of| ($M<2^{64}$) and \lean|prim_ofB| ($M<2^b$) prove that $w$ has order $M$: with $M=\prod q_i^{a_i}$ for primes $q_i$, $\mathrm{pm}(p,64,w,M)=1$ and $\mathrm{pm}(p,64,w,M/q_i)\ne1$ for each $i$; \lean|prime_mem_of_dvd| and \lean|pow_eq_one_of| support them.
\end{theorem}

\begin{definition}[Residue tables and factors]\label{def:lean-core-RT}
\lean|RT| is a table of residues: \lean|leaf v|, \lean|nil|, or \lean|cons h t| (a level listed head to tail; \lean|RT.node cs| folds \lean|cons| over \lean|cs|). \lean|RT.child T k| is the $k$-th subtable ($\mathrm{leaf}\ v$ is its own child, $\mathrm{nil}$ gives $\mathrm{leaf}\ 0$); \lean|RT.get σ ls T| descends one level per label of $ls$ at $\sigma(l)$ and reads a leaf ($0$ if none). An \lean|RFac| has a \lean|scope : List ℕ| and a \lean|table : RT|; \lean|F.at σ| $=\mathrm{get}(\sigma,F.\mathit{scope},F.\mathit{table})$. \lean|rprod p Fs σ| $=\prod_F F.\mathrm{at}(\sigma)\in\Z/p$ and \lean|rfold| the same product mod $p$ in $\N$ (\lean|rfold_cast|). \lean|RT.pack| collapses a level whose entries are all one leaf (\lean|RT.allLeaf|), \lean|RT.isZ| recognizes $\mathrm{leaf}\ 0$, \lean|RT.rest| drops a level's first entry, \lean|RT.beq| is structural equality (\lean|RT.eq_of_beq|), and \lean|RFac.trim| empties the scope of a one-leaf factor.
\end{definition}

\begin{definition}[Checked leaves and eliminations]\label{def:lean-core-leafR}
With $\mathrm{pin}(U,\sigma,\tau)(x)=\sigma(x)$ for $x\in U$ and $\tau(x)$ otherwise:
\begin{itemize}
\item \lean|leafR S M P C p w U σ n F|: every leg of $n$ is in $F$'s scope or in $U$, and at every assignment $\tau$ of the scope (\lean|assigns|) $F.\mathrm{at}(\tau)=\mathrm{tmN}(p,w,M,\mathrm{valW}(\ldots,\mathrm{pin}(U,\sigma,\tau),n))$;
\item \lean|leafXR … n F g|: $n$ is an X spider with $d\mid M$, $d>0$; every leg is in the scope or in $U$; the one-level table $g$ holds at each $s<d$ the residue of $n$'s formal entry with the single output $s$ (\lean|tXT M (P n.ph) d [s] []|); and $F.\mathrm{at}(\tau)=g\bigl((\sum\mathrm{outs}+(d-1)\sum\mathrm{ins})\bmod d\bigr)$ (\lean|tmN_tXT|: the X entry reads its legs only through that sum);
\item \lean|leafXD … n F G b|: as \lean|leafXR| for an X spider whose legs are distinct, exactly the scope, and outside $U$, with the residues packed $b$ bits each in the natural number $G$ and the table equal to \lean|xdesc| (the table over the scope of $G$'s entry at the weighted leg sum mod $d$; \lean|xdesc_get|);
\item \lean|LeafH| selects \lean|plain|, \lean|x g|, or \lean|xd G b|, and \lean|leafG| runs the matching check.
\end{itemize}
\lean|elimR p D l Fs new| checks summing out label $l$: with $\mathit{ins}$ the factors whose scope contains $l$, $l\notin\mathit{new.scope}$; every label of $\mathit{new}$ is in some factor of $\mathit{ins}$; every label of $\mathit{ins}$ other than $l$ is in $\mathit{new}$; and at every assignment $\tau$ of $\mathit{new}$'s scope $\mathit{new}.\mathrm{at}(\tau)=\bigl(\sum_{j<D(l)}\mathrm{rfold}(p,\mathit{ins},\tau[l\mapsto j])\bigr)\bmod p$. \lean|elimS| is its scope part and \lean|elimV| its entry part (\lean|elimR_split|); \lean|elimC| checks the entries in chunks over a prefix of the scope (\lean|setL|, \lean|combos|, \lean|elimR_chunks|); \lean|elimT| checks them by one descent through all tables at once (\lean|descR|, \lean|botS|, cursors \lean|adv1|, \lean|adv1R|, \lean|advAll|, \lean|rfoldL|) and implies \lean|elimR| (\lean|elimT_ok|). \lean|elimStepR l Fs new| keeps the factors without $l$ and appends $\mathit{new}$ trimmed. \lean|chainR p D Fs steps final| runs the eliminations in order and requires the result to be \lean|final|; \lean|scopesR U Fs steps| requires every label of every factor to lie in $U$ or be eliminated later.
\end{definition}

\begin{theorem}[Contraction of residue factors]\label{thm:lean-core-eval-res}
\begin{lstlisting}[style=lean]
theorem eval_res (S : Shape) (M : ℕ) (P : ℕ → ℕ → ℕ × List ℤ) (C : ℕ → ℕ → Terms) (p w : ℕ)
    (hp : p ≠ 0) (hw : (w : ZMod p) ^ M = 1) (hM0 : 0 < M) (hMb : M ≤ 2 ^ 64) (N : Net)
    (σ : ℕ → ℕ) (leaves : List RFac) (hlen : leaves.length = N.nodes.length)
    (hleaf : (N.nodes.zip leaves).all (fun q => leafR S M P C p w N.bnd σ q.1 q.2) = true)
    (steps : List (ℕ × RFac)) (final : List RFac) (hperm : N.int.Perm (steps.map Prod.fst).reverse)
    (hchain : chainR p S.dim leaves steps final = true) (hscope : scopesR N.bnd leaves steps = true)
    (hdis : N.bnd.all (fun x => !(steps.map Prod.fst).contains x) = true)
    (hσ : ∀ x ∈ N.bnd, σ x < S.dim x) :
    tmv p w (sumLT S.dim N.int (prodW S M P C N.nodes) σ) = (rfold p final σ : ZMod p)
\end{lstlisting}
\lean|eval_resW| drops $M\le 2^{64}$; \lean|eval_resG| and \lean|eval_resGW| take one \lean|LeafH| per node and check it by \lean|leafG|. The proof: \lean|tmv_sumLT| and \lean|tmv_prodW| move the residue map inside the sums and products; \lean|prod_leavesR| (or \lean|prod_leavesG|) replaces each node's residue by its leaf; \lean|chainR_val| sums the labels out step by step by \lean|elimR_val| ($\sum_{j<D(l)}\mathrm{rprod}(Fs,\sigma[l\mapsto j])=\mathrm{rprod}(\mathrm{elimStepR}(l,Fs,\mathit{new}),\sigma)$), reordered by \lean|sumR_perm|.
\end{theorem}

\begin{theorem}[Refutation by residues]\label{thm:lean-core-refute-res}
\begin{lstlisting}[style=lean]
theorem refute_res {E : Env} {s₁ s₂ : ℂ} {N₁ N₂ : Net} (σ : ℕ → ℕ) (hin : InRange E N₁.bnd σ)
    (M p w : ℕ) [Fact p.Prime] (hM : 0 < M) (hMb : M ≤ 2 ^ 64)
    (hw : IsPrimitiveRoot (w : ZMod p) M) (K₁ K₂ : ℂ) (T₁ T₂ : Terms) (r₁ r₂ : ℕ)
    (h₁ : eval E N₁ σ = K₁ * tv (om M) [] T₁ ∧ tmv p w T₁ = r₁)
    (h₂ : eval E N₂ σ = K₂ * tv (om M) [] T₂ ∧ tmv p w T₂ = r₂)
    (he : Equal E s₁ N₁ s₂ N₂) (L : ℂ) (f₁ f₂ : List ℤ)
    (hf₁ : L * (s₁ * K₁) = cvz (om M) f₁) (hf₂ : L * (s₂ * K₂) = cvz (om M) f₂)
    (hne : tmN p w M (vterms f₁) * r₁ % p ≠ tmN p w M (vterms f₂) * r₂ % p) : False
\end{lstlisting}
\end{theorem}
\noindent Here $\mathrm{cvz}(z,[a_0,a_1,\ldots])=a_0+z\,(a_1+z(\ldots))$, \lean|vterms f| is the formal sum of the integer vector $f$ (\lean|tv_vterms|), and \lean|tneg| negates a formal sum (\lean|tv_tneg|, \lean|tmv_tneg|). The equation \lean|he| at $\sigma$, multiplied by $L$, gives $\mathrm{tv}(\omega_M,\mathrm{vterms}(f_1)\cdot T_1-\mathrm{vterms}(f_2)\cdot T_2)=0$; \cref{thm:lean-core-tmv-zero} sends it to $0$ in $\Z/p$, against \lean|hne|.

\subsection{\texttt{RefuteSym}: residues by substitution and geometric sums}\label{sec:lean-core-refutesym}

\file{RefuteSym.lean} (imports \file{RefuteMod} and Mathlib's \file{LucasPrimality}) sums labels out symbolically before the residue contraction of \cref{thm:lean-core-eval-res}.

\begin{definition}[Label expressions]\label{def:lean-core-LX}
\lean|LX| has the constructors \lean|v x| (label $x$), \lean|c k|, \lean|add a b|, \lean|mul k a|, \lean|div a k|, \lean|mod a k|, and \lean|neg a n|, with values at $\tau$: $\tau(x)$, $k$, $a+b$, $k\,a$, $\lfloor a/k\rfloor$, $a\bmod k$, and $\mathrm{negN}(n,a)=(n-a\bmod n)\bmod n$ (\lean|LX.eval|). \lean|LX.vars| lists the labels, \lean|LX.subst e l r| substitutes, and \lean|LX.ub D e| is an upper bound when each label $x$ is below $D(x)$ ($D(x)-1$ for a label; sums and products of bounds; $\lfloor\cdot/k\rfloor$; $k-1$ for $\bmod\ k$ with $k>0$; $n-1$ for \lean|neg|; \lean|LX.eval_le_ub|). The smart constructors fold constants: \lean|sadd| (also $0+b=b$, $a+0=a$), \lean|smul|, \lean|sneg|, \lean|smod0|, \lean|smod| (also $(x\bmod k)\bmod k=x\bmod k$ and $(kx+r)\bmod k=r\bmod k$), \lean|sdiv0|, and \lean|sdiv| (also $(x\bmod k)/k=0$ and $(kx+r)/k=x+r/k$ for $k>0$). \lean|LX.norm U σ e| replaces each label of $U$ by $\sigma$'s value and rebuilds bottom up with the smart constructors (\lean|LX.norm_eval|); \lean|LX.nz e| is \lean|e.norm [] (fun _ => 0)|.
\end{definition}

\begin{definition}[Substitution states]\label{def:lean-core-SSt}
\lean|SSt| has bindings \lean|bs : List (ℕ × LX)|, nodes \lean|ns|, divisibility constraints \lean|cs : List (ℕ × LX)|, and a constant \lean|c : ℕ|. \lean|texpr bs x| is the expression bound to $x$ (first binding; else \lean|v x|), \lean|tval bs τ x| its value. A constraint $(N,e)$ has value $[N\mid e(\tau)]$ (\lean|cval|), and
\[
\mathrm{sval}(st,\tau)=c\cdot\prod_{n\in st.\mathit{ns}}\mathrm{tmv}\bigl(p,w,\mathrm{valW}(\ldots,\mathrm{tval}(st.\mathit{bs},\tau),n)\bigr)\cdot\prod_{q\in st.\mathit{cs}}\mathrm{cval}(q,\tau)\in\Z/p.
\]
\lean|sym_start|: $\mathrm{tmv}(\mathrm{sumLT}(N.\mathit{int},\mathrm{prodW}(N.\mathit{nodes}),\sigma))=\mathrm{sumR}(N.\mathit{int},\mathrm{sval}(\langle[\,],N.\mathit{nodes},[\,],1\rangle),\sigma)$.
\end{definition}

\begin{definition}[Delta steps]\label{def:lean-core-dstep}
A source (\lean|Src|) is \lean|zz n x y| (two legs of a Z spider $n$), \lean|bs n pos| (leg \lean|pos| of a binder or splitter), or \lean|cn q| (a constraint). \lean|spair| gives the two expressions the source makes equal when its value is nonzero: the two legs' expressions for a Z spider (each $\bmod\ d^{\mu(m)}$ for a bundle \lean|Bun m Z|, \lean|spairB|); for a binder $B$ with output $k$ and inputs $a,b$ of radix $D=S.\mathrm{dimv}(n.ph)$: at position $0$, $k=Da+b$; at $1$, $a=\lfloor k/D\rfloor$; at $2$, $b=k\bmod D$; for a splitter $S$ with outputs $a,b$ and input $k$: at $2$, $k=Da+b$; at $0$, $a=\lfloor k/D\rfloor$; at $1$, $b=k\bmod D$. \lean|solveP D l| solves $(\mathit{lhs},t)$ for $l$ when $\mathit{lhs}$ is $l$, $\mathrm{neg}(l,D(l))$ (giving $\mathrm{neg}(t,D(l))$), or $l\bmod D(l)$; \lean|csolve| solves a constraint $N\mid l+b$ (giving $\mathrm{neg}(b,N)$) or $N\mid\mathrm{neg}(l,N)+b$ (giving $b\bmod N$) for $N=D(l)$; \lean|solve| dispatches. \lean|dOK S D U R st l src e| holds when $l$ is a summed label of $R$ not yet bound, the source is in the state (\lean|sin|), \lean|solve| gives $e$, $l\notin\mathrm{vars}(e)$, every label of $e$ is in $R$ or $U$, and $\mathrm{ub}(e)<D(l)$. \lean|dapply| binds $l$ to $e$, substitutes $e$ into every binding and constraint, and simplifies by \lean|nz|; \lean|dapplyU U σ| simplifies by \lean|norm U σ| instead.
\end{definition}

\begin{definition}[Geometric and X steps]\label{def:lean-core-gstep}
\lean|fleg bs l n| is, for an $F$ node with one output and one input of which exactly one reads \lean|v l|, the other leg's expression (free of $l$). \lean|gOK S M D R st l n₁ n₂ ns' X Z| holds when $l\in R$ is unbound, $n_1,n_2$ are $F$ nodes with \lean|fleg| giving $X$ and $Z$, both of dimension $D(l)>0$ with $D(l)\mid M$, the state's nodes are a permutation of $n_1::n_2::ns'$, no node of $ns'$ mentions $l$ (\lean|nmen|), and no constraint does. \lean|gapply| removes $n_1,n_2$, adds the constraint $(D(l),X+Z)$, and multiplies the constant by $D(l)$ mod $p$ (by $\sum_{j<N}u^{ja}=N\,[N\mid a]$ for $u$ of order $N$, \lean|geom_root|). \lean|xOK S M P st n ns'| holds for an X spider $n$ of dimension $d>0$, $d\mid M$, all of whose phase exponents are $0$; \lean|xapply| replaces it by the constraint $(d,\ \mathrm{xexpr})$, with $\mathrm{xexpr}$ the sum of its output expressions and $(d-1)$ times its input expressions (the input expressions alone for $d=2$), and multiplies the constant by $d$. \lean|xbOK|, \lean|xbapply| do the same for \lean|Bun m X| with one output and no inputs, with modulus \lean|xbN| $=d^{\mu(m)}$. \lean|SStep| is \lean|d l src e|, \lean|x n|, or \lean|xb n|; \lean|stepS| applies one when its check holds (removing $l$ from $R$ for a delta step), and \lean|runS| a list.
\end{definition}

\begin{theorem}[Steps over sums]\label{thm:lean-core-sym-steps}
\lean|sym_dstep| and \lean|sym_dstepU| (with $U$ disjoint from $R$ and in range): if $R$ is a permutation of $R'\mathbin{+\!\!+}[l]$ and \lean|dOK| holds, then $\mathrm{sumR}(R,\mathrm{sval}(st),\sigma)=\mathrm{sumR}(R',\mathrm{sval}(st'),\sigma)$ for $st'$ = \lean|dapply| (resp.\ \lean|dapplyU|). \lean|sym_gstep| (for $p$ prime and $w$ of order $M$): the same under \lean|gOK| and \lean|gapply|. \lean|sym_xstep| and \lean|sym_xbstep|: $\mathrm{sumR}(R,\mathrm{sval}(st))=\mathrm{sumR}(R,\mathrm{sval}(st'))$ under \lean|xOK|/\lean|xapply| and \lean|xbOK|/\lean|xbapply|. \lean|sym_run| composes them: \lean|runS … R st ss = some (R', st')| gives $\mathrm{sumR}(R,\mathrm{sval}(st),\sigma)=\mathrm{sumR}(R',\mathrm{sval}(st'),\sigma)$.
\end{theorem}

\begin{theorem}[Final contraction]\label{thm:lean-core-sym-final}
\lean|leafS S M P C p w U σ bs g sc F| checks the residue factor of a group of nodes $g$ under bindings $bs$ at every assignment of the labels $sc$ (Z spiders evaluated by \lean|valQ|, which reads a Z spider off its first leg), \lean|leafQ p D U σ q sc F| that of a constraint, and \lean|kleaf p c| $=\langle[\,],\mathrm{leaf}(c\bmod p)\rangle$ is the constant. \lean|sym_final|: when the state's nodes are a permutation of the groups $gss$, each group's leaf and each constraint's leaf check, and \lean|chainR| eliminates the remaining summed labels $R$ from the leaves, constraint leaves, and \lean|kleaf| to \lean|final| with \lean|scopesR U| and $U$ disjoint from the eliminated labels, then $\mathrm{sumR}(R,\mathrm{sval}(st),\sigma)=\mathrm{rfold}(p,\mathit{final},\sigma)$. \lean|leafT| reads the bindings from a binary trie (\lean|BT|, keyed by the bits of the label least significant first; \lean|toBT|, \lean|BT.look|, \lean|leafS_ofT|), and \lean|leafW| checks the table in read order (\lean|walkJ|, \lean|walkT|, \lean|leafS_ofW|).
\end{theorem}

\begin{theorem}[Lucas's primality test]\label{thm:lean-core-prime-lucas}
\begin{lstlisting}[style=lean]
theorem prime_lucas (p g : ℕ) (hp : 1 < p) (hpb : p < 2 ^ 64) (ps : List (ℕ × ℕ))
    (hfac : (ps.map fun x => x.1 ^ x.2).prod = p - 1) (hps : ∀ x ∈ ps, x.1.Prime)
    (h1 : pm p 64 g (p - 1) = 1) (h2 : ps.all (fun x => pm p 64 g ((p - 1) / x.1) != 1) = true) :
    p.Prime
\end{lstlisting}
\lean|prime_lucasB| is the same with $b$ rounds in place of $64$. The proof applies Mathlib's \lean|lucas_primality|~\cite{lucas1878theorie}.
\end{theorem}

\section{Diagrams and their meaning (\texttt{Diagram})}\label{sec:lean-core-diagram}

\texttt{Diagram} represents the engine's diagram as data: nodes with the dimension index of every
port, wires joining two ports, the boundary ports (outputs first), and \bboxes~\cite{carette2022bangboxes}.

\begin{lstlisting}[style=leancore]
structure Port where            -- one port of a node
  node : ℕ                      -- node id
  out  : Bool                   -- output side
  pos  : ℕ                      -- position on that side
  deriving DecidableEq, Repr

structure DNode where
  id : ℕ
  kind : Kind
  dim : ℕ                       -- dimension index
  ph : ℕ                        -- phase index; second dimension index of B, S
  outs : List ℕ                 -- dimension index of each output port
  ins : List ℕ                  -- dimension index of each input port
  deriving DecidableEq, Repr

structure DBox where            -- a !-box
  id : ℕ
  mult : ℕ                      -- multiplicity symbol
  parent : Option ℕ             -- enclosing box
  nodes : List ℕ                -- node scope
  ports : List Port             -- port scope
  deriving DecidableEq, Repr

structure Diagram where
  nodes : List DNode
  wires : List (Port × Port)
  outs : List Port              -- boundary outputs
  ins : List Port               -- boundary inputs
  boxes : List DBox := []
  deriving DecidableEq, Repr
\end{lstlisting}

A \bbox{} with a nonempty node scope is node-scoped; one with an empty node scope and a nonempty port
scope is port-scoped. Definitions in the namespace \lean|Diagram|, for a diagram $D$:
\begin{itemize}
\item \lean|bnd| $=$ \lean|outs ++ ins|: the boundary slots, slot $i$ the $i$-th port.
\item \lean|portDim p|: the dimension index of port $p$, read from the first node with id
  $p.\mathrm{node}$ (\lean|outs.getD pos 0| or \lean|ins.getD pos 0|), $0$ when no node has that id.
  \lean|wireDim j|: the dimension index of wire $j$'s first port ($0$ out of range).
  \lean|bdims|: the boundary dimension indices in slot order.
\item \lean|locate p|: \lean|some (inl i)| for the first boundary slot $i$ holding $p$, otherwise
  \lean|some (inr j)| for the first wire $j$ with $p$ at either end, otherwise \lean|none|.
\item \leanLLone|value σ τ p|: $\sigma(i)$ at a boundary slot, $\tau(j)$ on wire $j$, $0$ otherwise.
\item \lean|DNode.node n f|: the network node of $n$ with output $k$ labeled $f\langle n.id,\mathit{true},k\rangle$
  and input $k$ labeled $f\langle n.id,\mathit{false},k\rangle$, kind, dimension and phase index copied.
\end{itemize}

\begin{definition}[Diagram evaluation]
\[\begin{aligned}
 D.\mathrm{eval}(E,\sigma)=\mathrm{sumOver}\Bigl(&j\mapsto E.\mathrm{dimv}(D.\mathrm{wireDim}\,j),\
 [0,\dots,|wires|-1],\\
 &\tau\mapsto\prod_{n\in D.\mathrm{nodes}}\mathrm{val}\bigl(E,\mathrm{id},n.\mathrm{node}(D.\mathrm{value}\,\sigma\,\tau)\bigr),\
 \underline0\Bigr):
\end{aligned}\]
the sum over a value on every wire of the product of each node's tensor at its port values; a
boundary port reads slot $i$ of $\sigma$. \lean|D.InRange E σ| holds when
$\sigma(i)<E.\mathrm{dimv}(\mathrm{bdims}[i])$ for every slot $i$.
\end{definition}

\begin{definition}[Translation to a network]
\lean|label off p| is the slot $i$ for a boundary port, $off+j$ for a port on wire $j$, and $0$
otherwise. The network of the diagram is
\[\mathtt{toNet}\ off=\bigl\langle D.\mathrm{nodes}.\mathrm{map}(\cdot.\mathrm{node}(\mathrm{label}\,off)),\
[off,\dots,off+|wires|-1],\ [0,\dots,|bnd|-1]\bigr\rangle.\]
\lean|okPort p| holds when $p$ is located on a
slot below $|bnd|$ or a wire below $|wires|$. \lean|check L off| is the conjunction of
\begin{enumerate}
\item $|bnd|\le off$;
\item every port of every node is \lean|okPort|;
\item the two ends of every wire have equal \lean|portDim|;
\item $L.\mathrm{lookup}(off+j)=$ \lean|some (wireDim j)| for every wire $j$;
\item $(L.\mathrm{lookup}\,i)_{i<|bnd|}=$ \lean|bdims.map some|.
\end{enumerate}
\end{definition}

\lean|DNode.mapDim f| and \lean|Diagram.mapDim f| read every dimension index through $f$ (the node's
own, every port's, and for $B$ and $S$ the phase index).

\begin{definition}[Diagram equality]
\lean|DEqual E s₁ D₁ s₂ D₂| $:\iff D_1.\mathrm{bdims}=D_2.\mathrm{bdims}\land\forall\sigma,\ D_1.\mathrm{InRange}(E,\sigma)\to
s_1D_1.\mathrm{eval}(E,\sigma)=s_2D_2.\mathrm{eval}(E,\sigma)$.
\end{definition}

\begin{theorem}
\lean|sumOver_shift| (summing labels $off+j$ is summing $j$ for integrands that agree on such
assignments), \lean|node_rename|, \lean|okPort_label| ($a(\mathrm{label}\,off\,p)=\mathrm{value}\,\sigma\,b\,p$ when
$a(off+j)=b(j)$ and $a(l)=\sigma(l)$ for $l<off$), \lean|eval_toNet|: if
\lean|D.check E.labs off = true| then $\mathrm{eval}(E,D.\mathrm{toNet}\,off,\sigma)=D.\mathrm{eval}(E,\sigma)$;
\lean|labDim_bnd|: slot $i$'s network dimension is \lean|E.dimv (bdims.getD i 0)|.
\lean|DEqual_iff|: if \lean|E.labs = L| and both diagrams pass \lean|check L|, then
\lean|DEqual E s₁ D₁ s₂ D₂ ↔ Equal E s₁ (D₁.toNet o₁) s₂ (D₂.toNet o₂)|.
\end{theorem}
The exporter proves its network equation, then rewrites with \lean|DEqual_iff| (its \lean|check|
hypotheses by \lean|decide|) to state the final theorem over the two diagrams.

%======================================================================
\section{Expansion of \texorpdfstring{\bboxes}{!-boxes} (\texttt{Expand})}\label{sec:lean-core-expand}

\lean|Diagram.expand ks D| instantiates every \bbox{}, the box $b$ at the multiplicity
$k=ks(b.\mathrm{mult})$. It mirrors the engine's expansion (\cref{ch:diagrams}) and is the reference
against which \cref{ch:lean-boxes} proves bundled networks correct.

\paragraph{Auxiliary definitions.}
\begin{itemize}
\item \lean|portKey p| $=(p.\mathrm{node},\ [p.\mathrm{out}],\ p.\mathrm{pos})$ (inputs before outputs);
  \lean|keyLe| its lexicographic order. \lean|fresh D| $=1+\max$ node id ($1$ for no nodes).
\item \lean|splice refs old maps|: with $here$ the ports of $refs$ in $old$: $refs$ when $here$ is
  empty; else, at the index of the first port of $refs$ in $old$, the block
  $\bigl[\langle\mathrm{lookup}_{mp}(r.\mathrm{node}),r.\mathrm{out},r.\mathrm{pos}\rangle: mp\in maps,\ r\in here\bigr]$
  (copy by copy, each node id through the copy's map, unmapped ids kept) between the ports before and
  after that index that are not in $old$.
\item \lean|scopeBnd D scope|: the boundary ports (inputs, then outputs) on nodes of $scope$,
  merge-sorted by \lean|keyLe|.
\item \lean|gather D b|: walks the ancestors of $b$ from its parent upward (at most $|boxes|$ steps);
  for each ancestor with a nonempty node scope, splices \lean|D.ins| and \lean|D.outs| with
  $old=$ \lean|scopeBnd| of its nodes and the identity map: each node-box ancestor's boundary ports
  become one block at the first of them. Port-scoped ancestors are skipped.
\item \lean|shiftPos g k p| $=p.\mathrm{pos}-c+kc$, $c$ the number of positions $i<p.\mathrm{pos}$ on the same
  side with $g\langle p.\mathrm{node},p.\mathrm{out},i\rangle$; \lean|growSide g k x out ds| replaces each
  dimension index at a position $g$ holds by $k$ copies; \lean|growImg g k p| is the $k$ ports
  $\langle p.\mathrm{node},p.\mathrm{out},\mathrm{shiftPos}+j\rangle_{j<k}$ when $g\,p$, else the one port at
  \lean|shiftPos|.
\item For a node box $b$: \lean|cpId D b j x| $=\mathrm{fresh}(D)+j\cdot|b.\mathrm{nodes}|+$ (index of $x$ in
  the sorted scope); \lean|inS b p| ($p$'s node in the scope); \lean|crossP D b p| ($p$ outside, wired
  to a port inside); \lean|nimg D b k p|: the $k$ copies $\langle\mathrm{cpId}\,j\,p.\mathrm{node},p.\mathrm{out},p.\mathrm{pos}\rangle$
  for an inside port, else \lean|growImg (crossP b) k p|; \lean|spliceN D b k refs|: the boundary with
  its inside ports as one block at the first of them, copy by copy, every other port mapped by
  \lean|nimg|.
\end{itemize}

\begin{algorithm}[Box instantiation]\label{alg:lean-core-expand}
\begin{enumerate}
\item \lean|copyNodes D b k|: every node in $b$'s scope becomes $k$ copies with ids \lean|cpId|; every
  other node's sides grow by \lean|growSide (crossP b) k| (a port wired into the box fans out into
  $k$ ports); every wire $(p,q)$ becomes \lean|(nimg p).zip (nimg q)|; outputs and inputs become
  \lean|spliceN|; every other box drops $b$'s nodes from its scope, gains all copies when it held any
  of them, maps its ports through \lean|nimg|, and is removed when both scopes end empty.
\item \lean|expandNode D b k| $=(\text{if }k=0\text{ then gather}(D,b)\text{ else }D).\mathrm{copyNodes}(b,k)$.
\item \lean|growPorts D b k| with $g(p)=[p\in b.\mathrm{ports}]$: every node's sides grow by
  \lean|growSide g k|, wires zip their images under \lean|growImg g k|, outputs and inputs are
  \lean|flatMap|ped through it, every other box maps its ports and is removed when empty.
  \lean|expandPort D b k| applies \lean|gather| first when $k=0$.
\item \lean|expandFuel ks fuel D|: with fuel left, take the leaf boxes (no box has them as parent)
  and pick the one of least id; none: return $D$; else replace $D$ by \lean|expandNode| (empty port
  scope) or \lean|expandPort| at $k=ks(b.\mathrm{mult})$ and continue with one less fuel.
\item \lean|expand ks D| $=$ \lean|expandFuel ks (D.boxes.length + 1) D|.
\end{enumerate}
\end{algorithm}

%======================================================================
\section{Digit sums of boxed diagrams}\label{sec:lean-core-digits}

The modules \texttt{Elim} to \texttt{OrderK} prove that a boxed diagram's bundled network evaluates
like its expansion. Each states the diagram as a constrained sum over port values or over the digits
of bundled ports, and relates those sums across the instantiation steps. The box chains they use
(\lean|walk|, \lean|nodeBoxes|, \lean|portBoxes|, \lean|nodeChain|, \lean|partner|, \lean|tail|,
\lean|toNetB|, \lean|WF|, \lean|step|, \lean|pick|, \leanLLone|φP|, \lean|gN|) are defined in
\cref{ch:lean-boxes}; in brief, \lean|D.nodeChain x| lists the multiplicity symbols of the node boxes
holding $x$, outermost first, and \lean|D.tail q| lists the multiplicities the port $q$ bundles
beyond its node's own copies (the deeper end's extra node boxes for a wired port, the port boxes for
a boundary port).

\subsection{\texttt{Elim}: constrained sums over ports}

\begin{itemize}
\item \lean|Port.code p| $=2\cdot\mathrm{pair}(p.\mathrm{node},\ 2\,p.\mathrm{pos}+[p.\mathrm{out}])$ (Cantor pairing
  \lean|Nat.pair|), an even label; \lean|Port.decode| inverts it; \lean|decode_code|, \lean|code_inj|,
  \lean|code_even|. Slot $i$ is the odd label $2i+1$; \lean|slots σ c| is $\sigma(c/2)$ at odd $c$, $0$
  at even.
\item \lean|DNode.legs n|: its output ports then its input ports; \lean|Diagram.legs| all of them.
\item \lean|pins|: each wire's second port pinned to its first, each boundary port to its slot.
\item \leanLLone|evalG δ nv σ| sums over a value below $\delta(p)$ on every port $p$ (label
  \lean|p.code|) the product $\prod_n nv(n,\nu\circ\mathrm{code})$ where every pin holds, else $0$.
\item \lean|checkC|: the wire ends and boundary ports are a permutation of the legs, the legs are
  distinct, and wire ends have equal dimension index.
\end{itemize}
\lean|sumOver_map_on|, \lean|pins_iff|, \lean|lookup_of_mem|, \lean|value_wire|, \lean|value_bnd| lead
to \lean|evalG_eq|: under the conditions of \lean|checkC| and in-range slots, the constrained sum
equals the sum over wire values of \lean|Diagram.eval|.

\subsection{\texttt{Digits}: digit sums}

\begin{definition}
\lean|cidx ks c| lists the copy indices along the chain $c$ in lexicographic order, outermost
coordinate first: $\mathrm{cidx}([\,])=[[\,]]$, $\mathrm{cidx}(m{::}c)=[i{::}t: i<ks(m),\ t\in\mathrm{cidx}(c)]$.
\lean|rcidx ks c| $=(\mathrm{cidx}\,ks\,c.\mathrm{reverse}).\mathrm{map\ reverse}$ (innermost coordinate slowest).
Digit $t$ of port $p$ is the label \lean|dlab p t| $=2\,\mathrm{pair}(p.\mathrm{code},\mathrm{encode}\,t)$, digit $t$ of
slot $i$ is \lean|slab i t| $=2\,\mathrm{pair}(i,\mathrm{encode}\,t)+1$ (\lean|Encodable.encode| of the list),
\lean|dport| recovers the port and \lean|slotsD σ| reads slot digits at odd labels.
\end{definition}
For a diagram $D$: \lean|legChain p| $=$ \lean|nodeChain p.node ++ tail p|; \lean|dlabs ks| lists every
port's digits; \lean|dnode ks n t| is copy $t$ of node $n$, each port read as its digits under $t$,
its tail's digits in \lean|rcidx| order; \leanLLone|pinsD ks ν| says every wire's two ends agree digit by
digit and every boundary port's digits equal its slot's; \leanLLone|prodD E ks ν| is the product over the
nodes and their copies ($t\in\mathrm{cidx}(\mathrm{nodeChain})$) of the copies' values, and
\[\begin{aligned}
 D.\mathrm{evalD}(E,ks,\sigma)=\mathrm{sumOver}\bigl(&l\mapsto E.\mathrm{dimv}(D.\mathrm{portDim}(\mathrm{dport}\,l)),\
 D.\mathrm{dlabs}\,ks,\\
 &\nu\mapsto[\mathrm{pinsD}\,ks\,\nu]\cdot\mathrm{prodD}\,E\,ks\,\nu,\ \mathrm{slotsD}\,\sigma\bigr)
\end{aligned}\]
(classical decidability of \lean|pinsD|), the diagram at the multiplicities $ks$ over digits, slot
$i$'s digit $t$ being $\sigma(i,t)$. Facts: \lean|mem_cidx| (membership iff each coordinate is below
its multiplicity), \lean|cidx_nodup|, \lean|mem_rcidx|, \lean|rcidx_append_single|,
\lean|mem_cidx_append|, \lean|mem_cidx_single|, \lean|dlab_inj|, \lean|dlab_ne_slab|,
\lean|dport_dlab|, \lean|slotsD_slab|, \lean|slotsD_dlab|, \lean|mem_dlabs|. \lean|evalD_relabel|:
two digit sums agree when the digits of one are an injective relabeling $\rho$ of the other's that
keeps dimensions, the slots read alike, and pins and products agree under $\rho$.

\subsection{\texttt{ExpandFacts}, \texttt{Relabel}, \texttt{Gather}, \texttt{Split}}

\begin{itemize}
\item \texttt{ExpandFacts}: \lean|wsum w n| $=\sum_{i<n}w(i)$, \lean|wsum_succ|, \lean|range_blocks|
  (consecutive blocks of widths $w$ fill $[0,\mathrm{wsum}\,w\,n)$), \lean|find?_unique|,
  \lean|find?_none|; \lean|walk_eq| and the predicate \lean|ChainOK D pr start L| (every box of $L$ is
  in $D$ and satisfies $pr$; every box of $D$ satisfying $pr$ whose parent is $start$ or a box of $L$ is
  in $L$; box $i$'s parent is $start$ for $i=0$, else box $i-1$; the starts are distinct), with
  \lean|ChainOK.length_le|, \lean|ChainOK.walk| (\lean|walk| returns $L$), \lean|filter_last|,
  \lean|ChainOK.transport|; \lean|shiftPos_eq|, \lean|growSide_length|, \lean|side_grow|,
  \lean|legs_grow| (a grown node's legs are the \lean|growImg| images of its legs).
\item \texttt{Relabel}: \lean|dtup l| (a digit label's copy index), \leanLLone|relab φ l| (a digit label
  through a map of digits $\varphi:\mathrm{Port}\to\mathrm{List}\,\N\to\mathrm{Port}\times\mathrm{List}\,\N$),
  \lean|Port.dflt| $=\langle0,\mathit{false},0\rangle$, \leanLLone|bndMap D D' φ σ' i t| (slot digit $(i,t)$ of
  $D$ read at $D'$'s slot of the image of its port); \leanLLone|evalD_relabelφ|: if $\varphi$ maps the
  digits of $D$ one to one onto those of $D'$, wires to wires, boundary to boundary, keeping
  dimensions, and node products agree, then $D'.\mathrm{evalD}(ks',\sigma')=D.\mathrm{evalD}(ks,\mathrm{bndMap}\,\varphi\,\sigma')$.
\item \texttt{Gather}: \lean|gather| only reorders the boundary (\lean|gather_go|, \lean|gather_eq|,
  \lean|splice_ident_perm|, \lean|lookup_ident|); congruence lemmas for diagrams with equal nodes,
  wires and boxes (\lean|walk_congr| through \lean|ChainOK_congr|, \lean|prodD_congr|);
  \lean|bnd_perm|, \lean|WF_perm_bnd|, \lean|evalD_perm_bnd| (a permuted boundary reads the slots
  through \lean|bndMap| with the identity map); \lean|portMap D b k σ'| composes the reading through
  \lean|gather| and \lean|growPorts|; \lean|evalD_expandPort|: instantiating a well-formed leaf port box
  keeps the digit sum with the slots read through \lean|portMap|, and keeps \lean|WF|.
\item \texttt{Split}: \leanLLone|φS img g p t| sends digit $t$ of a split port ($g\,p$) to its image number
  $t.\mathrm{last}$ with $t$ minus its last coordinate, and any other port to its single image;
  \lean|SplitOK D D' img g ks m| bundles the hypotheses (distinct legs, images partition the new legs,
  split ports have chain $c\mathbin{+\!\!+}[m]$ and $ks(m)$ distinct images of chain $c$, kept ports one image
  of the same chain, dimensions kept, wires map to zips of images, boundary to boundary), with
  \lean|SplitOK.into|, \lean|.onto|, \lean|.inj|; \lean|evalD_split|: then the digit sums agree with the
  slots read through \leanLLone|bndMap (φS img g)|. Helpers \lean|mem_zip_idx|, \lean|getElem_mem_zip|.
\end{itemize}

\subsection{\texttt{Positional}, \texttt{Tower}, \texttt{Pin}, \texttt{NetDigits}}

\begin{itemize}
\item \texttt{Positional}: \lean|size ks c| $=\prod_{m\in c}ks(m)$; \lean|pos ks c t| is $t$'s place
  in \lean|cidx ks c| ($\mathrm{pos}(m{::}c,i{::}t)=i\cdot\mathrm{size}(c)+\mathrm{pos}(c,t)$, $0$ on empty lists);
  \lean|size_cons|, \lean|size_append|, \lean|cidx_pos| (positions enumerate $[0,\mathrm{size})$),
  \lean|pos_lt|, \lean|pos_inj|, \lean|pos_append|, \lean|cidx_length|, \lean|getElem_cidx_pos|;
  \lean|digit_digit| ($j<k$: $\mathrm{digit}(B,j,\mathrm{digit}(B^k,i,v))=\mathrm{digit}(B,ik+j,v)$);
  \lean|setDigits ls d v σ| sets label $ls[i]$ to digit $i$ of $v$; \lean|digit_zero_add|,
  \lean|digit_succ_add|; \lean|sumOver_digits|: a sum over $n$ distinct labels of dimension $d>0$ is
  one sum over $v<d^n$ with the labels set to the digits of $v$.
\item \texttt{Tower}: \lean|blkC μ d p cs vs| (every value's block $p$ in base $d$ to its chain's
  size); \lean|tower_outCsC|, \lean|tower_inCsC| (the chains of \lean|C.foldr Kind.Bun b| are $b$'s with
  $C$ prepended), \lean|outCs_length_ge|, \lean|inCs_length_ge|, list helpers; \lean|tensor_towerC|:
  a generator copied along the chain $C$ is the product over $t\in\mathrm{cidx}(C)$ of the generator at
  every value's block $\mathrm{pos}(C,t)$, for values below $d$ to their chains' sizes; \lean|digsC_pos|:
  $\mathrm{digsC}(\mu,d,c,w)$ is $[w]$ for $c=[\,]$, else the digits at the places $\mathrm{pos}(c,s)$.
\item \texttt{Pin}: \lean|dtl|, \lean|wdl| (the digits of the wires' first ends), \lean|rdl| (of the
  second ends and boundary ports), \lean|pinD| (a boundary digit to its slot digit, any other to its
  partner's digit), \lean|idxOf?_get|, \lean|mem_rdl|, \lean|mem_wdl|, \lean|pinD_wire|,
  \lean|pinD_bnd|, \lean|dlabs_perm|; \lean|evalD_pin|: the digit sum is the sum over the first ends'
  digits alone, the others read through \lean|pinD|.
\item \texttt{NetDigits}: \lean|sumOver_blocks| (summing labels whose values are numbers of digits is
  summing the digits), \lean|flatMap_range_off|, \lean|idxOf_cidx|; \lean|NodeS| (a node is plain, a
  boxed node is no binder or splitter, a node with tails is Z, X or a W node without input tails,
  boxed or tailed nodes have all ports of the node's dimension, positive dimensions);
  \lean|label_eq|; \lean|slotDig E β i t| (digit $\mathrm{pos}(\mathrm{legChain},t)$ of $\beta(i)$ in base the
  slot's dimension); \lean|eval_toNetB| and \lean|eval_toNetB_evalD|: for a well-formed diagram whose
  nodes satisfy \lean|NodeS| and whose wire labels have dimension $d^{\mathrm{size}(\mathrm{legChain})}$, the
  bundled network \lean|D.toNetB off| at a bundled boundary assignment $\beta$ (in range) equals
  \lean|D.evalD E E.mult (D.slotDig E β)|.
\end{itemize}

\subsection{\texttt{Order}, \texttt{OrderK}: the slot order of the expansion}

\begin{itemize}
\item \texttt{Order}: \lean|SDom D ks| lists every slot with every copy index of its chain;
  \leanLLone|ixMap D D' φ x| reads a slot digit of $D$ as one of $D'$ through $\varphi$; \lean|bndMap_ix|;
  \leanLLone|BndBij D D' ks φ| ($\varphi$ maps boundary digits one to one onto $D'$'s, keeping dimensions,
  both boundaries without duplicates) and \lean|IxBij D D' ks f| (the same for an index map), with
  \lean|BndBij.mem|, \lean|.inj'|, \lean|.onto'|, \lean|.dim'|, \lean|.ix|, \lean|IxBij.id|,
  \lean|IxBij.comp|, \lean|WF.bnd_nodup|, \lean|bij_perm|, \lean|bij_port|, \lean|bij_node|,
  \lean|idxOf_getElem'|; \lean|stepIx D ks b| (one instantiation step at the index level),
  \lean|stepMap_ix|, \lean|stepIx_bij|; \lean|bidxFuel| (the composite over \lean|pick| and \lean|step|),
  \lean|bmapFuel_read|, \lean|bidxFuel_bij|; \lean|bidx D ks x| $=$ \lean|bidxFuel ks (D.boxes.length + 1) D x|,
  \lean|bmap_read|, \lean|bidx_bij|: for a well-formed $D$, \lean|bidx| is a bijection from $D$'s slot
  digits onto the expansion's slots, keeping dimensions.
\item \texttt{OrderK}: the expansion with its boundary in key order. \lean|mpB y o| is the least slot
  on side $o$ whose node lies in box $y$; \lean|slotSig i| is slot $i$'s key data (the least slots of
  its node boxes, outermost first, up to the first equal to $i$; then the remaining copy coordinates:
  the remaining node boxes' outermost first, then the port boxes' innermost first); \lean|keySig|;
  \lean|keyOf sg i t| $=$ the interleaving of $sg.1$ with $t$, then $i$, then $t$ read at the positions
  $sg.2$; \lean|lexLe|; \lean|SL ks| (all slot digits, insertion-sorted by \lean|lexLe| of their
  keys); \lean|expandK D ks| (the expansion with outputs and inputs replaced by the first
  $|outs|$ and the rest of the key-ordered slots); \lean|readK ks σ'' i t| $=\sigma''(\text{index of }(i,t)\text{ in SL})$.
  Lemmas \lean|SDom_nodup|, \lean|SDom_boxfree|, \lean|idxOf_map_inj|, \lean|idxOf?_of_mem|,
  \lean|value_perm|, \lean|eval_perm| (a permuted boundary reads slots through the permutation),
  \lean|range_getD|, \lean|SL_perm|, \lean|bidx_slot|, \lean|L_perm|, \lean|L_inj|, \lean|bdims_expandK|,
  and \lean|eval_expandK_evalD|: for a well-formed $D$ and $\sigma''$ in range,
  $(D.\mathrm{expandK}\,\mu).\mathrm{eval}(E,\sigma'')=D.\mathrm{evalD}(E,\mu,D.\mathrm{readK}\,\mu\,\sigma'')$ with
  $\mu=$ \lean|E.mult|.
\end{itemize}

%======================================================================
\section{Dimensions assumed equal (\texttt{Assume})}\label{sec:lean-core-assume}

A proof may merge dimension indices that hold equal values, for example two symbols the user assumes
equal. \lean|dimMap ps x| $=(ps.\mathrm{lookup}\,x).\mathrm{getD}\,x$; \lean|Node.mapDim f| maps the node's
dimension index (and for $B$, $S$ the phase index); \lean|Net.mapDim f| maps every node;
\lean|mapLabs f labs| maps each label's index; \lean|lookup_mapLabs|; \lean|val_mapDim| (if
$E.\mathrm{dimv}\circ f=E.\mathrm{dimv}$, values are unchanged).

\begin{theorem}[\lean|Equal.of_mapDim|, \lean|Equal.of_mapDim'|]
If every pair $(a,b)\in ps$ has $E.\mathrm{dimv}\,b=E.\mathrm{dimv}\,a$, an equation between the merged networks
$N.\mathrm{mapDim}(\mathrm{dimMap}\,ps)$, $M.\mathrm{mapDim}(\mathrm{dimMap}\,ps)$ in the environment with labels
$\mathrm{mapLabs}(\mathrm{dimMap}\,ps)$ gives \lean|Equal E s N s' M|. The primed form takes the merged
environment and networks as named arguments with equations to them.
\end{theorem}
%======================================================================
\section{Lean statements of the principal declarations}\label{sec:lean-core-listings}

This section reproduces, verbatim from the source, the Lean text of the declarations that define the semantics and of the principal lemmas of each module of this chapter: a definition with its body, a theorem with its statement up to (not including) its proof. Together with the summaries of the preceding sections they determine the library up to proofs.

\subsection{\texttt{Net}}
The definitions that fix the semantics, verbatim.
\begin{lstlisting}[style=leancore]
noncomputable def tensor (E : Env) : Kind → ℕ → ℕ → List ℕ → List ℕ → ℂ
  | .Z, d, ph, os, is => tZ d (E.phase ph) (os ++ is)
  | .X, d, ph, os, is => tX d (E.phase ph) os is
  | .F, d, _, [o], [i] => fou d o i
  | .T, _, _, [o], [i] => tT o i
  | .Ti, _, _, [o], [i] => tTi o i
  | .W, d, _, os, [i] => tW d os i
  | .I, d, _, [a], [b] => tI d a b
  | .ZW, d, ph, os, is => tZw d (E.coef ph) (os ++ is)
  | .XW, d, ph, os, is => tXw d (E.coef ph) os is
  | .ZB cs, d, ph, os, is => tZ d (E.phase ph) (os ++ bundled E.mult d is cs)
  | .XB co ci, d, ph, os, is => tX d (E.phase ph) (bundled E.mult d os co) (bundled E.mult d is ci)
  | .WB cs, d, _, os, [w] => tW d (bundled E.mult d os cs) w
  | .ZP _ _ c, d, _, os, is => tP d c (os ++ is)
  | .Bun m b, d, ph, os, is => ((List.range (E.mult m)).map fun i =>
      tensor E b d ph (copyOuts E b d i os) (copyIns E b d i is)).prod
  | _, _, _, _, _ => 0
\end{lstlisting}
\begin{lstlisting}[style=leancore]
noncomputable def val (E : Env) (σ : ℕ → ℕ) (n : Node) : ℂ :=
  let d := E.dimv n.dim
  let os := n.outs.map σ
  let is := n.ins.map σ
  match n.kind, os, is with
  | .Z, os, is => tZ d (E.phase n.ph) (os ++ is)
  | .X, os, is => tX d (E.phase n.ph) os is
  | .F, [o], [i] => fou d o i
  | .T, [o], [i] => tT o i
  | .Ti, [o], [i] => tTi o i
  | .W, os, [i] => tW d os i
  | .B, [k], [a, b] => tB d (E.dimv n.ph) k a b
  | .S, [a, b], [k] => tB d (E.dimv n.ph) k a b
  | .I, [a], [b] => tI d a b
  | .ZW, os, is => tZw d (E.coef n.ph) (os ++ is)
  | .XW, os, is => tXw d (E.coef n.ph) os is
  | .ZB cs, os, is => tZ d (E.phase n.ph) (os ++ bundled E.mult d is cs)
  | .XB co ci, os, is => tX d (E.phase n.ph) (bundled E.mult d os co) (bundled E.mult d is ci)
  | .WB cs, os, [w] => tW d (bundled E.mult d os cs) w
  | .Bun m b, os, is => tensor E (.Bun m b) d n.ph os is
  | .ZP _ _ c, os, is => tP d c (os ++ is)
  | _, _, _ => 0
\end{lstlisting}
\begin{lstlisting}[style=leancore]
noncomputable def tP (M : ℕ) (c : ℤ) : List ℕ → ℂ
  | [x, y] => ex ((c : ℝ) * x * y / M)
  | _ => 0
\end{lstlisting}
\begin{lstlisting}[style=leancore]
noncomputable def sumOver (D : ℕ → ℕ) : List ℕ → ((ℕ → ℕ) → ℂ) → (ℕ → ℕ) → ℂ
  | [], f, σ => f σ
  | l :: ls, f, σ => ∑ j ∈ range (D l), sumOver D ls f (Function.update σ l j)
\end{lstlisting}
\begin{lstlisting}[style=leancore]
noncomputable def eval (E : Env) (N : Net) (σ : ℕ → ℕ) : ℂ :=
  sumOver (labDim E) N.int (prodVal E N.nodes) σ
\end{lstlisting}
\begin{lstlisting}[style=leancore]
def InRange (E : Env) (ls : List ℕ) (σ : ℕ → ℕ) : Prop := ∀ l ∈ ls, σ l < labDim E l
\end{lstlisting}
\begin{lstlisting}[style=leancore]
theorem frame (E : Env) (C L R : List Node) (iC iL iR bnd iface : List ℕ) (kL kR : ℂ)
    (hloc : LocalEq E L iL R iR iface kL kR)
    (hL : ∀ x ∈ labelsOf C, x ∉ iL) (hR : ∀ x ∈ labelsOf C, x ∉ iR)
    (hface : ∀ x ∈ iface, x ∈ iC ∨ x ∈ bnd) (σ : ℕ → ℕ) (hσ : InRange E bnd σ) :
    kL * eval E ⟨C ++ L, iC ++ iL, bnd⟩ σ = kR * eval E ⟨C ++ R, iC ++ iR, bnd⟩ σ
\end{lstlisting}
\begin{lstlisting}[style=leancore]
theorem rewrite (E : Env) (C L R : List Node) (iC iL iR bnd iface : List ℕ) (s k s' : ℂ)
    (hloc : LocalEq E L iL R iR iface 1 k)
    (hL : ∀ x ∈ labelsOf C, x ∉ iL) (hR : ∀ x ∈ labelsOf C, x ∉ iR)
    (hface : ∀ x ∈ iface, x ∈ iC ∨ x ∈ bnd) (hs : s' = s * k) :
    Equal E s ⟨C ++ L, iC ++ iL, bnd⟩ s' ⟨C ++ R, iC ++ iR, bnd⟩
\end{lstlisting}
\begin{lstlisting}[style=leancore]
theorem Equal.mul_trans {E : Env} {a b c e : ℂ} {X Y Z : Net} (h : Equal E a X b Y)
    (h' : Equal E c Y e Z) : Equal E (a * c) X (b * e) Z
\end{lstlisting}
\begin{lstlisting}[style=leancore]
theorem Equal.rescale {E : Env} {a b s t : ℂ} {X Y : Net} (h : Equal E a X b Y) (ha : a ≠ 0)
    (hst : s * b = t * a) : Equal E s X t Y
\end{lstlisting}

\subsection{\texttt{Iso}}
\begin{lstlisting}[style=leancore]
def nodeEquiv (n m : Node) : Bool :=
  n.kind == m.kind && n.dim == m.dim && n.ph == m.ph &&
    match n.kind with
    | .Z | .ZW => (n.outs ++ n.ins).isPerm (m.outs ++ m.ins)
    | .X | .XW => n.outs.isPerm m.outs && n.ins.isPerm m.ins
    | _ => n.outs == m.outs && n.ins == m.ins
\end{lstlisting}
\begin{lstlisting}[style=leancore]
def isoOK (labs : List (ℕ × ℕ)) (A B : Net) (π : List (ℕ × ℕ)) (p : List ℕ) : Bool :=
  let f := ren π
  let K := π.map Prod.fst
  let A' := A.nodes.map (Node.rename f)
  let P := p.map fun j => A'.getD j Node.none
  decide K.Nodup && (K.map f).isPerm K &&
    (labelsOf A.nodes).all (fun x => decide (x ∈ A.int) || !decide (x ∈ K)) &&
    A.int.all (fun x => labIdx labs (f x) == labIdx labs x) &&
    A.bnd == B.bnd && B.int.isPerm (A.int.map f) && P.isPerm A' &&
    P.length == B.nodes.length && (List.zip P B.nodes).all fun q => nodeEquiv q.1 q.2
\end{lstlisting}
\begin{lstlisting}[style=leancore]
theorem iso (E : Env) (s : ℂ) (A B : Net) (π : List (ℕ × ℕ)) (p : List ℕ)
    (h : isoOK E.labs A B π p = true) : Equal E s A s B
\end{lstlisting}

\subsection{\texttt{Rules}}
\begin{lstlisting}[style=leancore]
theorem wire_elim (E : Env) (s : ℂ) (ns₁ ns₂ : List Node) (int bnd : List ℕ) (a b dz : ℕ)
    (hb : b ∈ int) (hab : a ≠ b) (ha : a ∈ int ∨ a ∈ bnd)
    (hda : labIdx E.labs a = dz) (hdb : labIdx E.labs b = dz) :
    Equal E s ⟨ns₁ ++ ⟨.I, [a], [b], dz, 0⟩ :: ns₂, int, bnd⟩
      s ⟨(ns₁ ++ ns₂).map (Node.rename (ren [(b, a)])), int.erase b, bnd⟩
\end{lstlisting}
\begin{lstlisting}[style=leancore]
theorem fusionZ (E : Env) (a b c dz w : ℕ) (o₁ i₁ o₂ i₂ o₃ i₃ iface : List ℕ)
    (h₁ : w ∈ o₁ ++ i₁) (h₂ : w ∈ o₂ ++ i₂)
    (hm : (o₃ ++ i₃).Perm ((o₁ ++ i₁ ++ (o₂ ++ i₂)).filter (· ≠ w)))
    (hw : labIdx E.labs w = dz)
    (hph : ∀ k < E.dimv dz, ex (E.phase c k) = ex (E.phase a k) * ex (E.phase b k)) :
    LocalEq E [⟨.Z, o₁, i₁, dz, a⟩, ⟨.Z, o₂, i₂, dz, b⟩] [w] [⟨.Z, o₃, i₃, dz, c⟩] [] iface 1 1
\end{lstlisting}
\begin{lstlisting}[style=leancore]
theorem fusionX (E : Env) (a b c dz w : ℕ) (o₁ i₁ o₂ i₂ o₃ i₃ iface : List ℕ)
    (hd : 0 < E.dimv dz) (h₁ : w ∈ o₁) (h₁' : w ∉ o₁.erase w) (h₁i : w ∉ i₁)
    (h₂ : w ∈ i₂) (h₂' : w ∉ i₂.erase w) (h₂o : w ∉ o₂)
    (ho : o₃.Perm (o₁.erase w ++ o₂)) (hi : i₃.Perm (i₁ ++ i₂.erase w))
    (hw : labIdx E.labs w = dz)
    (hph : ∀ k < E.dimv dz, ex (E.phase c k) = ex (E.phase a k) * ex (E.phase b k)) :
    LocalEq E [⟨.X, o₁, i₁, dz, a⟩, ⟨.X, o₂, i₂, dz, b⟩] [w] [⟨.X, o₃, i₃, dz, c⟩] [] iface 1 1
\end{lstlisting}
\begin{lstlisting}[style=leancore]
theorem identityZ (E : Env) (a dz p q w x : ℕ) (iface : List ℕ)
    (hpq : (p = w ∧ q = x) ∨ (p = x ∧ q = w))
    (hph : ∀ k < E.dimv dz, ex (E.phase a k) = 1) :
    LocalEq E [⟨.Z, [p], [q], dz, a⟩] [] [⟨.I, [x], [w], dz, 0⟩] [] iface 1 1
\end{lstlisting}
\begin{lstlisting}[style=leancore]
theorem identityX (E : Env) (a dz p q w x : ℕ) (iface : List ℕ)
    (hd : 0 < E.dimv dz) (hpq : (p = w ∧ q = x) ∨ (p = x ∧ q = w))
    (hface : p ∈ iface ∧ q ∈ iface) (hp : labIdx E.labs p = dz) (hq : labIdx E.labs q = dz)
    (hph : ∀ k < E.dimv dz, ex (E.phase a k) = 1) :
    LocalEq E [⟨.X, [p], [q], dz, a⟩] [] [⟨.I, [x], [w], dz, 0⟩] [] iface 1 1
\end{lstlisting}

\subsection{\texttt{ZXRules}}
\begin{lstlisting}[style=leancore]
theorem cap (E : Env) (a b dz w : ℕ) (iface : List ℕ) (hd : 0 < E.dimv dz)
    (hw : labIdx E.labs w = dz) (ha : ∀ k < E.dimv dz, ex (E.phase a k) = 1)
    (hb : ∀ k < E.dimv dz, ex (E.phase b k) = 1) :
    LocalEq E [⟨.Z, [w], [], dz, a⟩, ⟨.X, [], [w], dz, b⟩] [w] [] [] iface 1 (rt (E.dimv dz))
\end{lstlisting}
\begin{lstlisting}[style=leancore]
theorem colorState (E : Env) (a b dz w x : ℕ) (iface : List ℕ) (hd : 0 < E.dimv dz)
    (hw : labIdx E.labs w = dz) (hwx : w ≠ x) (ha : ∀ k < E.dimv dz, ex (E.phase a k) = 1)
    (hb : ∀ k < E.dimv dz, ex (E.phase b k) = 1) :
    LocalEq E [⟨.Z, [w], [], dz, a⟩, ⟨.F, [x], [w], dz, 0⟩] [w] [⟨.X, [x], [], dz, b⟩] [] iface
      1 1
\end{lstlisting}
\begin{lstlisting}[style=leancore]
theorem colorStateX (E : Env) (a b dz w x : ℕ) (iface : List ℕ) (hd : 0 < E.dimv dz)
    (hw : labIdx E.labs w = dz) (hwx : w ≠ x) (hx : x ∈ iface) (hxd : labIdx E.labs x = dz)
    (ha : ∀ k < E.dimv dz, ex (E.phase a k) = 1) (hb : ∀ k < E.dimv dz, ex (E.phase b k) = 1) :
    LocalEq E [⟨.X, [w], [], dz, a⟩, ⟨.F, [x], [w], dz, 0⟩] [w] [⟨.Z, [x], [], dz, b⟩] [] iface
      1 1
\end{lstlisting}
\begin{lstlisting}[style=leancore]
theorem stateCopy (E : Env) (a b c dz w : ℕ) (os iface : List ℕ) (k : ℂ) (hd : 0 < E.dimv dz)
    (hw : labIdx E.labs w = dz) (hwo : w ∉ os)
    (hos : ∀ o ∈ os, o ∈ iface ∧ labIdx E.labs o = dz)
    (ha : ∀ k < E.dimv dz, ex (E.phase a k) = 1) (hb : ∀ k < E.dimv dz, ex (E.phase b k) = 1)
    (hc : ∀ k < E.dimv dz, ex (E.phase c k) = 1)
    (hk : k * rt (E.dimv dz) ^ os.length = rt (E.dimv dz)) :
    LocalEq E [⟨.X, [w], [], dz, a⟩, ⟨.Z, os, [w], dz, b⟩] [w]
      (os.map fun o => ⟨.X, [o], [], dz, c⟩) [] iface 1 k
\end{lstlisting}
\begin{lstlisting}[style=leancore]
theorem triangle (E : Env) (dz w o i x y : ℕ) (iface : List ℕ)
    (hxy : (x = o ∧ y = i) ∨ (x = i ∧ y = o)) (hw : labIdx E.labs w = dz) (hwo : w ≠ o)
    (hwi : w ≠ i) (ho : o ∈ iface ∧ labIdx E.labs o = dz) (hi : i ∈ iface ∧ labIdx E.labs i = dz) :
    LocalEq E [⟨.T, [w], [i], dz, 0⟩, ⟨.Ti, [o], [w], dz, 0⟩] [w] [⟨.I, [x], [y], dz, 0⟩] []
      iface 1 1
\end{lstlisting}
\begin{lstlisting}[style=leancore]
theorem fourier4 (E : Env) (a b c dz w₁ w₂ w₃ : ℕ) (iface : List ℕ) (hd : 0 < E.dimv dz)
    (h₁ : labIdx E.labs w₁ = dz) (h₂ : labIdx E.labs w₂ = dz) (h₃ : labIdx E.labs w₃ = dz)
    (h12 : w₁ ≠ w₂) (h13 : w₁ ≠ w₃) (h23 : w₂ ≠ w₃) (ha1 : a ≠ w₁) (ha2 : a ≠ w₂)
    (ha3 : a ≠ w₃) (hb1 : b ≠ w₁) (hb2 : b ≠ w₂) (hb3 : b ≠ w₃)
    (ha : a ∈ iface ∧ labIdx E.labs a = dz) (hb : b ∈ iface ∧ labIdx E.labs b = dz)
    (hc : ∀ k < E.dimv dz, ex (E.phase c k) = 1) :
    LocalEq E [⟨.F, [w₁], [a], dz, 0⟩, ⟨.F, [w₂], [w₁], dz, 0⟩, ⟨.F, [w₃], [w₂], dz, 0⟩,
      ⟨.F, [b], [w₃], dz, 0⟩] [w₁, w₂, w₃] [⟨.Z, [b], [a], dz, c⟩] [] iface 1 1
\end{lstlisting}

\subsection{\texttt{Bialgebra}}
\begin{lstlisting}[style=leancore]
theorem bialgebra (E : Env) (px pz q₁ q₂ q₃ q₄ dz w a b c e n₀₀ n₀₁ n₁₀ n₁₁ : ℕ)
    (iface : List ℕ) (hd : 0 < E.dimv dz) (hw : labIdx E.labs w = dz)
    (hn : ∀ n ∈ [n₀₀, n₀₁, n₁₀, n₁₁], labIdx E.labs n = dz)
    (hnd : [n₀₀, n₀₁, n₁₀, n₁₁, a, b, c, e].Nodup) (hwd : [w, a, b, c, e].Nodup)
    (hf : ∀ x ∈ [a, b, c, e], x ∈ iface ∧ labIdx E.labs x = dz)
    (hpx : ∀ k < E.dimv dz, ex (E.phase px k) = 1) (hpz : ∀ k < E.dimv dz, ex (E.phase pz k) = 1)
    (h₁ : ∀ k < E.dimv dz, ex (E.phase q₁ k) = 1) (h₂ : ∀ k < E.dimv dz, ex (E.phase q₂ k) = 1)
    (h₃ : ∀ k < E.dimv dz, ex (E.phase q₃ k) = 1) (h₄ : ∀ k < E.dimv dz, ex (E.phase q₄ k) = 1) :
    LocalEq E [⟨.X, [w], [a, b], dz, px⟩, ⟨.Z, [c, e], [w], dz, pz⟩] [w]
      [⟨.Z, [n₀₀, n₀₁], [a], dz, q₁⟩, ⟨.Z, [n₁₀, n₁₁], [b], dz, q₂⟩,
        ⟨.X, [c], [n₀₀, n₁₀], dz, q₃⟩, ⟨.X, [e], [n₀₁, n₁₁], dz, q₄⟩]
      [n₀₀, n₀₁, n₁₀, n₁₁] iface 1 (rt (E.dimv dz))
\end{lstlisting}

\subsection{\texttt{Hopf}}
\begin{lstlisting}[style=leancore]
theorem hopf (E : Env) (pz px dz u v₁ v₂ v₃ : ℕ) (zo zi xo xi zo' zi' xo' xi' iface : List ℕ)
    (k : ℂ) (hd : 0 < E.dimv dz) (hz : (zo ++ zi).Perm (u :: v₁ :: (zo' ++ zi')))
    (hxo : xo.Perm xo') (hxi : xi.Perm (u :: v₃ :: xi')) (hnd : [u, v₁, v₂, v₃].Nodup)
    (hfree : ∀ x ∈ zo' ++ zi' ++ xo' ++ xi', x ∉ [u, v₁, v₂, v₃])
    (hdim : ∀ x ∈ [u, v₁, v₂, v₃], labIdx E.labs x = dz) (hk : k * E.dimv dz = 1) :
    LocalEq E [⟨.Z, zo, zi, dz, pz⟩, ⟨.X, xo, xi, dz, px⟩, ⟨.F, [v₂], [v₁], dz, 0⟩,
      ⟨.F, [v₃], [v₂], dz, 0⟩] [u, v₁, v₂, v₃] [⟨.Z, zo', zi', dz, pz⟩, ⟨.X, xo', xi', dz, px⟩]
      [] iface 1 k
\end{lstlisting}
\begin{lstlisting}[style=leancore]
theorem hopfSwapped (E : Env) (px pz dz u v₁ v₂ v₃ : ℕ)
    (xo xi zo zi xo' xi' zo' zi' iface : List ℕ) (k : ℂ) (hd : 0 < E.dimv dz)
    (hxo : xo.Perm (u :: v₁ :: xo')) (hxi : xi.Perm xi')
    (hz : (zo ++ zi).Perm (u :: v₃ :: (zo' ++ zi'))) (hnd : [u, v₁, v₂, v₃].Nodup)
    (hfree : ∀ x ∈ zo' ++ zi' ++ xo' ++ xi', x ∉ [u, v₁, v₂, v₃])
    (hdim : ∀ x ∈ [u, v₁, v₂, v₃], labIdx E.labs x = dz) (hk : k * E.dimv dz = 1) :
    LocalEq E [⟨.X, xo, xi, dz, px⟩, ⟨.Z, zo, zi, dz, pz⟩, ⟨.F, [v₂], [v₁], dz, 0⟩,
      ⟨.F, [v₃], [v₂], dz, 0⟩] [u, v₁, v₂, v₃] [⟨.X, xo', xi', dz, px⟩, ⟨.Z, zo', zi', dz, pz⟩]
      [] iface 1 k
\end{lstlisting}

\subsection{\texttt{WRules}}
\begin{lstlisting}[style=leancore]
theorem wFusion (E : Env) (dz w i : ℕ) (o₁ o₂ o₃ iface : List ℕ)
    (hw : labIdx E.labs w = dz) (h₁ : w ∈ o₁) (h₁' : w ∉ o₁.erase w) (h₂ : w ∉ o₂)
    (hwi : w ≠ i) (hi : i ∈ iface ∧ labIdx E.labs i = dz) (hd : 0 < E.dimv dz)
    (ho : o₃.Perm (o₁.erase w ++ o₂)) :
    LocalEq E [⟨.W, o₁, [i], dz, 0⟩, ⟨.W, o₂, [w], dz, 0⟩] [w] [⟨.W, o₃, [i], dz, 0⟩] [] iface
      1 1
\end{lstlisting}
\begin{lstlisting}[style=leancore]
theorem wIdentity (E : Env) (dz p q w x : ℕ) (iface : List ℕ)
    (hpq : (p = w ∧ q = x) ∨ (p = x ∧ q = w)) (hface : p ∈ iface ∧ q ∈ iface)
    (hp : labIdx E.labs p = dz) (hq : labIdx E.labs q = dz) :
    LocalEq E [⟨.W, [p], [q], dz, 0⟩] [] [⟨.I, [x], [w], dz, 0⟩] [] iface 1 1
\end{lstlisting}
\begin{lstlisting}[style=leancore]
theorem wZeroCopy (E : Env) (a c dz w : ℕ) (os iface : List ℕ) (k : ℂ) (hd : 0 < E.dimv dz)
    (hw : labIdx E.labs w = dz) (hwo : w ∉ os)
    (hos : ∀ o ∈ os, o ∈ iface ∧ labIdx E.labs o = dz)
    (ha : ∀ k < E.dimv dz, ex (E.phase a k) = 1) (hc : ∀ k < E.dimv dz, ex (E.phase c k) = 1)
    (hk : k * rt (E.dimv dz) ^ os.length = rt (E.dimv dz)) :
    LocalEq E [⟨.X, [w], [], dz, a⟩, ⟨.W, os, [w], dz, 0⟩] [w]
      (os.map fun o => ⟨.X, [o], [], dz, c⟩) [] iface 1 k
\end{lstlisting}
\begin{lstlisting}[style=leancore]
theorem triZeroState (E : Env) (a dz w o : ℕ) (iface : List ℕ) (hd : 0 < E.dimv dz)
    (hw : labIdx E.labs w = dz) (hwo : w ≠ o) (ho : o ∈ iface ∧ labIdx E.labs o = dz)
    (ha : ∀ k < E.dimv dz, ex (E.phase a k) = 1) :
    LocalEq E [⟨.X, [w], [], dz, a⟩, ⟨.T, [o], [w], dz, 0⟩] [w] [⟨.X, [o], [], dz, a⟩] []
      iface 1 1
\end{lstlisting}
\begin{lstlisting}[style=leancore]
theorem triZeroEffect (E : Env) (a c dz w i : ℕ) (iface : List ℕ) (hd : 0 < E.dimv dz)
    (hw : labIdx E.labs w = dz) (hwi : w ≠ i) (hi : i ∈ iface ∧ labIdx E.labs i = dz)
    (ha : ∀ k < E.dimv dz, ex (E.phase a k) = 1) (hc : ∀ k < E.dimv dz, ex (E.phase c k) = 1) :
    LocalEq E [⟨.T, [w], [i], dz, 0⟩, ⟨.X, [], [w], dz, a⟩] [w] [⟨.Z, [], [i], dz, c⟩] []
      iface 1 (rt (E.dimv dz))
\end{lstlisting}
\begin{lstlisting}[style=leancore]
theorem wZEffect (E : Env) (a dz w i v : ℕ) (os os' iface : List ℕ) (hd : 0 < E.dimv dz)
    (hw : labIdx E.labs w = dz) (hv : labIdx E.labs v = dz) (h₁ : w ∈ os)
    (h₁' : w ∉ os.erase w) (hwi : w ≠ i) (hvi : v ≠ i) (hvo : v ∉ os')
    (hi : i ∈ iface ∧ labIdx E.labs i = dz) (ho : os'.Perm (os.erase w))
    (hrest : w ∉ os') (ha : ∀ k < E.dimv dz, ex (E.phase a k) = 1) :
    LocalEq E [⟨.W, os, [i], dz, 0⟩, ⟨.Z, [], [w], dz, a⟩] [w]
      [⟨.T, [v], [i], dz, 0⟩, ⟨.W, os', [v], dz, 0⟩] [v] iface 1 1
\end{lstlisting}
\begin{lstlisting}[style=leancore]
theorem wZeroEffect (E : Env) (a dz w i : ℕ) (os os' iface : List ℕ) (hd : 0 < E.dimv dz)
    (hw : labIdx E.labs w = dz) (h₁ : w ∈ os) (h₁' : w ∉ os.erase w) (hwi : w ≠ i)
    (hi : i ∈ iface ∧ labIdx E.labs i = dz) (ho : os'.Perm (os.erase w))
    (ha : ∀ k < E.dimv dz, ex (E.phase a k) = 1) :
    LocalEq E [⟨.W, os, [i], dz, 0⟩, ⟨.X, [], [w], dz, a⟩] [w] [⟨.W, os', [i], dz, 0⟩] []
      iface 1 (rt (E.dimv dz))
\end{lstlisting}

\subsection{\texttt{WAssoc}}
\begin{lstlisting}[style=leancore]
noncomputable def sumW (d : ℕ) : ℕ → (List ℕ → ℂ) → ℂ
  | 0, f => f []
  | M + 1, f => ∑ x ∈ range d, sumW d M fun xs => f (x :: xs)
\end{lstlisting}
\begin{lstlisting}[style=leancore]
noncomputable def wTerm (d : ℕ) (w : ℕ → ℂ) : List ℕ → List (List ℕ) → ℂ
  | x :: xs, B :: Bs => w x * tW d B x * wTerm d w xs Bs
  | _, _ => 1
\end{lstlisting}
\begin{lstlisting}[style=leancore]
theorem wAssoc_nil (d : ℕ) (hd : 0 < d) (w : ℕ → ℂ) (hw : w 0 = 1) (i : ℕ) (hi : i < d)
    (Bs : List (List ℕ)) :
    sumW d Bs.length (fun xs => tW d xs i * wTerm d w xs Bs) = w i * tW d Bs.flatten i
\end{lstlisting}

\subsection{\texttt{Connective}}
\begin{lstlisting}[style=leancore]
theorem connStates (E : Env) (ds dt dk p₁ p₂ c a b k : ℕ) (iface : List ℕ)
    (hs : 0 < E.dimv ds) (ht : 0 < E.dimv dt) (hdk : E.dimv dk = E.dimv ds * E.dimv dt)
    (ha : labIdx E.labs a = ds) (hb : labIdx E.labs b = dt) (hab : a ≠ b) (hak : a ≠ k)
    (hbk : b ≠ k) (hk : k ∈ iface ∧ labIdx E.labs k = dk)
    (h₁ : ∀ k < E.dimv ds, ex (E.phase p₁ k) = 1) (h₂ : ∀ k < E.dimv dt, ex (E.phase p₂ k) = 1)
    (hc : ∀ k < E.dimv dk, ex (E.phase c k) = 1) :
    LocalEq E [⟨.B, [k], [a, b], ds, dt⟩, ⟨.Z, [a], [], ds, p₁⟩, ⟨.Z, [b], [], dt, p₂⟩] [a, b]
      [⟨.Z, [k], [], dk, c⟩] [] iface 1 1
\end{lstlisting}
\begin{lstlisting}[style=leancore]
theorem connInverse (E : Env) (ds dt dk a b k₁ k₂ x y : ℕ) (iface : List ℕ)
    (ht : 0 < E.dimv dt) (hdk : E.dimv dk = E.dimv ds * E.dimv dt)
    (hxy : (x = k₁ ∧ y = k₂) ∨ (x = k₂ ∧ y = k₁)) (ha : labIdx E.labs a = ds)
    (hb : labIdx E.labs b = dt) (hd : [a, b, k₁, k₂].Nodup)
    (h₁ : k₁ ∈ iface ∧ labIdx E.labs k₁ = dk) (h₂ : k₂ ∈ iface ∧ labIdx E.labs k₂ = dk) :
    LocalEq E [⟨.S, [a, b], [k₁], ds, dt⟩, ⟨.B, [k₂], [a, b], ds, dt⟩] [a, b]
      [⟨.I, [x], [y], dk, 0⟩] [] iface 1 1
\end{lstlisting}
\begin{lstlisting}[style=leancore]
theorem connInverseBind (E : Env) (ds dt dk k a₁ b₁ a₂ b₂ x₁ y₁ x₂ y₂ : ℕ) (iface : List ℕ)
    (hdk : E.dimv dk = E.dimv ds * E.dimv dt) (hk : labIdx E.labs k = dk)
    (hk' : k ∉ [a₁, b₁, a₂, b₂])
    (hxa : (x₁ = a₁ ∧ y₁ = a₂) ∨ (x₁ = a₂ ∧ y₁ = a₁))
    (hxb : (x₂ = b₁ ∧ y₂ = b₂) ∨ (x₂ = b₂ ∧ y₂ = b₁))
    (ha₁ : a₁ ∈ iface ∧ labIdx E.labs a₁ = ds) (ha₂ : a₂ ∈ iface ∧ labIdx E.labs a₂ = ds)
    (hb₁ : b₁ ∈ iface ∧ labIdx E.labs b₁ = dt) (hb₂ : b₂ ∈ iface ∧ labIdx E.labs b₂ = dt) :
    LocalEq E [⟨.B, [k], [a₁, b₁], ds, dt⟩, ⟨.S, [a₂, b₂], [k], ds, dt⟩] [k]
      [⟨.I, [x₁], [y₁], ds, 0⟩, ⟨.I, [x₂], [y₂], dt, 0⟩] [] iface 1 1
\end{lstlisting}

\subsection{\texttt{WZ}}
\begin{lstlisting}[style=leancore]
theorem wzBialgebra (E : Env) (pz q dz w : ℕ) (zi wo iface : List ℕ)
    (rows cols : List (List ℕ)) (pin : List (ℕ × ℕ)) (hd : 0 < E.dimv dz)
    (hw : labIdx E.labs w = dz) (hzi : zi ≠ []) (hwz : w ∉ zi) (hwo : w ∉ wo)
    (hrange : ∀ x ∈ zi ++ wo, x ∈ iface ∧ labIdx E.labs x = dz)
    (hpin : rows.flatten = pin.map Prod.fst) (hnd : rows.flatten.Nodup)
    (hto : ∀ p ∈ pin, p.2 ∈ wo) (hdisj : ∀ x ∈ zi ++ wo ++ [w], x ∉ rows.flatten)
    (hcol : ((wo.zip cols).flatMap fun p => p.2.map fun e => (e, p.1)).Perm pin)
    (hrow : ∀ r ∈ rows, r.map (ren pin) = wo) (hlen : zi.length = rows.length)
    (hcl : wo.length = cols.length) (hdims : ∀ e ∈ rows.flatten, labIdx E.labs e = dz)
    (hpz : ∀ k < E.dimv dz, ex (E.phase pz k) = 1) (hq : ∀ k < E.dimv dz, ex (E.phase q k) = 1) :
    LocalEq E [⟨.Z, [w], zi, dz, pz⟩, ⟨.W, wo, [w], dz, 0⟩] [w]
      ((zi.zip rows).map (fun p => ⟨.W, p.2, [p.1], dz, 0⟩) ++
        (wo.zip cols).map (fun p => ⟨.Z, [p.1], p.2, dz, q⟩)) rows.flatten iface 1 1
\end{lstlisting}

\subsection{\texttt{Completeness}}
\begin{lstlisting}[style=leancore]
theorem wPerm (E : Env) (a b dz i : ℕ) (os os' iface : List ℕ) (h : os.Perm os') :
    LocalEq E [⟨.W, os, [i], dz, a⟩] [] [⟨.W, os', [i], dz, b⟩] [] iface 1 1
\end{lstlisting}
\begin{lstlisting}[style=leancore]
theorem zPerm (E : Env) (a dz : ℕ) (os is os' is' iface : List ℕ)
    (h : (os ++ is).Perm (os' ++ is')) :
    LocalEq E [⟨.Z, os, is, dz, a⟩] [] [⟨.Z, os', is', dz, a⟩] [] iface 1 1
\end{lstlisting}
\begin{lstlisting}[style=leancore]
theorem wAso (E : Env) (dz w w' i : ℕ) (o₁ o₂ p₁ p₂ o₃ iface : List ℕ)
    (hw : labIdx E.labs w = dz) (h₁ : w ∈ o₁) (h₁' : w ∉ o₁.erase w) (h₂ : w ∉ o₂)
    (hwi : w ≠ i) (hi : i ∈ iface ∧ labIdx E.labs i = dz) (hd : 0 < E.dimv dz)
    (ho : o₃.Perm (o₁.erase w ++ o₂))
    (gw : labIdx E.labs w' = dz) (g₁ : w' ∈ p₁) (g₁' : w' ∉ p₁.erase w') (g₂ : w' ∉ p₂)
    (gwi : w' ≠ i) (go : o₃.Perm (p₁.erase w' ++ p₂)) :
    LocalEq E [⟨.W, o₁, [i], dz, 0⟩, ⟨.W, o₂, [w], dz, 0⟩] [w]
      [⟨.W, p₁, [i], dz, 0⟩, ⟨.W, p₂, [w'], dz, 0⟩] [w'] iface 1 1
\end{lstlisting}
\begin{lstlisting}[style=leancore]
theorem kCopy (E : Env) (a b c dz w j : ℕ) (os iface : List ℕ) (k : ℂ) (hd : 0 < E.dimv dz)
    (hw : labIdx E.labs w = dz) (hwo : w ∉ os)
    (hos : ∀ o ∈ os, o ∈ iface ∧ labIdx E.labs o = dz)
    (ha : ∀ q < E.dimv dz, ex (E.phase a q) = om (E.dimv dz) ^ (q * j))
    (hb : ∀ q < E.dimv dz, ex (E.phase b q) = 1)
    (hc : ∀ q < E.dimv dz, ex (E.phase c q) = om (E.dimv dz) ^ (q * j))
    (hk : k * rt (E.dimv dz) ^ os.length = rt (E.dimv dz)) :
    LocalEq E [⟨.X, [w], [], dz, a⟩, ⟨.Z, os, [w], dz, b⟩] [w]
      (os.map fun o => ⟨.X, [o], [], dz, c⟩) [] iface 1 k
\end{lstlisting}

\subsection{\texttt{Weights}}
A representative selection; the other rule lemmas of the module follow the same pattern.
\begin{lstlisting}[style=leancore]
noncomputable def zc (E : Env) (n : Node) (k : ℕ) : ℂ :=
  match n.kind with
  | .Z => 1 * ex (E.phase n.ph k)
  | .ZW => E.coef n.ph k
  | _ => 0
\end{lstlisting}
\begin{lstlisting}[style=leancore]
theorem zcCongr (E : Env) (n m : Node) (iface : List ℕ) (hn : isZc n = true)
    (hm : isZc m = true) (hd : m.dim = n.dim) (hl : (n.outs ++ n.ins).Perm (m.outs ++ m.ins))
    (hc : ∀ k < E.dimv n.dim, zc E m k = zc E n k) :
    LocalEq E [n] [] [m] [] iface 1 1
\end{lstlisting}
\begin{lstlisting}[style=leancore]
theorem zcFusion (E : Env) (na nb m : Node) (dz w : ℕ) (iface : List ℕ) (ha : isZc na = true)
    (hb : isZc nb = true) (hm : isZc m = true) (hda : na.dim = dz) (hdb : nb.dim = dz)
    (hdm : m.dim = dz) (h₁ : w ∈ na.outs ++ na.ins) (h₂ : w ∈ nb.outs ++ nb.ins)
    (hl : (m.outs ++ m.ins).Perm ((na.outs ++ na.ins ++ (nb.outs ++ nb.ins)).filter (· ≠ w)))
    (hw : labIdx E.labs w = dz)
    (hc : ∀ k < E.dimv dz, zc E m k = zc E na k * zc E nb k) :
    LocalEq E [na, nb] [w] [m] [] iface 1 1
\end{lstlisting}
\begin{lstlisting}[style=leancore]
theorem zcScalar (E : Env) (n : Node) (iface : List ℕ) (s : ℂ) (hn : isZc n = true)
    (ho : n.outs = []) (hi : n.ins = []) (hs : s = ∑ k ∈ range (E.dimv n.dim), zc E n k) :
    LocalEq E [n] [] [] [] iface 1 s
\end{lstlisting}
\begin{lstlisting}[style=leancore]
theorem zcAd (E : Env) (dz x y i : ℕ) (na nb m : Node) (iface : List ℕ) (ha : isZc na = true)
    (hb : isZc nb = true) (hm : isZc m = true) (hda : na.dim = dz) (hdb : nb.dim = dz)
    (hdm : m.dim = dz) (hla : na.outs ++ na.ins = [x]) (hlb : nb.outs ++ nb.ins = [y])
    (hlm : m.outs ++ m.ins = [i]) (hxy : x ≠ y) (hxi : x ≠ i) (hyi : y ≠ i)
    (hx : labIdx E.labs x = dz) (hy : labIdx E.labs y = dz)
    (hi : i ∈ iface ∧ labIdx E.labs i = dz)
    (h0 : zc E m 0 = zc E na 0 * zc E nb 0)
    (hc : ∀ k < E.dimv dz, k ≠ 0 → zc E m k = zc E na k * zc E nb 0 + zc E na 0 * zc E nb k) :
    LocalEq E [⟨.W, [x, y], [i], dz, 0⟩, na, nb] [x, y] [m] [] iface 1 1
\end{lstlisting}
\begin{lstlisting}[style=leancore]
theorem zcK2 (E : Env) (dz p j w w' o i : ℕ) (na nb : Node) (s : ℂ) (iface : List ℕ)
    (ha : isZc na = true) (hb : isZc nb = true) (hda : na.dim = dz) (hdb : nb.dim = dz)
    (hla : na.outs ++ na.ins = [o, w]) (hlb : nb.outs ++ nb.ins = [w', i])
    (hw : labIdx E.labs w = dz) (hw' : labIdx E.labs w' = dz)
    (ho : o ∈ iface ∧ labIdx E.labs o = dz) (hi : i ∈ iface ∧ labIdx E.labs i = dz)
    (hwo : w ≠ o) (hwi : w ≠ i) (hw'o : w' ≠ o) (hw'i : w' ≠ i) (hd : 0 < E.dimv dz)
    (hp : ∀ q < E.dimv dz, ex (E.phase p q) = om (E.dimv dz) ^ (q * j))
    (hc : ∀ k < E.dimv dz, zc E na k = s * zc E nb ((k + j) % E.dimv dz)) :
    LocalEq E [⟨.X, [w], [i], dz, p⟩, na] [w] [nb, ⟨.X, [o], [w'], dz, p⟩] [w'] iface 1 s
\end{lstlisting}
\begin{lstlisting}[style=leancore]
theorem zcTi (E : Env) (dz e o i : ℕ) (m : Node) (iface : List ℕ) (hm : isZc m = true)
    (hdm : m.dim = dz) (hl : m.outs ++ m.ins = [e]) (he : labIdx E.labs e = dz)
    (ho : o ∈ iface ∧ labIdx E.labs o = dz) (hi : i ∈ iface ∧ labIdx E.labs i = dz)
    (heo : e ≠ o) (hei : e ≠ i) (h0 : zc E m 0 = 1)
    (hc : ∀ k < E.dimv dz, k ≠ 0 → zc E m k = -1) :
    LocalEq E [⟨.W, [e, o], [i], dz, 0⟩, m] [e] [⟨.Ti, [o], [i], dz, 0⟩] [] iface 1 1
\end{lstlisting}
\begin{lstlisting}[style=leancore]
theorem fourierMult (E : Env) (pz px dz a b w : ℕ) (ws iface : List ℕ) (hd : 0 < E.dimv dz)
    (hnd : ws.Nodup) (ha : a ∉ ws) (hb : b ∉ ws) (hw : w ≠ a ∧ w ≠ b)
    (hab : a ∈ iface ∧ labIdx E.labs a = dz) (hbb : b ∈ iface ∧ labIdx E.labs b = dz)
    (hwl : labIdx E.labs w = dz) (hws : ∀ x ∈ ws, labIdx E.labs x = dz) (hne : ws ≠ [])
    (he : (ws.length + 1) % E.dimv dz = 0)
    (hz : ∀ k < E.dimv dz, ex (E.phase pz k) = 1) (hx : ∀ k < E.dimv dz, ex (E.phase px k) = 1)
    (s : ℂ) (hs : s = rt (E.dimv dz) ^ (ws.length - 1)) :
    LocalEq E [⟨.F, [w], [a], dz, 0⟩, ⟨.F, [b], [w], dz, 0⟩] [w]
      [⟨.Z, ws, [a], dz, pz⟩, ⟨.X, [b], ws, dz, px⟩] ws iface 1 s
\end{lstlisting}
\begin{lstlisting}[style=leancore]
theorem bindDef (E : Env) (ds dt dk a b c x y z p₁ p₂ p₃ : ℕ) (iface : List ℕ)
    (hdk : E.dimv dk = E.dimv ds * E.dimv dt)
    (hnd : [x, y, z, a, b, c].Nodup) (hx : labIdx E.labs x = ds) (hy : labIdx E.labs y = dt)
    (hz : labIdx E.labs z = dk) (ha : a ∈ iface ∧ labIdx E.labs a = ds)
    (hb : b ∈ iface ∧ labIdx E.labs b = dt) (hc : c ∈ iface ∧ labIdx E.labs c = dk)
    (h₁ : ∀ k < E.dimv ds, ex (E.phase p₁ k) = 1) (h₂ : ∀ k < E.dimv dt, ex (E.phase p₂ k) = 1)
    (h₃ : ∀ k < E.dimv dk, ex (E.phase p₃ k) = 1) :
    LocalEq E [⟨.B, [c], [a, b], ds, dt⟩] []
      [⟨.S, [x, y], [z], ds, dt⟩, ⟨.Z, [x], [a], ds, p₁⟩, ⟨.Z, [y], [b], dt, p₂⟩,
        ⟨.Z, [c], [z], dk, p₃⟩] [x, y, z] iface 1 1
\end{lstlisting}
\begin{lstlisting}[style=leancore]
theorem zxwD1 (E : Env) (dz i o w v v' w' : ℕ) (m m' : Node) (iface : List ℕ)
    (hd : 0 < E.dimv dz) (hm : isZc m = true) (hm' : isZc m' = true) (hdm : m.dim = dz)
    (hdm' : m'.dim = dz) (hl : m.outs ++ m.ins = [w, i]) (hl' : m'.outs ++ m'.ins = [o, w'])
    (hnd : [w, v, v', w', i, o].Nodup) (hw : ∀ x ∈ [w, v, v', w'], labIdx E.labs x = dz)
    (hi : i ∈ iface ∧ labIdx E.labs i = dz) (ho : o ∈ iface ∧ labIdx E.labs o = dz)
    (hc : ∀ k < E.dimv dz, zc E m' k = zc E m ((E.dimv dz - k) % E.dimv dz)) :
    LocalEq E [m, ⟨.F, [v], [w], dz, 0⟩, ⟨.F, [o], [v], dz, 0⟩] [w, v]
      [⟨.F, [v'], [i], dz, 0⟩, ⟨.F, [w'], [v'], dz, 0⟩, m'] [v', w'] iface 1 1
\end{lstlisting}
\begin{lstlisting}[style=leancore]
def MixL.nodes (l : MixL) : List Node :=
  if l.out then [⟨.X, [], [l.a], l.dr, l.p⟩, ⟨.S, [l.a, l.v], [l.K], l.dr, l.dt⟩]
  else [⟨.X, [l.a], [], l.dr, l.p⟩, ⟨.B, [l.K], [l.a, l.v], l.dr, l.dt⟩]
\end{lstlisting}
\begin{lstlisting}[style=leancore]
theorem mixZ (E : Env) (dD c : ℕ) (legs : List MixL) (ko ki : List ℕ) (σ : ℕ → ℕ)
    (hperm : (ko ++ ki).Perm (legs.map MixL.K))
    (hnd : (legs.flatMap fun l => [l.a, l.K]).Nodup)
    (hv : ∀ l ∈ legs, l.v ∉ legs.flatMap fun l => [l.a, l.K])
    (hr : ∀ l ∈ legs, 0 < E.dimv l.dr)
    (hφ : ∀ l ∈ legs, ∀ k < E.dimv l.dr, ex (E.phase l.p k) = 1)
    (ha : ∀ l ∈ legs, labDim E l.a = E.dimv l.dr)
    (hK : ∀ l ∈ legs, σ l.v < labDim E l.K) (hV : ∀ l ∈ legs, σ l.v < E.dimv l.dt) :
    sumOver (labDim E) (legs.flatMap fun l => [l.a, l.K])
      (prodVal E (legs.flatMap MixL.nodes ++ [⟨.ZW, ko, ki, dD, c⟩])) σ =
      (legs.map fun l => rt (E.dimv l.dr)).prod *
        tZw (E.dimv dD) (E.coef c) (legs.map fun l => σ l.v)
\end{lstlisting}
\begin{lstlisting}[style=leancore]
theorem qfS1 (E : Env) (LA LB LF : List MixL) (mA mB mF : Node) (W : List ℕ) (iface : List ℕ)
    (s : ℂ) (hA : MixOK E mA LA) (hB : MixOK E mB LB) (hF : MixOK E mF LF)
    (hW : W ≠ []) (hWnd : W.Nodup) (hWA : ∀ w ∈ W, w ∈ LA.map (·.v))
    (hWB : ∀ w ∈ W, w ∈ LB.map (·.v))
    (hWd : ∀ l ∈ LA ++ LB, l.v ∈ W → labDim E l.v = E.dimv l.dt)
    (hfree : ∀ l ∈ LA ++ LB, l.v ∉ W → l.v ∈ iface ∧ labDim E l.v = E.dimv l.dt)
    (hFv : (LF.map (·.v)).Perm (((LA ++ LB).filter fun l => l.v ∉ W).map (·.v)))
    (hFd : ∀ l ∈ LF, ∃ l' ∈ LA ++ LB, l'.v = l.v ∧ E.dimv l'.dt = E.dimv l.dt)
    (hLF : LF ≠ [])
    (hAB : ∀ x ∈ labelsOf (LA.flatMap MixL.nodes ++ [mA]), x ∉ LB.flatMap fun l => [l.a, l.K])
    (hBA : ∀ x ∈ labelsOf (LB.flatMap MixL.nodes ++ [mB]), x ∉ LA.flatMap fun l => [l.a, l.K])
    (hcA : ∀ j < E.dimv mA.dim, (∃ l ∈ LA, E.dimv l.dt ≤ j) → zc E mA j = 0)
    (hcB : ∀ j < E.dimv mB.dim, (∃ l ∈ LB, E.dimv l.dt ≤ j) → zc E mB j = 0)
    (hcF : ∀ j < E.dimv mF.dim, (∀ l ∈ LF, j < E.dimv l.dt) →
      zc E mF j = if ∀ l ∈ LA ++ LB, j < E.dimv l.dt then zc E mA j * zc E mB j else 0)
    (hs : s * (LF.map fun l => rt (E.dimv l.dr)).prod =
      (LA.map fun l => rt (E.dimv l.dr)).prod * (LB.map fun l => rt (E.dimv l.dr)).prod) :
    LocalEq E ((LA.flatMap MixL.nodes ++ [mA]) ++ (LB.flatMap MixL.nodes ++ [mB]))
      (W ++ ((LA.flatMap fun l => [l.a, l.K]) ++ LB.flatMap fun l => [l.a, l.K]))
      (LF.flatMap MixL.nodes ++ [mF]) (LF.flatMap fun l => [l.a, l.K]) iface 1 s
\end{lstlisting}

\subsection{\texttt{Unfuse}}
\begin{lstlisting}[style=leancore]
theorem unfuseZB (E : Env) (cs : List (List ℕ)) (c : List ℕ) (a cph dz y b : ℕ)
    (os hd tl o₂ i₂ iface : List ℕ) (htl : tl.length = cs.length)
    (hy : labIdx E.labs y = dz) (hyl : y ∉ os ++ hd ++ tl ++ [b])
    (hm : (o₂ ++ i₂).Perm (os ++ hd ++ [y]))
    (hc : ∀ k < E.dimv dz, ex (E.phase cph k) = 1) :
    LocalEq E [⟨.ZB (cs ++ [c]), os, hd ++ tl ++ [b], dz, a⟩] []
      [⟨zk cs, o₂, i₂ ++ tl, dz, a⟩, ⟨.ZB [c], [], [y, b], dz, cph⟩] [y] iface 1 1
\end{lstlisting}
\begin{lstlisting}[style=leancore]
theorem unfuseXB (E : Env) (co ci : List (List ℕ)) (c : List ℕ) (a cph dz y b : ℕ)
    (hd tl is o₂ iface : List ℕ) (htl : tl.length = co.length) (hd0 : 0 < E.dimv dz)
    (hy : labIdx E.labs y = dz) (hyl : y ∉ hd ++ tl ++ [b] ++ is)
    (ho : o₂.Perm (hd ++ [y]))
    (hc : ∀ k < E.dimv dz, ex (E.phase cph k) = 1) :
    LocalEq E [⟨.XB (co ++ [c]) ci, hd ++ tl ++ [b], is, dz, a⟩] []
      [⟨xk co ci, o₂ ++ tl, is, dz, a⟩, ⟨.XB [c] [], [b], [y], dz, cph⟩] [y] iface 1 1
\end{lstlisting}

\subsection{\texttt{FuseB}}
\begin{lstlisting}[style=leancore]
theorem fusionZG (E : Env) (cs₁ cs₂ : List (List ℕ)) (a b c dz w : ℕ)
    (o₁ i₁ t₁ o₂ i₂ t₂ o₃ i₃ iface : List ℕ) (ht₁ : t₁.length = cs₁.length)
    (ht₂ : t₂.length = cs₂.length)
    (h₁ : w ∈ o₁ ++ i₁) (h₁' : w ∉ (o₁ ++ i₁).erase w)
    (h₂ : w ∈ o₂ ++ i₂) (h₂' : w ∉ (o₂ ++ i₂).erase w) (hwt : w ∉ t₁ ++ t₂)
    (hm : (o₃ ++ i₃).Perm ((o₁ ++ i₁).erase w ++ (o₂ ++ i₂).erase w))
    (hw : labIdx E.labs w = dz)
    (hph : ∀ k < E.dimv dz, ex (E.phase c k) = ex (E.phase a k) * ex (E.phase b k)) :
    LocalEq E [⟨zk cs₁, o₁, i₁ ++ t₁, dz, a⟩, ⟨zk cs₂, o₂, i₂ ++ t₂, dz, b⟩] [w]
      [⟨zk (cs₁ ++ cs₂), o₃, i₃ ++ (t₁ ++ t₂), dz, c⟩] [] iface 1 1
\end{lstlisting}
\begin{lstlisting}[style=leancore]
theorem bunI (E : Env) (ms : List ℕ) (dz bz x w : ℕ) (iface : List ℕ)
    (hb : E.dimv bz = E.dimv dz ^ (ms.map E.mult).prod) (hd : 0 < E.dimv dz)
    (hx : labIdx E.labs x = bz) (hw : labIdx E.labs w = bz) (hxi : x ∈ iface) (hwi : w ∈ iface) :
    LocalEq E [⟨Kind.tower ms, [x], [w], dz, 0⟩] [] [⟨.I, [x], [w], bz, 0⟩] [] iface 1 1
\end{lstlisting}

\subsection{\texttt{WFuse}}
\begin{lstlisting}[style=leancore]
theorem wFuseB (E : Env) (b : Kind) (hb : b = .Z ∨ b = .ZW ∨ b = .W) (w : ℕ → ℂ)
    (dz m bz pw pb pz pw' t a a' : ℕ) (ys : List ℕ)
    (hd : 0 < E.dimv dz) (hbz : E.dimv bz = E.dimv dz ^ ([m].map E.mult).prod)
    (ht : labIdx E.labs t = bz)
    (ha : labIdx E.labs a = dz) (ha' : labIdx E.labs a' = dz)
    (hta : t ≠ a) (hty : t ∉ ys) (haa : a' ≠ a) (hay : a' ∉ ys)
    (hcopy : ∀ x < E.dimv dz, ∀ vs : List ℕ, vs.length = ys.length →
      tensor E b (E.dimv dz) pb vs [x] = w x * tW (E.dimv dz) vs x)
    (hw0 : w 0 = 1) (hz : ∀ k < E.dimv dz, E.coef pz k = w k) :
    LocalEq E [⟨.WB [[m]], [t], [a], dz, pw⟩, ⟨.Bun m b, ys, [t], dz, pb⟩] [t]
      [⟨.ZW, [a'], [a], dz, pz⟩, ⟨.WB (ys.map fun _ => [m]), ys, [a'], dz, pw'⟩] [a']
      (a :: ys) 1 1
\end{lstlisting}

\subsection{\texttt{Mixed}}
\begin{lstlisting}[style=leancore]
noncomputable def MNet.cpm (E : Env) (M : MNet) (σ σ' : ℕ → ℕ) : ℂ :=
  envSum (labDim E) M.env (fun a b => eval E M.net a * (starRingEnd ℂ) (eval E M.net b)) σ σ'
\end{lstlisting}
\begin{lstlisting}[style=leancore]
theorem Equal.cpm {E : Env} {s₁ s₂ : ℂ} {N₁ N₂ : Net} (h : Equal E s₁ N₁ s₂ N₂)
    (env : List ℕ) (σ σ' : ℕ → ℕ) (hσ : InRange E N₁.bnd σ) (hσ' : InRange E N₁.bnd σ') :
    s₁ * (starRingEnd ℂ) s₁ * (MNet.mk N₁ env).cpm E σ σ' =
      s₂ * (starRingEnd ℂ) s₂ * (MNet.mk N₂ env).cpm E σ σ'
\end{lstlisting}
\begin{lstlisting}[style=leancore]
theorem cpm_phase {E : Env} {s : ℂ} {N₁ N₂ : Net} (h : Equal E 1 N₁ s N₂)
    (hs : s * (starRingEnd ℂ) s = 1) (env : List ℕ) (σ σ' : ℕ → ℕ)
    (hσ : InRange E N₁.bnd σ) (hσ' : InRange E N₁.bnd σ') :
    (MNet.mk N₁ env).cpm E σ σ' = (MNet.mk N₂ env).cpm E σ σ'
\end{lstlisting}

\subsection{\texttt{Cyclo}}
\begin{lstlisting}[style=leancore]
def pmodStep (phi a : List ℤ) : List ℤ :=
  if a.length < phi.length then a
  else match a.getLast? with
    | some top => ktrim (kadd a (List.replicate (a.length - phi.length) 0 ++ pscale (-top) phi))
    | none => a
\end{lstlisting}
\begin{lstlisting}[style=leancore]
def pmod (phi : List ℤ) : ℕ → List ℤ → List ℤ
  | 0, a => a
  | n + 1, a => pmod phi n (pmodStep phi (ktrim a))
\end{lstlisting}
\begin{lstlisting}[style=leancore]
theorem cyc_ne (M : ℕ) (hM : 0 < M) (r : List ℤ) (hlen : r.length ≤ M.totient)
    (hnz : r.any (· ≠ 0) = true) : cvz (om M) r ≠ 0
\end{lstlisting}
\begin{lstlisting}[style=leancore]
theorem cvz_eq_zero_of_pmod (z : ℂ) (phi a : List ℤ) (h : cvz z phi = 0) (n : ℕ)
    (hz : (pmod phi n a).all (· == 0) = true) : cvz z a = 0
\end{lstlisting}
\begin{lstlisting}[style=leancore]
theorem ex_atoms (M p : ℕ) (z : ℤ) (gs : List ℝ) (ks : List ℤ) {θ : ℝ} (hM : 0 < M)
    (h : θ = p / M + z + (List.zipWith (fun (g : ℝ) (k : ℤ) => (k : ℝ) * g) gs ks).sum) :
    ex θ = om M ^ p * mono (gs.map ex) ks
\end{lstlisting}

\subsection{\texttt{GaussSign}}
\begin{lstlisting}[style=leancore]
theorem gauss_rt (p N : ℕ) (hp : 1 < p) (hodd : p % 2 = 1)
    (hprod : (pmod (List.replicate p 1) N
      (kadd (lgVec p) (pscale (-1) (gprod p ((p - 1) / 2))))).all (· == 0) = true)
    (hsq : (pmod (List.replicate p 1) N
      (kadd (pmul (lgVec p) (lgVec p)) [-((-1) ^ ((p - 1) / 2) * p)])).all (· == 0) = true) :
    ∑ k ∈ range p, (QSem.lg p k : ℂ) * om p ^ k = I ^ ((p - 1) / 2) * (-1) ^ gneg p * rt p
\end{lstlisting}

\subsection{\texttt{Ball}}
\begin{lstlisting}[style=leancore]
theorem near_mul {a b : ℂ} {x y r u v s : ℝ} (ha : Near a x y r) (hb : Near b u v s) :
    Near (a * b) (x * u - y * v) (x * v + y * u) ((|x| + |y|) * s + (|u| + |v|) * r + r * s)
\end{lstlisting}
\begin{lstlisting}[style=leancore]
theorem near_exp (a A ε : ℝ) (n : ℕ) (hn : 0 < n) (ha : |a - A| ≤ ε) (hA : |A| ≤ 1) :
    Near (Complex.exp ((a : ℂ) * I)) (cosT A n) (sinT A n)
      (2 * (ε + |A| ^ (2 * n) * (((2 * n).succ : ℝ) * (((2 * n).factorial : ℝ) *
        ((2 * n : ℕ) : ℝ))⁻¹)))
\end{lstlisting}

\subsection{\texttt{NumField}}
\begin{lstlisting}[style=leancore]
theorem root_unique (cs : List ℤ) (lo hi : ℝ) (hlh : lo < hi) (hlo : pev cs lo < 0) (hhi : 0 < pev cs hi)
    (hd : ∀ x, lo < x → x < hi → 0 < pdv cs x) :
    ∃! x, lo < x ∧ x < hi ∧ pev cs x = 0
\end{lstlisting}
\begin{lstlisting}[style=leancore]
theorem netCoef (E : Env) (A B : Net) (hi : A.int = B.int) (hb : A.bnd = B.bnd)
    (h : List.Forall₂ (fun n m => ∀ σ, val E σ n = val E σ m) A.nodes B.nodes) :
    Equal E 1 A 1 B
\end{lstlisting}

\subsection{\texttt{Refute}}
\begin{lstlisting}[style=leancore]
def gadd (M : ℕ) (k : List ℤ) (e : ℕ) (c : ℤ) : Grp → Grp
  | [] => [(k, addAt (List.replicate M 0) (e % M) c)]
  | g :: gs =>
    if g.1 = k then (g.1, addAt g.2 (e % M) c) :: gs
    else if klt k g.1 then (k, addAt (List.replicate M 0) (e % M) c) :: g :: gs
    else g :: gadd M k e c gs
\end{lstlisting}
\begin{lstlisting}[style=leancore]
def tXT (M : ℕ) (P : ℕ → ℕ × List ℤ) (d : ℕ) (os is : List ℕ) : Terms :=
  (List.range d).map fun k =>
    (1, (P k).1 + M / d * (k * os.sum) + M / d * ((d - 1) * (k * is.sum)), (P k).2)
\end{lstlisting}
\begin{lstlisting}[style=leancore]
def nodeOK (E : Shape) (M : ℕ) (keys : List (ℕ × ℕ)) (n : Node) : Bool :=
  isBS n || (n.kind.noZW && M % E.dimv n.dim == 0 &&
    (List.range (E.dimv n.dim)).all fun j => keys.contains (n.ph, j))
\end{lstlisting}
\begin{lstlisting}[style=leancore]
theorem eval_cyc (E : Env) (M : ℕ) (gs : List ℝ) (P : ℕ → ℕ → ℕ × List ℤ)
    (keys : List (ℕ × ℕ)) (N : Net) (σ : ℕ → ℕ) (hM : 0 < M)
    (hok : N.nodes.all (nodeOK E.shape M keys) = true) (hph : PhOK E M gs P keys) :
    eval E N σ = sqK E.shape N * gval (om M) (gs.map ex) (evalG E.shape M P N σ)
\end{lstlisting}
\begin{lstlisting}[style=leancore]
theorem refute {E : Env} {s₁ s₂ : ℂ} {N₁ N₂ : Net} (σ : ℕ → ℕ) (hin : InRange E N₁.bnd σ)
    (h : s₁ * eval E N₁ σ ≠ s₂ * eval E N₂ σ) : ¬ Equal E s₁ N₁ s₂ N₂
\end{lstlisting}

\subsection{\texttt{Contract}}
\begin{lstlisting}[style=leancore]
def elimOK (M : ℕ) (D : ℕ → ℕ) (l : ℕ) (Fs Fs' : List Fac) : Bool :=
  let ins := Fs.filter fun F => F.scope.contains l
  let outs := Fs.filter fun F => !F.scope.contains l
  match Fs'.reverse with
  | [] => false
  | new :: rest =>
    rest.reverse == outs && !new.scope.contains l &&
      (new.scope.all fun x => ins.any fun F => F.scope.contains x) &&
      (ins.all fun F => F.scope.all fun x => x == l || new.scope.contains x) &&
      (assigns D new.scope).all fun τ => new.at D τ ==
        compact M ((List.range (D l)).flatMap fun j => ffold D ins (Function.update τ l j))
\end{lstlisting}
\begin{lstlisting}[style=leancore]
def scopesOK (U : List ℕ) : List Fac → List (ℕ × List Fac) → Bool
  | Fs, [] => Fs.all fun F => F.scope.all (U.contains ·)
  | Fs, (l, Fs') :: steps =>
    (Fs.all fun F => F.scope.all fun x => U.contains x || x == l ||
      (steps.map Prod.fst).contains x) && scopesOK U Fs' steps
\end{lstlisting}
\begin{lstlisting}[style=leancore]
theorem eval_ve_at (E : Env) (M : ℕ) (gs : List ℝ) (P : ℕ → ℕ → ℕ × List ℤ)
    (keys : List (ℕ × ℕ)) (N : Net) (hM : 0 < M)
    (hok : N.nodes.all (nodeOK E.shape M keys) = true) (hph : PhOK E M gs P keys)
    (σ : ℕ → ℕ) (leaves : List Fac) (hlen : leaves.length = N.nodes.length)
    (hleaf : (N.nodes.zip leaves).all (fun p => leafAt E.shape M P N.bnd σ p.1 p.2) = true)
    (steps : List (ℕ × Fac)) (final : List Fac)
    (hperm : N.int.Perm (steps.map Prod.fst).reverse)
    (hchain : chainNew M E.shape.dim leaves steps final = true)
    (hscope : scopesOK N.bnd leaves (expand leaves steps) = true)
    (hdis : N.bnd.all (fun x => !(steps.map Prod.fst).contains x) = true)
    (hσ : ∀ x ∈ N.bnd, σ x < E.shape.dim x) :
    eval E N σ = sqK E.shape N * fprod (om M) (gs.map ex) E.shape.dim final σ
\end{lstlisting}

\subsection{\texttt{Diagram}}
\begin{lstlisting}[style=leancore]
def check (L : List (ℕ × ℕ)) (off : ℕ) : Bool :=
  decide (D.bnd.length ≤ off) &&
  D.nodes.all (fun n => (List.range n.outs.length).all (fun k => D.okPort ⟨n.id, true, k⟩) &&
    (List.range n.ins.length).all (fun k => D.okPort ⟨n.id, false, k⟩)) &&
  D.wires.all (fun w => D.portDim w.1 = D.portDim w.2) &&
  (List.range D.wires.length).all (fun j => L.lookup (off + j) = some (D.wireDim j)) &&
  decide ((List.range D.bnd.length).map (fun i => L.lookup i) = D.bdims.map some)
\end{lstlisting}
\begin{lstlisting}[style=leancore]
theorem eval_toNet (E : Env) (off : ℕ) (hc : D.check E.labs off = true) (σ : ℕ → ℕ) :
    ArchytasZX.eval E (D.toNet off) σ = D.eval E σ
\end{lstlisting}
\begin{lstlisting}[style=leancore]
theorem DEqual_iff {E : Env} {s₁ s₂ : ℂ} {D₁ D₂ : Diagram} {o₁ o₂ : ℕ} {L : List (ℕ × ℕ)}
    (hL : E.labs = L) (h₁ : D₁.check L o₁ = true) (h₂ : D₂.check L o₂ = true) :
    DEqual E s₁ D₁ s₂ D₂ ↔ Equal E s₁ (D₁.toNet o₁) s₂ (D₂.toNet o₂)
\end{lstlisting}

\subsection{\texttt{Expand}}
The expansion, verbatim.
\begin{lstlisting}[style=leancore]
def splice (refs old : List Port) (maps : List (List (ℕ × ℕ))) : List Port :=
  let here := refs.filter (· ∈ old)
  if here = [] then refs
  else
    let at_ := refs.findIdx (· ∈ old)
    let block := maps.flatMap fun mp => here.map fun r => ⟨(mp.lookup r.node).getD r.node, r.out, r.pos⟩
    (refs.take at_).filter (· ∉ old) ++ block ++ (refs.drop at_).filter (· ∉ old)
\end{lstlisting}
\begin{lstlisting}[style=leancore]
def gather (D : Diagram) (b : DBox) : Diagram :=
  go D b.parent D.boxes.length
where
  go (D : Diagram) (anc : Option ℕ) : ℕ → Diagram
  | 0 => D
  | fuel + 1 =>
    match anc.bind fun a => D.boxes.find? (·.id = a) with
    | none => D
    | some outer =>
      if outer.nodes = [] then go D outer.parent fuel
      else
        let ident := outer.nodes.map fun x => (x, x)
        let cr := D.scopeBnd outer.nodes
        go { D with ins := splice D.ins cr [ident], outs := splice D.outs cr [ident] }
          outer.parent fuel
\end{lstlisting}
\begin{lstlisting}[style=leancore]
def shiftPos (g : Port → Bool) (k : ℕ) (p : Port) : ℕ :=
  let c := ((List.range p.pos).filter fun i => g ⟨p.node, p.out, i⟩).length
  p.pos - c + k * c
\end{lstlisting}
\begin{lstlisting}[style=leancore]
def growSide (g : Port → Bool) (k x : ℕ) (out : Bool) (ds : List ℕ) : List ℕ :=
  ds.zipIdx.flatMap fun (d, i) => if g ⟨x, out, i⟩ then List.replicate k d else [d]
\end{lstlisting}
\begin{lstlisting}[style=leancore]
def growImg (g : Port → Bool) (k : ℕ) (p : Port) : List Port :=
  if g p then (List.range k).map fun j => ⟨p.node, p.out, shiftPos g k p + j⟩
  else [⟨p.node, p.out, shiftPos g k p⟩]
\end{lstlisting}
\begin{lstlisting}[style=leancore]
def cpId (D : Diagram) (b : DBox) (j x : ℕ) : ℕ :=
  D.fresh + j * b.nodes.length + (b.nodes.mergeSort (· ≤ ·)).idxOf x
\end{lstlisting}
\begin{lstlisting}[style=leancore]
def nimg (D : Diagram) (b : DBox) (k : ℕ) (p : Port) : List Port :=
  if inS b p then (List.range k).map fun j => ⟨D.cpId b j p.node, p.out, p.pos⟩
  else growImg (D.crossP b) k p
\end{lstlisting}
\begin{lstlisting}[style=leancore]
def spliceN (D : Diagram) (b : DBox) (k : ℕ) (refs : List Port) : List Port :=
  let here := refs.filter (inS b)
  if here = [] then refs.flatMap (D.nimg b k)
  else
    let at_ := refs.findIdx (inS b)
    (refs.take at_).flatMap (D.nimg b k) ++
      (List.range k).flatMap (fun j => here.map fun p => (⟨D.cpId b j p.node, p.out, p.pos⟩ : Port)) ++
      ((refs.drop at_).filter (!inS b ·)).flatMap (D.nimg b k)
\end{lstlisting}
\begin{lstlisting}[style=leancore]
def copyNodes (D : Diagram) (b : DBox) (k : ℕ) : Diagram :=
  let copies := (List.range k).flatMap fun j => (b.nodes.mergeSort (· ≤ ·)).map (D.cpId b j)
  { nodes := D.nodes.flatMap fun n =>
      if n.id ∈ b.nodes then (List.range k).map fun j => { n with id := D.cpId b j n.id }
      else [{ n with outs := growSide (D.crossP b) k n.id true n.outs,
                     ins := growSide (D.crossP b) k n.id false n.ins }],
    wires := D.wires.flatMap fun w => (D.nimg b k w.1).zip (D.nimg b k w.2),
    outs := D.spliceN b k D.outs, ins := D.spliceN b k D.ins,
    boxes := ((D.boxes.filter (·.id ≠ b.id)).map fun x =>
      { x with nodes := x.nodes.filter (· ∉ b.nodes) ++
                 (if x.nodes.any (· ∈ b.nodes) then copies else []),
               ports := x.ports.flatMap (D.nimg b k) }).filter fun x => x.nodes ≠ [] ∨ x.ports ≠ [] }
\end{lstlisting}
\begin{lstlisting}[style=leancore]
def growPorts (D : Diagram) (b : DBox) (k : ℕ) : Diagram :=
  let g := fun (p : Port) => decide (p ∈ b.ports)
  let img := growImg g k
  { nodes := D.nodes.map fun n =>
      { n with outs := growSide g k n.id true n.outs, ins := growSide g k n.id false n.ins },
    wires := D.wires.flatMap fun w => (img w.1).zip (img w.2),
    outs := D.outs.flatMap img, ins := D.ins.flatMap img,
    boxes := ((D.boxes.filter (·.id ≠ b.id)).map fun x => { x with ports := x.ports.flatMap img }).filter
      fun x => x.nodes ≠ [] ∨ x.ports ≠ [] }
\end{lstlisting}
\begin{lstlisting}[style=leancore]
def expandFuel (ks : ℕ → ℕ) : ℕ → Diagram → Diagram
  | 0, D => D
  | fuel + 1, D =>
    let leaves := D.boxes.filter fun b => !D.boxes.any (·.parent = some b.id)
    match leaves.foldr (fun b acc => match acc with
        | none => some b
        | some c => if b.id < c.id then some b else some c) none with
    | none => D
    | some b =>
      let k := ks b.mult
      expandFuel ks fuel (if b.ports = [] then D.expandNode b k else D.expandPort b k)
\end{lstlisting}

\subsection{\texttt{NetDigits}}
\begin{lstlisting}[style=leancore]
theorem eval_toNetB_evalD (hW : WF D) (E : Env) (off : ℕ) (hoff : D.bnd.length ≤ off)
    (hN : ∀ n ∈ D.nodes, D.NodeS E n)
    (hlab : ∀ j (hj : j < D.wires.length), labDim E (off + j) =
      E.dimv (D.portDim D.wires[j].1) ^ size E.mult (D.legChain D.wires[j].1))
    (β : ℕ → ℕ) (hβ : ∀ i (hi : i < D.bnd.length),
      β i < E.dimv (D.portDim D.bnd[i]) ^ size E.mult (D.legChain D.bnd[i])) :
    ArchytasZX.eval E (D.toNetB off) β = D.evalD E E.mult (D.slotDig E β)
\end{lstlisting}

\subsection{\texttt{OrderK}}
\begin{lstlisting}[style=leancore]
def slotSig (i : ℕ) : List ℕ × List ℕ :=
  let q := D.bnd.getD i Port.dflt
  let bs := D.nodeBoxes q.node
  let mps := bs.map (D.mpB · (decide (i < D.outs.length)))
  let r := mps.findIdx (· = i)
  (mps.take r, List.range' r (bs.length - r) ++ (List.range' bs.length (D.portBoxes q).length).reverse)
\end{lstlisting}
\begin{lstlisting}[style=leancore]
def keyOf (sg : List ℕ × List ℕ) (i : ℕ) (t : List ℕ) : List ℕ :=
  (sg.1.zip t).flatMap (fun p => [p.1, p.2]) ++ [i] ++ sg.2.map (t.getD · 0)
\end{lstlisting}
\begin{lstlisting}[style=leancore]
def expandK (ks : ℕ → ℕ) : Diagram :=
  let E := D.expand ks
  let L := (D.SL ks).map fun x => E.bnd.getD (D.bidx ks x).1 Port.dflt
  { E with outs := L.take E.outs.length, ins := L.drop E.outs.length }
\end{lstlisting}
\begin{lstlisting}[style=leancore]
theorem eval_expandK_evalD (hW : WF D) (E : Env) (σ'' : ℕ → ℕ)
    (hσ : (D.expandK E.mult).InRange E σ'') :
    (D.expandK E.mult).eval E σ'' = D.evalD E E.mult (D.readK E.mult σ'')
\end{lstlisting}

\subsection{\texttt{Assume}}
\begin{lstlisting}[style=leancore]
theorem Equal.of_mapDim (E : Env) (ps : List (ℕ × ℕ)) (hps : ∀ p ∈ ps, E.dimv p.2 = E.dimv p.1)
    {s s' : ℂ} {N M : Net}
    (h : Equal { E with labs
\end{lstlisting}

